% Les systèmes politiques — Allemand Tle A — Cours signé OnBuch+
% Programme officiel MINESEC. Compiler avec : tectonic AL24-les-systemes-politiques.tex
\newcommand{\DOCMATIERE}{Allemand}
\newcommand{\DOCNIVEAU}{Tle A}
\newcommand{\DOCMODULE}{Chapitre 5 : Citoyenneté}
\newcommand{\DOCLECON}{AL24}
\newcommand{\DOCTITRE}{Les systèmes politiques}
% OnBuch+ — préambule « cours signés » ENRICHI (compilé avec tectonic / XeLaTeX)
% Système visuel « Pop » d'OnBuch+ : fond crème (jamais blanc), bordures encre
% épaisses, ombres « dures » décalées, titres Archivo Black en capitales.
% Version MATHÉMATIQUES : graphes (pgfplots/tikz), boîtes pédagogiques
% (définition, propriété, méthode, attention, le savais-tu, exercice résolu,
% exercice, TP, bilan), page de garde et signature OnBuch+ ancrée en bas.
%
% Macros attendues dans main.tex AVANT \input{preamble} :
%   \DOCMATIERE  \DOCNIVEAU  \DOCMODULE  \DOCLECON (numéro)  \DOCTITRE
\documentclass[11pt]{article}

\usepackage{fontspec}
\usepackage[ngerman,french]{babel} % français = langue principale ; l'allemand via \all{..} / textebox
\usepackage[a4paper,margin=2cm,top=3.4cm,bottom=2.8cm,headheight=1.4cm,footskip=1.3cm]{geometry}
\usepackage{amsmath,amssymb,amsfonts}
\usepackage{mathtools} % \xmapsto, \coloneqq... souvent utilisés par les modèles
\usepackage{cancel} % simplifications barrées (bilans de forces)
% \abs{z} (module, valeur absolue) et \norm{u} (norme), utilisés par les modèles
\providecommand{\abs}{}\let\abs\relax\DeclarePairedDelimiter{\abs}{\lvert}{\rvert}
\providecommand{\norm}{}\let\norm\relax\DeclarePairedDelimiter{\norm}{\lVert}{\rVert}
\usepackage[table]{xcolor} % [table] : fournit aussi \rowcolors (alternance de lignes)
\usepackage{pagecolor}
\usepackage{tikz}
\usetikzlibrary{arrows.meta,positioning,calc,shapes.geometric,shapes.misc,shapes.arrows,decorations.pathmorphing,decorations.pathreplacing,decorations.markings,patterns,shadows,mindmap,trees,fit,backgrounds,matrix,intersections,angles,quotes}
\usepackage{pgfplots}
\usepackage[version=4]{mhchem} % équations nucléaires \ce{^{235}_{92}U}
\usepackage[european,siunitx]{circuitikz} % schémas électriques
\pgfplotsset{compat=1.18}
\usepgfplotslibrary{fillbetween,groupplots} % aires entre courbes (calcul intégral)
\usepackage{enumitem}
\usepackage{fancyhdr}
% [most] charge toutes les bibliothèques tcolorbox (listings, minted, raster,
% documentation...) et gonfle inutilement la mémoire TeX sur un document long
% (~300 boîtes) : on ne charge que ce qui est réellement utilisé.
\usepackage[skins,breakable]{tcolorbox}
\usepackage{graphicx}
\usepackage{colortbl}
\usepackage{booktabs}
\usepackage{tabularx}
\usepackage{diagbox} % case d'en-tête diagonale des tableaux à double entrée
% Colonnes extensibles centrée / gauche / droite (C, L, R), souvent utilisées par les modèles
\newcolumntype{C}{>{\centering\arraybackslash}X}
\newcolumntype{L}{>{\raggedright\arraybackslash}X}
\newcolumntype{R}{>{\raggedleft\arraybackslash}X}
\usepackage{array}
\usepackage{multicol}
\usepackage{multirow}
\usepackage{siunitx}
% parse-numbers=false : accepte aussi \SI{2,5 \times 10^{-3}}{...}
\sisetup{locale=FR,output-decimal-marker={,},group-minimum-digits=4,parse-numbers=false}
\DeclareSIUnit\ohm{\ensuremath{\symup{\Omega}}} % U+2126 absent de la police
\DeclareSIUnit\division{div} % réglages d'oscilloscope (ms/div, V/div)
\DeclareSIUnit\tr{tr} % tours (vitesse de rotation, ex. \SI{1500}{\tr\per\minute})
\DeclareSIUnit\molar{M} % concentration molaire (ex. \SI{3}{\molar}, \SI{1}{\micro\molar})
\DeclareSIUnit\an{an} % année (le modèle confond parfois avec \year, un compteur TeX) — ex. \SI{5}{\kilo\meter\per\an}
\DeclareSIUnit\FCFA{FCFA} % franc CFA (ex. \SI{500}{\FCFA\per\kilo\gram})
\DeclareSIUnit\degreeBrix{\textdegree Brix} % teneur en sucre d'un sirop/jus (réfractométrie)
\DeclareSIUnit\month{mois} % le modèle utilise parfois \month, un compteur TeX, comme unité de durée
\usepackage{qrcode}
% \degree : commande fréquemment utilisée par les modèles pour un angle ou une
% température en dehors de \SI{}{} (non fournie par les paquets chargés).
\providecommand{\degree}{^{\circ}}
% \qed (fin de démonstration), utilisé par les modèles sans amsthm
\providecommand{\qed}{\hfill$\square$}
% \barre : double barre des valeurs interdites dans un tableau de variations (tkz-tab)
\providecommand{\barre}{\ensuremath{\|}}
% environnement proof (amsthm non chargé) utilisé par les modèles
\ifdefined\proof\else\newenvironment{proof}[1][Démonstration]{\par\noindent\textit{#1.}\ }{\qed\par}\fi
\usepackage{lastpage}
\usepackage{hyperref}
\hypersetup{hidelinks}

% ============================================================
% Palette « Pop » OnBuch+ — mêmes hex que la classe P (plus_ui.dart)
% ============================================================
\definecolor{poporange}{HTML}{F59321}
\definecolor{popdark}{HTML}{E07A0C}
\definecolor{popcream}{HTML}{FFF6E8}
\definecolor{popink}{HTML}{1C1714}
\definecolor{poppurple}{HTML}{7A5AE0}
\definecolor{popgreen}{HTML}{2E9E6A}
\definecolor{poppink}{HTML}{E0457B}
\definecolor{popblue}{HTML}{2F7DE1}
\definecolor{popgold}{HTML}{C98A12}
\definecolor{popmuted}{HTML}{8A7F76}
\definecolor{popline}{HTML}{EADFD2}
\definecolor{popgray}{HTML}{8A7F76}
\colorlet{popgrey}{popgray}
% Alias tolérants (le modèle nomme parfois les couleurs différemment)
\colorlet{popwhite}{white}
\colorlet{popblack}{popink}
\colorlet{popred}{poppink}
\colorlet{popyellow}{popgold}
% Teintes claires (fonds de tableaux/encadrés) — le modèle nomme parfois
% les couleurs avec un suffixe L (ex. popblueL) sans les avoir définies.
\colorlet{poporangeL}{poporange!20}
\colorlet{popdarkL}{popdark!20}
\colorlet{poppurpleL}{poppurple!20}
\colorlet{popgreenL}{popgreen!20}
\colorlet{poppinkL}{poppink!20}
\colorlet{popblueL}{popblue!20}
\colorlet{popgoldL}{popgold!20}
\colorlet{popredL}{popred!20}
\colorlet{popgrayL}{popgray!20}
% Teintes claires pour fonds de tableaux / zones
\colorlet{popblueL}{popblue!10!white}
\colorlet{poporangeL}{poporange!14!white}
\colorlet{popgreenL}{popgreen!12!white}
\colorlet{poppurpleL}{poppurple!12!white}
\colorlet{poppinkL}{poppink!12!white}
\colorlet{popgoldL}{popgold!14!white}

\pagecolor{popcream}

% ============================================================
% Typographie
% ============================================================
\defaultfontfeatures{Ligatures=TeX}
\setmainfont{PlusJakartaSans-Medium.ttf}[Path=../../../fonts/,BoldFont=PlusJakartaSans-Bold.ttf]
% Symboles, API et emojis absents de la police principale : repli sur DejaVu Sans (caractères actifs)
\newfontface\symfont{DejaVuSans.ttf}[Path=../../../fonts/]
\makeatletter
\newcommand\symactive[1]{\catcode#1=13 \begingroup\lccode`\~=#1 \lowercase{\endgroup\def~}{{\symfont\char#1}}}
\newcommand\symblank[1]{\catcode#1=13 \begingroup\lccode`\~=#1 \lowercase{\endgroup\def~}{}}
\newcommand\symhyph[1]{\catcode#1=13 \begingroup\lccode`\~=#1 \lowercase{\endgroup\def~}{\mbox{-}}}
\newcount\symc
\newcommand\symrange[2]{\symc=#1 \@whilenum\symc<#2 \do{\expandafter\symactive\expandafter{\the\symc}\advance\symc by 1 }}
\newcommand\symblankrange[2]{\symc=#1 \@whilenum\symc<#2 \do{\expandafter\symblank\expandafter{\the\symc}\advance\symc by 1 }}
\makeatother
\symrange{"0250}{"02FF}\symactive{"2713}\symactive{"2714}\symactive{"2717}\symactive{"2718}\symactive{"2610}\symactive{"2611}\symactive{"2612}
\symhyph{"2011}\symhyph{"2E1A}\symblankrange{"1F300}{"1FAFF}\symblank{"FE0F}
% Maths en sans-serif (Fira Math) avec chiffres et lettres droites de Plus
% Jakarta, pour que les formules se fondent dans le texte.
\usepackage[math-style=ISO,bold-style=ISO]{unicode-math}
\setmathfont{FiraMath-Regular.otf}[Path=../../../fonts/]
\setmathfont{PlusJakartaSans-Medium.ttf}[Path=../../../fonts/,range={up/num,up/{Latin,latin},\mathup}]
\usepackage{icomma} % virgule décimale sans espace en mode math
% Caractères mathématiques Unicode absents des polices de texte (Plus Jakarta,
% Archivo Black) : rendus par leur équivalent mathématique, en texte comme en
% formule — sinon ils s'affichent en carré vide (ex. « Arithmétique dans ℤ »).
\usepackage{newunicodechar}
\newunicodechar{ℝ}{\ensuremath{\mathbb{R}}}
\newunicodechar{ℤ}{\ensuremath{\mathbb{Z}}}
\newunicodechar{ℂ}{\ensuremath{\mathbb{C}}}
\newunicodechar{ℕ}{\ensuremath{\mathbb{N}}}
\newunicodechar{ℚ}{\ensuremath{\mathbb{Q}}}
\newunicodechar{𝔻}{\ensuremath{\mathbb{D}}}
\newunicodechar{α}{\ensuremath{\alpha}}
\newunicodechar{β}{\ensuremath{\beta}}
\newunicodechar{γ}{\ensuremath{\gamma}}
\newunicodechar{δ}{\ensuremath{\delta}}
\newunicodechar{ε}{\ensuremath{\varepsilon}}
\newunicodechar{θ}{\ensuremath{\theta}}
\newunicodechar{λ}{\ensuremath{\lambda}}
\newunicodechar{μ}{\ensuremath{\mu}}
\newunicodechar{σ}{\ensuremath{\sigma}}
\newunicodechar{τ}{\ensuremath{\tau}}
\newunicodechar{φ}{\ensuremath{\varphi}}
\newunicodechar{ψ}{\ensuremath{\psi}}
\newunicodechar{ω}{\ensuremath{\omega}}
\newunicodechar{Σ}{\ensuremath{\Sigma}}
% majuscules produites par \MakeUppercase dans les titres (α-aminés -> Α)
\newunicodechar{Α}{\ensuremath{\alpha}}
\newunicodechar{Β}{\ensuremath{\beta}}
\newunicodechar{Γ}{\ensuremath{\Gamma}}
\newunicodechar{Δ}{\ensuremath{\Delta}}
\newunicodechar{Ε}{\ensuremath{\varepsilon}}
\newunicodechar{Θ}{\ensuremath{\Theta}}
\newunicodechar{Λ}{\ensuremath{\Lambda}}
\newunicodechar{Μ}{\ensuremath{\mu}}
\newunicodechar{Τ}{\ensuremath{\tau}}
\newunicodechar{Φ}{\ensuremath{\Phi}}
\newunicodechar{Ψ}{\ensuremath{\Psi}}
\newunicodechar{Ω}{\ensuremath{\Omega}}
\newunicodechar{↦}{\ensuremath{\mapsto}}
\newunicodechar{∈}{\ensuremath{\in}}
\newunicodechar{∉}{\ensuremath{\notin}}
\newunicodechar{∧}{\ensuremath{\wedge}}
\newunicodechar{⇔}{\ensuremath{\Leftrightarrow}}
\newunicodechar{⟺}{\ensuremath{\Longleftrightarrow}}
\newunicodechar{⇒}{\ensuremath{\Rightarrow}}
\newunicodechar{⟹}{\ensuremath{\Longrightarrow}}
\newunicodechar{∘}{\ensuremath{\circ}}
\newunicodechar{∪}{\ensuremath{\cup}}
\newunicodechar{∩}{\ensuremath{\cap}}
\newunicodechar{≡}{\ensuremath{\equiv}}
\newunicodechar{⊂}{\ensuremath{\subset}}
\newunicodechar{⊥}{\ensuremath{\perp}}
\newunicodechar{∃}{\ensuremath{\exists}}
\newunicodechar{∀}{\ensuremath{\forall}}
\newunicodechar{∛}{\ensuremath{\sqrt[3]{\,}}}
\newunicodechar{∜}{\ensuremath{\sqrt[4]{\,}}}
\newunicodechar{⃗}{\ensuremath{{}^{\to}}}
\newunicodechar{ⁿ}{\ifmmode^{n}\else\textsuperscript{\ensuremath{n}}\fi}
\newunicodechar{ˣ}{\ifmmode^{x}\else\textsuperscript{\ensuremath{x}}\fi}
\newunicodechar{ᵇ}{\ifmmode^{b}\else\textsuperscript{\ensuremath{b}}\fi}
\newunicodechar{ʳ}{\ifmmode^{r}\else\textsuperscript{\ensuremath{r}}\fi}
\newunicodechar{ᵘ}{\ifmmode^{u}\else\textsuperscript{\ensuremath{u}}\fi}
\newunicodechar{ʸ}{\ifmmode^{y}\else\textsuperscript{\ensuremath{y}}\fi}
\newunicodechar{ᵃ}{\ifmmode^{a}\else\textsuperscript{\ensuremath{a}}\fi}
\newunicodechar{ᵉ}{\ifmmode^{e}\else\textsuperscript{\ensuremath{e}}\fi}
\newunicodechar{ᵏ}{\ifmmode^{k}\else\textsuperscript{\ensuremath{k}}\fi}
\newunicodechar{ᵖ}{\ifmmode^{p}\else\textsuperscript{\ensuremath{p}}\fi}
\newunicodechar{ᴺ}{\ifmmode^{N}\else\textsuperscript{\ensuremath{N}}\fi}
\newunicodechar{ᶜ}{\ifmmode^{c}\else\textsuperscript{\ensuremath{c}}\fi}
\newunicodechar{ʰ}{\ifmmode^{h}\else\textsuperscript{\ensuremath{h}}\fi}
\newunicodechar{ᵠ}{\ifmmode^{q}\else\textsuperscript{\ensuremath{q}}\fi}
\newunicodechar{ₙ}{\ifmmode_{n}\else\textsubscript{\ensuremath{n}}\fi}
\newunicodechar{ₐ}{\ifmmode_{a}\else\textsubscript{\ensuremath{a}}\fi}
\newunicodechar{ᵢ}{\ifmmode_{i}\else\textsubscript{\ensuremath{i}}\fi}
\newunicodechar{ᵣ}{\ifmmode_{r}\else\textsubscript{\ensuremath{r}}\fi}
\newunicodechar{ₖ}{\ifmmode_{k}\else\textsubscript{\ensuremath{k}}\fi}
\newunicodechar{ᵦ}{\ifmmode_{\beta}\else\textsubscript{\ensuremath{\beta}}\fi}
\newunicodechar{ₓ}{\ifmmode_{x}\else\textsubscript{\ensuremath{x}}\fi}
\newfontfamily\popdisplay{ArchivoBlack-Regular.ttf}[Path=../../../fonts/]
\newfontfamily\poplogo{SpaceGrotesk-Bold.ttf}[Path=../../../fonts/]
\newfontfamily\popsemi{PlusJakartaSans-SemiBold.ttf}[Path=../../../fonts/]
\newfontfamily\popxbold{PlusJakartaSans-ExtraBold.ttf}[Path=../../../fonts/]
\newfontfamily\popxboldls{PlusJakartaSans-ExtraBold.ttf}[Path=../../../fonts/,LetterSpace=3.0]

\setlist[enumerate]{leftmargin=1.7em,itemsep=5pt,topsep=4pt}
\setlist[itemize]{leftmargin=1.7em,itemsep=4pt,topsep=4pt,label={\color{poporange}\textbullet}}
\setlength{\parindent}{0pt}
\setlength{\parskip}{5pt}
\renewcommand{\arraystretch}{1.25}

% pgfplots : style Pop par défaut
\pgfplotsset{
  every axis/.append style={axis line style={popink,line width=1pt},tick style={popink},
    grid style={popline,line width=0.5pt},label style={font=\small},tick label style={font=\footnotesize}},
  popcurve/.style={poporange,line width=1.8pt,no markers,smooth},
}
% Même style disponible dans un \draw TikZ simple (hors pgfplots)
\tikzset{popcurve/.style={poporange,line width=1.8pt}}
\tikzset{popgrid/.style={popline,very thin}}
\colorlet{popcurve}{poporange} % le nom du style est parfois utilisé comme couleur

% ============================================================
% Pill / squiggle
% ============================================================
\newcommand{\poppill}[3]{%
  \tikz[baseline=-0.55ex]\node[draw=popink,line width=1.2pt,rounded corners=2.6pt,%
    fill=#2,inner xsep=6.5pt,inner ysep=3pt]{%
    \popxboldls\fontsize{7}{7}\selectfont\color{#3}\MakeUppercase{#1}};%
}
\newcommand{\popsquiggle}[1]{%
  \tikz[baseline=-0.4ex]{
    \draw[#1,line width=2.6pt,line cap=round]
      plot [smooth] coordinates {(0,0.09) (0.34,0.01) (0.68,0.21) (1.02,0.01) (1.36,0.10)};
  }%
}
% Mot-clé surligné (marqueur orange sous le texte)
\newcommand{\cle}[1]{{\popxbold\color{popdark}#1}}

% ============================================================
% En-tête / pied
% ============================================================
\pagestyle{fancy}
\fancyhf{}
\renewcommand{\headrulewidth}{0pt}
\renewcommand{\footrulewidth}{0pt}
\lhead{\raisebox{-0.15ex}{\poplogo\fontsize{13}{13}\selectfont\color{popink}OnBuch\color{poporange}+}}
\rhead{\poppill{\DOCMATIERE\ \textbullet\ \DOCNIVEAU\ \textbullet\ Leçon \DOCLECON}{white}{popink}}
\lfoot{\popxbold\fontsize{7.6}{7.6}\selectfont\color{popmuted}\MakeUppercase{Cours signé OnBuch+}\ \textbullet\ {\popsemi\DOCTITRE}}
\rfoot{\tikz[baseline=-0.6ex]\node[rounded corners=8.5pt,fill=popink,minimum height=17pt,minimum width=17pt,inner xsep=5pt,inner sep=0]{\popxbold\fontsize{8}{8}\selectfont\color{popcream}\thepage\,/\,\pageref*{LastPage}};}

% ============================================================
% Titres
% ============================================================
\newcommand{\chaptitre}[2]{%
  \begin{tcolorbox}[enhanced,colback=poporange,colframe=popink,boxrule=2.6pt,arc=11pt,
      drop shadow=popink,left=15pt,right=15pt,top=12pt,bottom=11pt,before skip=3pt,after skip=18pt]
    \poppill{#2}{popink}{popcream}\par\vspace{9pt}
    {\raggedright\hyphenpenalty=10000\exhyphenpenalty=10000\popdisplay\fontsize{22}{24}\selectfont\color{popcream}\MakeUppercase{#1}\par}
  \end{tcolorbox}
}
% Section numérotée : pastille orange avec le numéro + titre Archivo
\newcounter{popsec}
\newcommand{\coursec}[1]{%
  \stepcounter{popsec}\setcounter{popsub}{0}%
  \par\vspace{16pt}\noindent
  \begin{minipage}{\linewidth}
  \tikz[baseline=-0.7ex]\node[draw=popink,line width=1.6pt,fill=poporange,rounded corners=4pt,minimum size=22pt,inner sep=0]{\popdisplay\fontsize{12}{12}\selectfont\color{popcream}\thepopsec};%
  \hspace{8pt}{\popdisplay\fontsize{15}{17}\selectfont\color{popink}\MakeUppercase{#1}}\par
  \vspace{4pt}\noindent{\color{poporange}\rule{36pt}{3.6pt}}
  \end{minipage}\par\nopagebreak\vspace{8pt}
}
\newcounter{popsub}
\newcommand{\courssub}[1]{%
  \stepcounter{popsub}%
  \par\vspace{9pt}\noindent{\popxbold\fontsize{11.5}{13}\selectfont\color{popink}{\color{poporange}\thepopsec.\thepopsub}\hspace{6pt}#1}\par\nopagebreak\vspace{4pt}
}

% ============================================================
% Boîtes pédagogiques — même recette : fond blanc, bordure épaisse colorée,
% ombre dure encre, badge plein. Titre optionnel [#1].
% ============================================================
\tcbset{popbase/.style={breakable,enhanced,colback=white,boxrule=2.3pt,arc=9pt,
  shadow={3pt}{-3pt}{0pt}{popink},left=14pt,right=14pt,top=11pt,bottom=11pt,before skip=11pt,after skip=15pt,
  fontupper=\popsemi\fontsize{10.6}{14.5}\selectfont\color{popink},
  fontlower=\popsemi\fontsize{10.6}{14.5}\selectfont\color{popink}}}
% Coche dessinée (le glyphe ✓ n'existe pas dans les polices du document)
\newcommand{\popcheck}[1]{\tikz[baseline=-0.5ex]\draw[#1,line width=1.6pt,line cap=round,line join=round](-0.13,0.0)--(-0.04,-0.09)--(0.13,0.1);}
\providecommand{\coche}{\popcheck{popgreen}}
\providecommand{\resultat}[1]{\par\smallskip\noindent\fcolorbox{popgreen}{popgreenL}{\parbox{\dimexpr\linewidth-2\fboxsep-2\fboxrule\relax}{\textbf{Résultat.} #1}}\par\smallskip}
\newcommand{\popboxhead}[3]{\poppill{#1}{#2}{white}\ifx\relax#3\relax\else\hspace{6pt}{\popxbold\fontsize{10.6}{12}\selectfont\color{popink}#3}\fi\par\vspace{7pt}}

\newtcolorbox{aretenir}[1][]{popbase,colframe=popgreen,before upper={\popboxhead{À retenir}{popgreen}{#1}}}
% Texte en allemand (document de lecture, dialogue, extrait) : cadre
% d'étude. Un titre optionnel (source, auteur) s'affiche dans l'en-tête.
\newtcolorbox{textebox}[1][]{popbase,colframe=popblue,before upper={\popboxhead{Text}{popblue}{#1}\begin{otherlanguage}{ngerman}},after upper={\end{otherlanguage}}}
% Marques ✓ / ✗ (forme correcte / fautive) souvent utilisées par les modèles
\providecommand{\cross}{\textcolor{poppink}{\ensuremath{\times}}}
\providecommand{\xmark}{\cross}\providecommand{\crossmark}{\cross}\providecommand{\cmark}{\ensuremath{\checkmark}}
% Ligne de réponse / justification à compléter (inventée par les modèles)
\providecommand{\justification}[1]{\par\noindent\textit{Justification~:} \dotfill\par}
% \textipa (package tipa non chargé) : la police ne couvre pas l'API -> simple passage
\providecommand{\textipa}[1]{#1}
% Soulignement (ulem non chargé) : \uline, \ul
\providecommand{\uline}[1]{\underline{#1}}\providecommand{\ul}[1]{\underline{#1}}
% \glossaire{\item ...} inventée par les modèles : liste de glossaire
\providecommand{\glossaire}[1]{\begin{itemize}#1\end{itemize}}
% \alert (beamer) inventée par les modèles : texte mis en évidence
\providecommand{\alert}[1]{\textbf{\textcolor{poppink}{#1}}}
% variantes ASCII d'ordinaux écrites en macro
\providecommand{\iere}{\textsuperscript{re}}\providecommand{\ieme}{\textsuperscript{e}}\providecommand{\ere}{\textsuperscript{re}}\providecommand{\eme}{\textsuperscript{e}}\providecommand{\er}{\textsuperscript{er}}\providecommand{\ie}{\textsuperscript{e}}
% Ordinaux français écrits en macro par les modèles : 1\ière, 3\ième
\newcommand{\ière}{\textsuperscript{re}}\newcommand{\ième}{\textsuperscript{e}}
% \exemple (inventée par les modèles) : étiquette « Exemple : »
\providecommand{\exemple}{\textbf{Exemple~:}\ }
% \square utilisable aussi hors mode math (cases à cocher)
\AtBeginDocument{\let\squareMath\square\renewcommand{\square}{\ensuremath{\squareMath}}}
% Mot ou phrase en allemand à l'intérieur d'un paragraphe français
\newcommand{\all}[1]{\ifmmode\text{\textit{#1}}\else\begingroup\selectlanguage{ngerman}\itshape #1\endgroup\fi}
\let\esp\all % alias toléré
\newtcolorbox{exemplebox}[1][]{popbase,colframe=poporange,before upper={\popboxhead{Exemple}{poporange}{#1}}}
\newtcolorbox{definition}[1][]{popbase,colframe=popblue,before upper={\popboxhead{Définition}{popblue}{#1}}}
\newtcolorbox{propriete}[1][]{popbase,colframe=poppurple,before upper={\popboxhead{Propriété}{poppurple}{#1}}}
\newtcolorbox{methode}[1][]{popbase,colframe=popgold,before upper={\popboxhead{Méthode}{popgold}{#1}}}
\newtcolorbox{attention}[1][]{popbase,colframe=poppink,before upper={\popboxhead{Attention}{poppink}{#1}}}
\newtcolorbox{experience}[1][]{popbase,colframe=popblue,colback=popblueL,before upper={\popboxhead{Expérience}{popblue}{#1}}}
% « Le savais-tu ? » : carte violette pleine (contexte, culture, Cameroun)
\newtcolorbox{savaistu}[1][]{popbase,colback=poppurple,colframe=popink,
  fontupper=\popsemi\fontsize{10.4}{14.2}\selectfont\color{white},
  before upper={\poppill{Le savais-tu\,?}{white}{poppurple}\ifx\relax#1\relax\else\hspace{6pt}{\popxbold\color{white}#1}\fi\par\vspace{7pt}}}
% Exercice résolu : énoncé en haut, \tcblower puis la solution sur fond crème
\newcounter{popexo}\newcounter{popexr}
\newtcolorbox{exoresolu}[1][]{popbase,colframe=poporange,colbacklower=popcream,
  segmentation style={popink,line width=1.2pt,dashed},
  before upper={\stepcounter{popexr}\popboxhead{Exercice résolu \thepopexr}{poporange}{#1}},
  before lower={\poppill{Solution}{popink}{popcream}\par\vspace{6pt}}}
% Exercice (non résolu en place) — la correction est dans la partie Corrigés
\newtcolorbox{exercice}[1][]{popbase,colframe=popink,
  before upper={\stepcounter{popexo}\popboxhead{Exercice \thepopexo}{popink}{#1}}}
% Corrigé
\newtcolorbox{corrige}[1][]{popbase,colframe=popgreen,colback=popgreenL,
  before upper={\popboxhead{Corrigé}{popgreen}{#1}}}
% Objectifs / prérequis / bilan
\newtcolorbox{objectifs}{popbase,colframe=popink,colback=poporangeL,
  before upper={\popboxhead{Objectifs de la leçon}{popink}{}}}
\newtcolorbox{prerequis}{popbase,colframe=popink,colback=popblueL,
  before upper={\popboxhead{Prérequis}{popblue}{}}}
\newtcolorbox{bilan}{popbase,colframe=popink,colback=white,boxrule=2.8pt,
  before upper={\popboxhead{Fiche bilan}{poporange}{L'essentiel à connaître pour l'examen}}}
% Figure encadrée avec légende
\newtcolorbox{popfigure}[1][]{enhanced,breakable=false,colback=white,colframe=popline,boxrule=1.4pt,arc=8pt,
  left=10pt,right=10pt,top=10pt,bottom=8pt,before skip=10pt,after skip=12pt,halign=center,
  lower separated=false,halign lower=center,
  fontlower=\popsemi\fontsize{9}{11.5}\selectfont\color{popmuted},
  #1}
\newcommand{\legende}[1]{\par\vspace{6pt}{\popsemi\fontsize{9}{11.5}\selectfont\color{popmuted}#1}}

% ============================================================
% Signature OnBuch+
% ============================================================
\newcommand{\poplogoinline}[1]{{\poplogo\fontsize{#1}{#1}\selectfont\color{popink}OnBuch\color{poporange}+}}
\newcommand{\signatureonbuch}{%
  \begin{tcolorbox}[enhanced,colback=popink,colframe=popink,boxrule=2.2pt,arc=10pt,
      drop shadow=poporange,left=14pt,right=14pt,top=10pt,bottom=10pt,before skip=0pt,after skip=0pt]
    \tikz[baseline=-0.7ex]\node[circle,fill=popgreen,draw=popcream,line width=1.3pt,minimum size=22pt,inner sep=0]{\popcheck{white}};%
    \hspace{9pt}\begin{minipage}[c]{0.62\linewidth}
      {\popxbold\fontsize{12}{13}\selectfont\color{popcream}Signé\ \textbullet\ {\poplogo OnBuch\color{poporange}+}}\par\vspace{2pt}
      {\raggedright\popsemi\fontsize{8}{10.5}\selectfont\color{popline}\MakeUppercase{Équipe pédagogique OnBuch+}\\\MakeUppercase{Conforme au programme officiel MINESEC}\par}
    \end{minipage}\hfill
    {\poplogo\fontsize{22}{22}\selectfont\color{popcream}OnBuch\color{poporange}+}
  \end{tcolorbox}%
}

% ============================================================
% Page de garde
% #1 titre · #2 matière · #3 niveau · #4 module · #5 n° leçon · #6 durée ·
% #7 description / objectifs (texte ou itemize)
% ============================================================
\newcommand{\pagegarde}[7]{%
  \thispagestyle{empty}
  \newgeometry{a4paper,margin=1.7cm,top=1.6cm,bottom=1.4cm}%
  \begin{tikzpicture}[remember picture,overlay]
    \fill[poporange!18] ([xshift=-3.2cm,yshift=-3.6cm]current page.north east) circle (4.6cm);
    \fill[poppurple!14] ([xshift=2.4cm,yshift=7.6cm]current page.south west) circle (3.8cm);
  \end{tikzpicture}%
  {\poplogo\fontsize{30}{30}\selectfont\color{popink}OnBuch\color{poporange}+}\hfill
  \raisebox{0.6ex}{\poppill{Cours signé \textbullet\ Premium}{popink}{popcream}}\par
  \vspace{1pt}\popsquiggle{poppurple}\par
  \vspace{8pt}
  \begin{tcolorbox}[enhanced,colback=poporange,colframe=popink,boxrule=2.8pt,arc=13pt,
      drop shadow=popink,left=18pt,right=18pt,top=12pt,bottom=13pt,after skip=10pt]
    \poppill{#2 \textbullet\ #3}{popink}{popcream}\hspace{5pt}\poppill{#4}{white}{popink}\par\vspace{8pt}
    {\popdisplay\fontsize{14}{14}\selectfont\color{popink}LEÇON #5}\par\vspace{4pt}
    {\raggedright\hyphenpenalty=10000\exhyphenpenalty=10000\popdisplay\fontsize{25}{27}\selectfont\color{popcream}\MakeUppercase{#1}\par}
    \vspace{8pt}
    {\popxbold\fontsize{9.5}{9.5}\selectfont\color{popink}\MakeUppercase{Durée conseillée : #6}}
  \end{tcolorbox}
  \begin{tcolorbox}[enhanced,colback=white,colframe=popink,boxrule=2.2pt,arc=10pt,
      drop shadow=popink,left=16pt,right=16pt,top=10pt,bottom=10pt,after skip=8pt]
    \poppill{Dans cette leçon}{poppurple}{white}\par\vspace{5pt}
    \popsemi\fontsize{9.3}{12.6}\selectfont\color{popink}#7
  \end{tcolorbox}
  \noindent\begin{minipage}[t]{0.487\linewidth}
    \begin{tcolorbox}[enhanced,colback=popgreen,colframe=popink,boxrule=2.2pt,arc=10pt,
        drop shadow=popink,left=11pt,right=11pt,top=10pt,bottom=10pt,height=3.1cm,valign=center]
      \tcbox[colback=white,colframe=popink,boxrule=1.3pt,arc=5pt,boxsep=0pt,nobeforeafter,
        left=3pt,right=3pt,top=3pt,bottom=3pt,valign=center]{\qrcode[height=1.9cm]{https://play.google.com/store/apps/details?id=com.luvvix.onbuch&hl=fr}}%
      \hspace{8pt}\begin{minipage}[c]{0.5\linewidth}
        {\popdisplay\fontsize{11}{12.5}\selectfont\color{white}\MakeUppercase{Télécharge OnBuch}}\par\vspace{4pt}
        {\raggedright\popsemi\fontsize{8.4}{11}\selectfont\color{white}Gratuit \textbullet\ Google Play\\Tuteur IA, annales, concours, résultats\par}
      \end{minipage}
    \end{tcolorbox}
  \end{minipage}\hfill
  \begin{minipage}[t]{0.487\linewidth}
    \begin{tcolorbox}[enhanced,colback=poppurple,colframe=popink,boxrule=2.2pt,arc=10pt,
        drop shadow=popink,left=11pt,right=11pt,top=10pt,bottom=10pt,height=3.1cm,valign=center]
      \tikz[baseline=-0.7ex]\node[circle,fill=white,draw=popink,line width=1.3pt,minimum size=1.6cm,inner sep=0]{\popdisplay\fontsize{24}{24}\selectfont\color{poppurple}+};%
      \hspace{8pt}\begin{minipage}[c]{0.6\linewidth}
        {\popdisplay\fontsize{11}{12.5}\selectfont\color{white}\MakeUppercase{Passe à OnBuch+}}\par\vspace{4pt}
        {\raggedright\popsemi\fontsize{8.4}{11}\selectfont\color{white}Tous les cours signés, Tuteur IA illimité, fiches PDF — onglet OnBuch+ de l'app\par}
      \end{minipage}
    \end{tcolorbox}
  \end{minipage}
  \vspace{8pt}
  \popcontenu
  \vfill
  \signatureonbuch
  \restoregeometry
  \newpage
}

% Rangée « Ce cours contient » (page de garde)
\newcommand{\poptile}[3]{% #1 couleur #2 gros texte #3 légende
  \begin{tcolorbox}[enhanced,colback=#1,colframe=popink,boxrule=2pt,arc=9pt,shadow={2.5pt}{-2.5pt}{0pt}{popink},
      left=6pt,right=6pt,top=9pt,bottom=9pt,halign=center,valign=center,height=1.75cm,width=\linewidth,nobeforeafter]
    {\popdisplay\fontsize{12.5}{13}\selectfont\color{white}\MakeUppercase{#2}}\par\vspace{3pt}
    {\centering\popsemi\fontsize{7.8}{9.5}\selectfont\color{white}#3\par}
  \end{tcolorbox}}
\newcommand{\popcontenu}{%
  \par\noindent{\popxboldls\fontsize{7.5}{7.5}\selectfont\color{popmuted}CE COURS SIGNÉ CONTIENT}\par\vspace{7pt}
  \noindent\begin{minipage}[t]{0.235\linewidth}\poptile{popblue}{Cours}{complet, illustré, pas à pas}\end{minipage}\hfill
  \begin{minipage}[t]{0.235\linewidth}\poptile{poporange}{Méthodes}{et exercices résolus}\end{minipage}\hfill
  \begin{minipage}[t]{0.235\linewidth}\poptile{poppink}{Exercices}{type examen + corrigés}\end{minipage}\hfill
  \begin{minipage}[t]{0.235\linewidth}\poptile{popgreen}{Fiche bilan}{l'essentiel à retenir}\end{minipage}\par}

% Sommaire
\newtcolorbox{sommaire}{enhanced,colback=white,colframe=popink,boxrule=2.2pt,arc=10pt,
  drop shadow=popink,left=18pt,right=18pt,top=13pt,bottom=13pt,before skip=6pt,after skip=20pt,
  before upper={\poppill{Sommaire}{poppurple}{white}\par\vspace{9pt}\popsemi\fontsize{10.6}{14.5}\selectfont\color{popink}}}

% Signature de fin de cours — poussée tout en bas de la dernière page
\newcommand{\findecours}{%
  \par\vfill\nopagebreak
  \begin{center}\popsquiggle{poporange}\end{center}\vspace{4pt}
  \signatureonbuch
}
\AtBeginDocument{\def\llbracket{\mathopen{[\mkern-3.5mu[}}\def\rrbracket{\mathclose{]\mkern-3.5mu]}}} % intervalles d'entiers (glyphe ⟦ absent de la police)
% Glyphes absents de Fira Math : redéfinis à partir de glyphes disponibles
\AtBeginDocument{%
  \def\setminus{\mathbin{\backslash}}\let\smallsetminus\setminus
  \def\vdots{\mathord{\vbox{\baselineskip3.2pt\lineskiplimit0pt\kern4pt\hbox{.}\hbox{.}\hbox{.}}}}%
  \def\ddots{\mathinner{\mkern1mu\raise6.5pt\hbox{.}\mkern2mu\raise3.6pt\hbox{.}\mkern2mu\raise0.7pt\hbox{.}\mkern1mu}}%
  \def\triangle{\mathord{\scalebox{0.9}{$\symup{\Delta}$}}}%
  \def\nabla{\mathord{\raisebox{\depth}{\scalebox{1}[-1]{$\symup{\Delta}$}}}}%
  \def\checkmark{\popcheck{popdark}}%
  \def\star{\mathbin{\ast}}\def\bigstar{\mathord{\ast}}%
  \def\diamond{\mathbin{\scalebox{0.6}{$\Diamond$}}}%
  \def\bigcup{\mathop{\mathchoice{\raisebox{-0.3em}{\scalebox{1.7}{$\cup$}}}{\scalebox{1.25}{$\cup$}}{\cup}{\cup}}\displaylimits}%
  \def\bigcap{\mathop{\mathchoice{\raisebox{-0.3em}{\scalebox{1.7}{$\cap$}}}{\scalebox{1.25}{$\cap$}}{\cap}{\cap}}\displaylimits}%
  \def\lesseqqgtr{\mathrel{\lessgtr}}\def\bigoplus{\mathop{\oplus}\displaylimits}%
}
\providecommand{\ssi}{si et seulement si} % abréviation fréquente
\AtBeginDocument{\providecommand{\wideparen}{\overparen}} % arc de cercle

%% FIGFIX-BEGIN v3 — correctifs globaux des figures (docs/onbuch-generation/tools/figfix)
\usepackage{adjustbox}\usepackage{etoolbox}
% toute figure TikZ/pgfplots plus large que la ligne est réduite à la largeur disponible
\BeforeBeginEnvironment{tikzpicture}{\begin{adjustbox}{max width=\linewidth}}
\AfterEndEnvironment{tikzpicture}{\end{adjustbox}}
% légendes pgfplots lisibles par-dessus les courbes
\pgfplotsset{every axis/.append style={legend style={fill=white,fill opacity=0.92,text opacity=1,draw=black!15,inner sep=3pt}}}
% étiquettes d'arêtes (syntaxe edge["..."]) avec fond blanc
\tikzset{every edge quotes/.append style={fill=white,inner sep=1.2pt}}
% bulles de cartes mentales : texte à la ligne dans la bulle, bulle un peu plus grande
\tikzset{every concept/.append style={text width=2.0cm,align=center,minimum size=2.3cm,inner sep=2pt}}
%% FIGFIX-END
\begin{document}

\pagegarde{Les systèmes politiques}{Allemand}{Tle A}{Chapitre 5 : Citoyenneté}{AL24}{2 h}{Cette leçon de type « texte » propose la lecture et le commentaire d'un texte original sur les systèmes politiques (démocratie, dictature, monarchie), l'acquisition du vocabulaire thématique (die Wahl, die Partei, die Verfassung...), ainsi que deux points de grammaire essentiels : le style nominal et les connecteurs logiques, avec le passif et le futur en conjugaison. Elle aboutit à une production guidée : résumé et court commentaire d'un texte politique.
\vspace{4pt}
\begin{multicols}{2}\small
\begin{itemize}
\item À la fin de cette leçon, je sais comprendre globalement puis en détail un texte allemand sur les systèmes politiques.
\item Je sais employer correctement le vocabulaire du thème (die Demokratie, die Diktatur, die Wahl, die Partei, der Präsident, die Verfassung) avec articles et pluriels.
\item Je sais reconnaître et former des familles de mots autour de wählen, herrschen, regieren.
\item Je sais identifier et produire le style nominal (Nominalstil) et le transformer en phrases verbales.
\item Je sais utiliser les connecteurs logiques (denn, weil, deshalb, aber, obwohl, zuerst... dann) avec la bonne place du verbe.
\item Je sais conjuguer au passif (Präsens et Präteritum) et au futur (Futur I).
\end{itemize}\end{multicols}}

\begin{sommaire}
\begin{multicols}{2}
\begin{enumerate}
\item Texte support : « Politische Systeme in der Welt »
\item Compréhension du texte
\item Lexique thématique : die politischen Systeme
\item Grammaire 1 : le style nominal (der Nominalstil)
\item Grammaire 2 : connecteurs logiques, passif et futur
\item Méthode du commentaire et production guidée
\item Activité d'intégration
\item Méthodes et exercices résolus
\item Exercices
\item Corrigés des exercices
\item Fiche bilan
\end{enumerate}
\end{multicols}
\end{sommaire}

\begin{objectifs}
\begin{itemize}
\item À la fin de cette leçon, je sais comprendre globalement puis en détail un texte allemand sur les systèmes politiques.
\item Je sais employer correctement le vocabulaire du thème (die Demokratie, die Diktatur, die Wahl, die Partei, der Präsident, die Verfassung) avec articles et pluriels.
\item Je sais reconnaître et former des familles de mots autour de wählen, herrschen, regieren.
\item Je sais identifier et produire le style nominal (Nominalstil) et le transformer en phrases verbales.
\item Je sais utiliser les connecteurs logiques (denn, weil, deshalb, aber, obwohl, zuerst... dann) avec la bonne place du verbe.
\item Je sais conjuguer au passif (Präsens et Präteritum) et au futur (Futur I).
\item Je sais répondre à des questions de compréhension en phrases complètes.
\item Je sais rédiger un résumé et un court commentaire structuré d'un texte politique.
\end{itemize}
\end{objectifs}

\begin{prerequis}
\begin{itemize}
\item Les déclinaisons de l'article et de l'adjectif (Nominatif, Accusatif, Datif) vues en Première.
\item La place du verbe : position 2 dans la phrase énonciative, verbe en fin dans la subordonnée.
\item Le Präsens et le Präteritum des verbes forts et faibles.
\item L'auxiliaire werden au Präsens (base du passif et du futur).
\item Le vocabulaire général de la vie en société vu en Première (der Staat, das Volk, die Regierung).
\end{itemize}
\end{prerequis}

\newpage

\coursec{Texte support : « Politische Systeme in der Welt »}

\courssub{Situation d'accroche et problématique}

En février, les Camerounais suivent avec passion les débats sur les élections et les institutions de leur pays : à la radio de Yaoundé, dans les journaux de Douala, dans les discussions au marché ou au carrefour des bendskins, on parle du président, des partis, de la constitution. En Allemagne, en Autriche et en Suisse, les citoyens votent aussi — mais dans des systèmes différents : république fédérale parlementaire à Berlin, démocratie directe à Berne, chancelière ou chancelier à Vienne. Comment dit-on « voter », « la constitution » ou « le président » en allemand ? Comment distinguer une démocratie d'une dictature dans un article de presse germanophone ?

\textbf{Problématique :} quelles clés linguistiques (vocabulaire, mots transparents, structures grammaticales) nous permettent de lire, comprendre et commenter un texte allemand sur les systèmes politiques ?

Ce thème est fréquent au Baccalauréat : il combine un lexique riche et régulier, des mots transparents faciles à repérer, et des structures typiques du texte journalistique (style nominal, connecteurs, passif) que nous étudierons dans les sections suivantes.

\courssub{Lecture globale : première approche du texte}

\begin{methode}[Lire un texte politique en trois étapes]
\begin{enumerate}
\item \textbf{Lecture globale sans dictionnaire.} Lisez le texte une première fois en silence, puis à voix haute. Posez-vous trois questions : \emph{qui ?} (les citoyens, le gouvernement, un monarque), \emph{quoi ?} (les systèmes politiques), \emph{où ?} (l'Allemagne, la Suisse, la Grande-Bretagne). Ne vous arrêtez pas sur les mots inconnus.
\item \textbf{Lecture détaillée avec glossaire.} Relisez phrase par phrase en utilisant le glossaire. Soulignez les mots du champ lexical politique et entourez les mots transparents.
\item \textbf{Reformulation.} Résumez l'idée principale en une seule phrase, en français puis en allemand : \all{Der Text beschreibt verschiedene politische Systeme und ihre Merkmale.} « Le texte décrit différents systèmes politiques et leurs caractéristiques. »
\end{enumerate}
\end{methode}

Avant de lire, observons le titre : \all{Politische Systeme in der Welt}. Deux mots transparents sautent aux yeux : \all{Systeme} ressemble à « systèmes » et \all{politische} à « politiques ». \all{Welt} (le monde) est un mot de base déjà connu. Nous pouvons donc prédire le contenu : un article d'information qui présente les différents régimes politiques dans le monde.

\begin{definition}[Ein politisches System]
\all{Ein politisches System ist die Art und Weise, wie ein Staat organisiert ist und wer die Macht ausübt.} — Un système politique est la manière dont un État est organisé et qui y exerce le pouvoir.
\begin{itemize}
\item \all{die Art und Weise} = la manière, la façon
\item \all{der Staat, -en} = l'État
\item \all{die Macht ausüben} = exercer le pouvoir
\end{itemize}
\end{definition}

\courssub{Le texte}

\begin{textebox}[Politische Systeme in der Welt]
Es gibt verschiedene politische Systeme in der Welt. Ein politisches System bestimmt, wie ein Staat organisiert ist und wer die Macht ausübt.

In einer Demokratie wählt das Volk seine Regierung in freien Wahlen. Alle Bürger haben das Recht zu wählen und gewählt zu werden. Deutschland ist zum Beispiel eine parlamentarische Demokratie : das Volk wählt den Bundestag, und der Bundestag wählt den Bundeskanzler. Die Verfassung, das Grundgesetz, garantiert die Rechte der Bürger und die Gewaltenteilung : die Regierung, das Parlament und die Gerichte kontrollieren sich gegenseitig. Auch die Schweiz ist eine Demokratie ; dort stimmen die Bürger mehrmals im Jahr direkt über wichtige Fragen ab.

In einer Diktatur dagegen hat eine Person oder eine Partei die ganze Macht. Es gibt keine freien Wahlen, und die Opposition wird unterdrückt. Die Presse ist nicht frei, und viele Menschen haben Angst, ihre Meinung zu sagen.

Eine Monarchie ist ein Staat mit einem König oder einer Königin an der Spitze. In einer absoluten Monarchie entscheidet der Monarch allein über alles. In einer parlamentarischen Monarchie wie Großbritannien dagegen hat der Monarch nur eine repräsentative Rolle : das gewählte Parlament und die Regierung regieren das Land.

In allen Systemen ist die Teilnahme der Bürger am politischen Leben wichtig für die Zukunft des Landes. Wer wählt, diskutiert und sich informiert, trägt zur Entwicklung seines Landes bei.
\end{textebox}

\textbf{Glossaire (mots difficiles) :}

\begin{center}
\begin{tabularx}{\linewidth}{|X|X|}
\hline
\rowcolor{popblueL} \textbf{Deutsch} & \textbf{Français} \\
\hline
\all{die Macht ausüben} & exercer le pouvoir \\
\hline
\all{das Volk, die Völker} & le peuple \\
\hline
\all{die Regierung, -en} & le gouvernement \\
\hline
\all{die Wahl, -en / wählen} & l'élection / élire, voter \\
\hline
\all{das Recht, -e} & le droit \\
\hline
\all{das Grundgesetz} & la Loi fondamentale (constitution allemande) \\
\hline
\all{die Gewaltenteilung} & la séparation des pouvoirs \\
\hline
\all{das Gericht, -e} & le tribunal \\
\hline
\all{sich gegenseitig} & mutuellement, l'un l'autre \\
\hline
\all{abstimmen über + Akk.} & voter sur, se prononcer sur \\
\hline
\all{unterdrücken} & opprimer, réprimer \\
\hline
\all{die Meinung, -en} & l'opinion, l'avis \\
\hline
\all{an der Spitze} & à la tête (de l'État) \\
\hline
\all{regieren} & gouverner, régner \\
\hline
\all{die Teilnahme an + Dat.} & la participation à \\
\hline
\all{beitragen zu + Dat. (trägt bei, trug bei, hat beigetragen)} & contribuer à \\
\hline
\end{tabularx}
\end{center}

\courssub{Compréhension globale : type de texte, source, idée générale}

\begin{itemize}
\item \textbf{Type de texte :} un article de presse à visée informative et explicative (un \all{Sachtext}, texte factuel). Il n'y a ni dialogue, ni récit : l'auteur définit, compare et classe.
\item \textbf{Source probable :} un article de journal ou de magazine pour jeunes lecteurs (type magazine scolaire allemand), car le ton est pédagogique et les exemples sont simples.
\item \textbf{Idée générale :} il existe plusieurs systèmes politiques dans le monde — démocratie, dictature, monarchie — et dans tous les cas la participation des citoyens est essentielle.
\end{itemize}

\textbf{Le plan du texte} (à retrouver dans vos cahiers) :

\begin{enumerate}
\item Définition d'un système politique (phrases 1 à 2).
\item La démocratie, avec les exemples de l'Allemagne et de la Suisse (paragraphe 2).
\item La dictature (paragraphe 3).
\item La monarchie : absolue contre parlementaire (paragraphe 4).
\item Conclusion : l'importance de la participation des citoyens (dernier paragraphe).
\end{enumerate}

\begin{exoresolu}[Questions de compréhension globale]
\textbf{Énoncé.} Répondez en allemand par des phrases complètes :
\begin{enumerate}
\item \all{Was ist ein politisches System?}
\item \all{Wer wählt den Bundeskanzler in Deutschland?}
\item \all{Was passiert mit der Opposition in einer Diktatur?}
\item \all{Welche Rolle hat der Monarch in einer parlamentarischen Monarchie?}
\end{enumerate}
\tcblower
\textbf{Solution détaillée.}
\begin{enumerate}
\item \all{Ein politisches System bestimmt, wie ein Staat organisiert ist und wer die Macht ausübt.} « Un système politique détermine comment un État est organisé et qui exerce le pouvoir. » (Reprise de la définition du premier paragraphe.)
\item \all{Der Bundestag wählt den Bundeskanzler.} « Le Bundestag élit le chancelier fédéral. » Attention : ce n'est pas le peuple directement — c'est le sens de « démocratie \emph{parlementaire} ».
\item \all{Die Opposition wird unterdrückt.} « L'opposition est opprimée. » On relève ici un passif (\all{wird unterdrückt}), que nous étudierons en grammaire.
\item \all{Er hat nur eine repräsentative Rolle ; das Parlament und die Regierung regieren das Land.} « Il n'a qu'un rôle représentatif ; le parlement et le gouvernement gouvernent le pays. »
\end{enumerate}
\end{exoresolu}

\courssub{Lecture détaillée : mots transparents et champ lexical politique}

\textbf{Stratégie n\textsuperscript{o} 1 : les mots transparents.} Ce sont des mots dont la forme allemande ressemble au français ; ils viennent souvent du latin ou du grec. Dans ce texte, on peut deviner sans dictionnaire :

\begin{center}
\begin{tabularx}{\linewidth}{|X|X|X|}
\hline
\rowcolor{popblueL} \textbf{Deutsch} & \textbf{Français} & \textbf{Remarque} \\
\hline
\all{die Demokratie} & la démocratie & féminin, comme en français \\
\hline
\all{die Diktatur} & la dictature & féminin \\
\hline
\all{die Monarchie} & la monarchie & féminin \\
\hline
\all{die Partei} & le parti & attention au genre : \all{die} ! \\
\hline
\all{der Präsident} & le président & masculin ; la présidente : \all{die Präsidentin} \\
\hline
\all{das System} & le système & neutre, contrairement au français \\
\hline
\all{die Opposition} & l'opposition & féminin \\
\hline
\all{die Verfassung} & la constitution & moins transparent, à mémoriser \\
\hline
\all{repräsentativ} & représentatif & adjectif transparent \\
\hline
\all{parlamentarisch} & parlementaire & adjectif transparent \\
\hline
\end{tabularx}
\end{center}

\textbf{Stratégie n\textsuperscript{o} 2 : le contexte.} Si vous ne connaissez pas \all{unterdrücken}, lisez la phrase entière : \all{Es gibt keine freien Wahlen, und die Opposition wird unterdrückt.} « Il n'y a pas d'élections libres, et l'opposition est... » Le contexte négatif (pas d'élections libres) permet de deviner : opprimée, réprimée.

\textbf{Stratégie n\textsuperscript{o} 3 : les connaissances du monde.} Vous savez que la Grande-Bretagne a un roi et un parlement : vous pouvez donc comprendre \all{parlamentarische Monarchie} même sans glossaire.

\begin{attention}[Faux amis et pièges de genre]
\begin{itemize}
\item \all{das Volk} signifie « le peuple », pas « la foule » ; \all{das Volk wählt} = « le peuple élit ».
\item \all{die Wahl} = l'élection, mais aussi « le choix » ; le verbe \all{wählen} = élire, voter, choisir.
\item \all{die Macht} = le pouvoir (politique) ; ne pas confondre avec \all{die Nacht} (la nuit) !
\item En français « le système » est masculin ; en allemand c'est \all{das System} (neutre). Apprenez toujours le nom avec son article.
\item \all{herrschen} = régner, dominer ; \all{der Herrscher} = le souverain, le dirigeant.
\end{itemize}
\end{attention}

\courssub{Carte mentale du texte}

\begin{popfigure}
\begin{tikzpicture}[mindmap, grow cyclic, every node/.style={concept}, root concept/.append style={concept color=popblue, font=\bfseries, text=white}, level 1/.append style={level distance=3.4cm, sibling angle=120}, level 2/.append style={level distance=2.4cm, sibling angle=50}]
\node[root concept]{Politische Systeme}
  child[concept color=popgreen] { node {Demokratie} 
      child { node[concept color=popgreenL, text=popdark, align=center]{freie Wahlen} }
      child { node[concept color=popgreenL, text=popdark, align=center]{Gewaltenteilung} } }
  child[concept color=poporange] { node {Diktatur} 
      child { node[concept color=poporangeL, text=popdark, align=center]{eine Partei an der Macht} }
      child { node[concept color=poporangeL, text=popdark, align=center]{Opposition unterdrückt} } }
  child[concept color=poppurple] { node {Monarchie} 
      child { node[concept color=poppurpleL, text=popdark, align=center]{König an der Spitze} }
      child { node[concept color=poppurpleL, text=popdark, align=center]{absolut oder parlamentarisch} } };
\end{tikzpicture}
\legende{Carte mentale : les trois systèmes politiques du texte et leurs caractéristiques.}
\end{popfigure}

\courssub{Relevé grammatical : préparer les sections suivantes}

Avant d'étudier la grammaire en détail, relevons dans le texte les structures que nous analyserons dans les sections 4 et 5. C'est une démarche inductive : observer d'abord, expliquer ensuite.

\begin{center}
\begin{tabularx}{\linewidth}{|X|X|X|}
\hline
\rowcolor{popblueL} \textbf{Structure} & \textbf{Exemple relevé dans le texte} & \textbf{Traduction} \\
\hline
Style nominal & \all{die Teilnahme der Bürger am politischen Leben} & la participation des citoyens à la vie politique \\
\hline
Style nominal & \all{die Gewaltenteilung} & la séparation des pouvoirs \\
\hline
Connecteur d'opposition & \all{In einer Diktatur dagegen...} & En revanche, dans une dictature... \\
\hline
Connecteur de conséquence & \all{...und viele Menschen haben Angst, ihre Meinung zu sagen} & ...et beaucoup de gens ont peur de dire leur avis \\
\hline
Passif (Präsens) & \all{die Opposition wird unterdrückt} & l'opposition est opprimée \\
\hline
Passif (infinitif) & \all{gewählt zu werden} & être élu \\
\hline
\end{tabularx}
\end{center}

\begin{aretenir}[Ce qu'il faut retenir de cette première lecture]
\begin{itemize}
\item Un \all{Sachtext} politique se lit en trois étapes : globale, détaillée, reformulation.
\item Les mots transparents (\all{die Demokratie, die Partei, der Präsident, das System}) donnent immédiatement le thème du texte.
\item Le contexte et les connaissances du monde permettent de deviner les mots inconnus (\all{unterdrücken, Gewaltenteilung}).
\item Le texte contient déjà les trois cibles grammaticales de la leçon : style nominal, connecteurs logiques et passif.
\end{itemize}
\end{aretenir}

\begin{savaistu}[La démocratie directe en Suisse]
La Suisse pratique la \all{direkte Demokratie} : les citoyens votent plusieurs fois par an sur des lois précises, lors de la \all{Volksabstimmung} (le référendum, littéralement « le vote du peuple » : \all{das Volk} + \all{die Abstimmung}, le vote). Ils peuvent aussi lancer une \all{Volksinitiative} pour proposer eux-mêmes une modification de la constitution. C'est un exemple célèbre de \all{Teilnahme der Bürger am politischen Leben} — exactement l'idée de la conclusion de notre texte !
\end{savaistu}

\begin{exercice}[À vous de jouer]
\begin{enumerate}
\item Retrouvez dans le texte cinq mots du champ lexical politique et donnez leur article et leur traduction.
\item Traduisez en français : \all{Die Verfassung garantiert die Rechte der Bürger.}
\item Devinez le sens de \all{die Presse ist nicht frei} grâce au contexte du paragraphe sur la dictature.
\item Résumez le texte en une phrase allemande commençant par \all{Der Text beschreibt...}
\end{enumerate}
\end{exercice}

\begin{corrige}[Exercice « À vous de jouer »]
\begin{enumerate}
\item Par exemple : \all{die Regierung} (le gouvernement), \all{die Wahl} (l'élection), \all{das Volk} (le peuple), \all{der König} (le roi), \all{das Parlament} (le parlement).
\item « La constitution garantit les droits des citoyens. »
\item « La presse n'est pas libre » : dans une dictature, l'opposition est opprimée, donc les journaux ne peuvent pas écrire librement.
\item \all{Der Text beschreibt verschiedene politische Systeme — Demokratie, Diktatur und Monarchie — und die wichtige Rolle der Bürger.} « Le texte décrit différents systèmes politiques — démocratie, dictature et monarchie — et le rôle important des citoyens. »
\end{enumerate}
\end{corrige}

\coursec{Compréhension du texte}

Après la lecture attentive du texte support \emph{« Politische Systeme in der Welt »} (section 1), nous passons maintenant à l'exploitation didactique de ce document. La compréhension de l'écrit est une compétence centrale du Bac : il ne s'agit pas seulement de « comprendre le sens général », mais de \textbf{repérer des informations précises}, de \textbf{reformuler avec ses propres mots} et de \textbf{justifier ses réponses par le texte}. Cette section vous guide pas à pas, du global au détaillé, jusqu'à une première expression personnelle argumentée.

\courssub{Le texte de référence (rappel avec numérotation des lignes)}

Pour faciliter le repérage (citations, numéros de ligne), voici le texte intégral découpé en segments numérotés. Gardez-le sous les yeux pour tous les exercices qui suivent.

\begin{textebox}[Politische Systeme in der Welt — Texte support numéroté]
(1) In der modernen Welt gibt es verschiedene politische Systeme, die das Zusammenleben der Menschen regeln. 
(2) Die zwei Hauptformen sind die Demokratie und die Diktatur. 
(3) In einer Demokratie geht die Staatsgewalt vom Volk aus. 
(4) Ein wichtiges Beispiel ist die Bundesrepublik Deutschland: Sie ist eine parlamentarische Demokratie. 
(5) Das Volk wählt alle vier Jahre den Bundestag. 
(6) Der Bundestag wählt dann den Bundeskanzler oder die Bundeskanzlerin. 
(7) Das Grundgesetz garantiert die Grundrechte wie Meinungsfreiheit, Pressefreiheit und Versammlungsfreiheit. 
(8) Eine Gewaltenteilung zwischen Legislative, Exekutive und Judikative sichert die Machtkontrolle. 
(9) Ganz anders ist die Situation in einer Diktatur. 
(10) Dort herrscht eine einzelne Person oder eine kleine Gruppe ohne Kontrolle durch das Volk. 
(11) Es gibt keine freien Wahlen, und die Opposition wird oft unterdrückt oder ins Gefängnis gesteckt. 
(12) Die Medien sind gleichgeschaltet und dürfen nicht kritisch berichten. 
(13) Ein drittes System ist die parlamentarische Monarchie, wie in Großbritannien, Schweden oder Spanien. 
(14) Dort gibt es einen Monarchen (König oder Königin) als Staatsoberhaupt. 
(15) Seine Rolle ist jedoch meist nur repräsentativ und zeremoniell. 
(16) Die eigentliche politische Macht liegt beim gewählten Parlament und der Regierung. 
(17) Der Monarch muss neutral bleiben und darf nicht in die Tagespolitik eingreifen. 
(18) Zusammenfassend lässt sich sagen: Eine Verfassung und die Trennung der Gewalten sind die Grundpfeiler 
(19) jeder stabilen Demokratie, während Willkür und Machtmissbrauch die Diktatur kennzeichnen.
\end{textebox}

\begin{aretenir}[Glossaire rapide du texte (mots difficiles)]
\begin{itemize}
\item \all{die Staatsgewalt, -en} : le pouvoir d'État ; \all{das Zusammenleben} : la vie en société.
\item \all{das Grundgesetz} : la Loi fondamentale (constitution allemande) ; \all{das Grundrecht, -e} : le droit fondamental.
\item \all{die Gewaltenteilung} : la séparation des pouvoirs ; \all{die Machtkontrolle, -n} : le contrôle du pouvoir.
\item \all{gleichgeschaltet} : mis au pas, aligné ; \all{unterdrücken} : réprimer, opprimer.
\item \all{das Staatsoberhaupt, -häupter} : le chef de l'État ; \all{die Willkür} : l'arbitraire ; \all{der Machtmissbrauch} : l'abus de pouvoir.
\item \all{eingreifen (griff ein, hat eingegriffen)} : intervenir ; \all{die Tagespolitik} : la politique quotidienne.
\end{itemize}
\end{aretenir}

\courssub{Compréhension globale (Gesamtverständnis)}

Avant d'entrer dans le détail, vérifiez que vous avez saisi la structure d'ensemble et le message principal. Répondez en français, en phrases complètes.

\begin{exercice}[Questions globales]
\begin{enumerate}
\item De quoi parle ce texte ? Donne le thème général et la thèse défendue.
\item Quels sont les trois systèmes politiques présentés ? Cite-les dans l'ordre du texte.
\item Quel système est présenté comme le modèle de référence pour la stabilité ? Justifie par une phrase du texte (numéro de ligne).
\end{enumerate}
\end{exercice}

\begin{corrige}[Exercice : Compréhension globale]
\begin{enumerate}
\item \textbf{Thème :} Le texte présente et compare les principaux systèmes politiques dans le monde moderne. \textbf{Thèse :} la démocratie (avec constitution et séparation des pouvoirs) est le garant de la liberté et de la stabilité, contrairement à la dictature caractérisée par l'arbitraire.
\item Les trois systèmes sont : 1. la démocratie (parlementaire, ex. Allemagne), 2. la dictature, 3. la monarchie parlementaire (ex. Royaume-Uni, Suède, Espagne).
\item C'est la démocratie. Justification (l. 18-19) : \all{„Eine Verfassung und die Trennung der Gewalten sind die Grundpfeiler jeder stabilen Demokratie…“} (« Une constitution et la séparation des pouvoirs sont les piliers de toute démocratie stable… »).
\end{enumerate}
\end{corrige}

\courssub{Compréhension détaillée (Detailverständnis) — Fragen auf Deutsch}

Au Bac, les questions de compréhension détaillée sont souvent posées \textbf{en allemand}. Vous devez y répondre \textbf{en allemand}, en phrases complètes (sujet + verbe), en reprenant le vocabulaire du texte mais sans le recopier servilement.

\begin{methode}[Répondre à une question de compréhension écrite (Bac standard)]
\begin{enumerate}
\item \textbf{Lire la question} et identifier le \cle{mot-clé} (Wer, Was, Warum, Welche Rolle, Wie…).
\item \textbf{Repérer la zone} du texte : chercher le mot-clé ou ses synonymes (ex. \all{Bundeskanzler} $\rightarrow$ l. 6 ; \all{Grundgesetz} $\rightarrow$ l. 7).
\item \textbf{Lire la phrase concernée} et la phrase d'avant / d'après (contexte).
\item \textbf{Formuler la réponse} : phrase complète, sujet + verbe conjugué, \emph{propres mots} (reformulation) si possible, sinon citation encadrée de guillemets allemands „…“.
\item \textbf{Vérifier} : la réponse répond-elle exactement à la question ? L'orthographe (majuscules des noms, Umlaute) est-elle correcte ?
\end{enumerate}
\end{methode}

\begin{exemplebox}[Modèles de réponses détaillées (Fragen \& Antworten)]
\textbf{Frage 1 :} \all{Wer wählt in Deutschland den Bundeskanzler?} \\
\textbf{Antwort :} \all{Der Bundestag wählt den Bundeskanzler.} (l. 6) — « Le Bundestag élit le Chancelier. »

\textbf{Frage 2 :} \all{Was garantiert das Grundgesetz?} \\
\textbf{Antwort :} \all{Das Grundgesetz garantiert die Grundrechte wie Meinungsfreiheit, Pressefreiheit und Versammlungsfreiheit.} (l. 7) — « La Loi fondamentale garantit les droits fondamentaux comme la liberté d'opinion, la liberté de la presse et la liberté de réunion. »

\textbf{Frage 3 :} \all{Was passiert mit der Opposition in einer Diktatur?} \\
\textbf{Antwort :} \all{In einer Diktatur wird die Opposition oft unterdrückt oder ins Gefängnis gesteckt.} (l. 11) — « Dans une dictature, l'opposition est souvent réprimée ou jetée en prison. »

\textbf{Frage 4 :} \all{Welche Rolle hat der Monarch in einer parlamentarischen Monarchie?} \\
\textbf{Antwort :} \all{Die Rolle des Monarchen ist nur repräsentativ und zeremoniell; die politische Macht liegt beim Parlament und bei der Regierung.} (l. 15-16) — « Le rôle du monarque est seulement représentatif et cérémoniel ; le pouvoir politique revient au Parlement et au gouvernement. »
\end{exemplebox}

\begin{attention}[Piège Bac : la réponse en un seul mot vaut zéro]
\begin{itemize}
\item \textbf{Interdit :} \all{Der Bundestag.} / \all{Grundrechte.} / \all{Repräsentativ.}
\item \textbf{Exigé :} une phrase complète avec sujet et verbe conjugué, \textbf{même pour une question « Wer / Was »}.
\item \textbf{Francophones, attention :} en français, on répond souvent « Le Bundestag. ». En allemand au Bac, il faut écrire : \all{Der Bundestag wählt den Bundeskanzler.}
\item \textbf{Astuce :} reprenez le sujet de la question (ou le pronom correspondant) + verbe + complément.
\end{itemize}
\end{attention}

\begin{exercice}[Compréhension détaillée — À vous de jouer !]
Répondez en allemand, en phrases complètes, en vous basant \textbf{uniquement} sur le texte.
\begin{enumerate}
\item \all{Wie oft wählt das Volk den Bundestag?}
\item \all{Welche drei Grundrechte nennt der Text explizit?}
\item \all{Wer herrscht in einer Diktatur?}
\item \all{Wie sind die Medien in einer Diktatur gesteuert?}
\item \all{Nenne zwei Länder, die als Beispiel für eine parlamentarische Monarchie dienen.}
\item \all{Warum muss der Monarch neutral bleiben?}
\end{enumerate}
\end{exercice}

\begin{corrige}[Exercice : Compréhension détaillée]
\begin{enumerate}
\item \all{Das Volk wählt den Bundestag alle vier Jahre.} (l. 5)
\item \all{Der Text nennt die Meinungsfreiheit, die Pressefreiheit und die Versammlungsfreiheit.} (l. 7)
\item \all{In einer Diktatur herrscht eine einzelne Person oder eine kleine Gruppe.} (l. 10)
\item \all{Die Medien sind gleichgeschaltet und dürfen nicht kritisch berichten.} (l. 12)
\item \all{Großbritannien, Schweden und Spanien werden als Beispiele genannt.} (l. 13)
\item \all{Er muss neutral bleiben, weil er nicht in die Tagespolitik eingreifen darf.} (l. 17)
\end{enumerate}
\end{corrige}

\courssub{Exercice Vrai / Faux — Richtig / Falsch mit Begründung}

Cet exercice type Bac teste la lecture fine. \textbf{Chaque réponse doit être justifiée par une citation précise du texte (numéro de ligne + extrait entre guillemets allemands).}

\begin{exercice}[Richtig oder Falsch? Begründen Sie mit Zeilennachweis!]
\begin{enumerate}
\item \all{In Deutschland wählt das Volk den Bundeskanzler direkt.}
\item \all{Das Grundgesetz sichert die Gewaltenteilung zu.}
\item \all{In einer Diktatur gibt es Wahlen, aber sie sind nicht frei.}
\item \all{Der Monarch in einer parlamentarischen Monarchie regiert das Land aktiv.}
\item \all{Die Trennung der Gewalten ist nur in einer Monarchie wichtig.}
\end{enumerate}
\end{exercice}

\begin{exoresolu}[Corrigé détaillé : Vrai / Faux]
\begin{enumerate}
\item \textbf{Falsch.} Z. 5-6 : \all{„Das Volk wählt alle vier Jahre den Bundestag. Der Bundestag wählt dann den Bundeskanzler…“} $\rightarrow$ Le peuple élit le Bundestag, qui élit ensuite le Chancelier : il s'agit d'une élection \emph{indirecte}.
\item \textbf{Richtig.} Z. 7-8 : \all{„Das Grundgesetz garantiert die Grundrechte… Eine Gewaltenteilung zwischen Legislative, Exekutive und Judikative sichert die Machtkontrolle.“} $\rightarrow$ Le texte lie la séparation des pouvoirs au cadre démocratique garanti par le Grundgesetz.
\item \textbf{Falsch.} Z. 11 : \all{„Es gibt keine freien Wahlen…“} $\rightarrow$ Le texte dit qu'il n'y a \textbf{pas} d'élections libres ; il ne dit pas qu'il existe des élections non libres.
\item \textbf{Falsch.} Z. 15-16 : \all{„Seine Rolle ist jedoch meist nur repräsentativ und zeremoniell. Die eigentliche politische Macht liegt beim gewählten Parlament…“} $\rightarrow$ Rôle cérémonial, pas de pouvoir politique réel.
\item \textbf{Falsch.} Z. 18-19 : \all{„Eine Verfassung und die Trennung der Gewalten sind die Grundpfeiler jeder stabilen Demokratie…“} $\rightarrow$ C'est le pilier de la \textbf{démocratie}, pas de la monarchie.
\end{enumerate}
\end{exoresolu}

\courssub{Reformulation et résumé guidé (Wortschatzarbeit)}

Savoir reformuler sans copier-coller est la clé de la \emph{Zusammenfassung} (résumé) et de l'\emph{Übersetzung} (traduction). Complétez le résumé ci-dessous avec \textbf{un seul mot} (ou une forme fléchie) \textbf{pris exactement dans le texte} (respectez la casse et l'orthographe).

\begin{exercice}[Lückentext — Résumé structuré]
\all{Der Text vergleicht drei politische \_\_\_\_\_\_\_\_\_\_ (1). In einer \_\_\_\_\_\_\_\_\_\_ (2) wie Deutschland geht die Gewalt vom \_\_\_\_\_\_\_\_\_\_ (3) aus. Das \_\_\_\_\_\_\_\_\_\_ (4) wählt den Bundestag, dieser wählt den \_\_\_\_\_\_\_\_\_\_ (5). Das \_\_\_\_\_\_\_\_\_\_ (6) schützt die Grundrechte. In einer \_\_\_\_\_\_\_\_\_\_ (7) gibt es keine freien \_\_\_\_\_\_\_\_\_\_ (8), die \_\_\_\_\_\_\_\_\_\_ (9) wird unterdrückt. Eine \_\_\_\_\_\_\_\_\_\_ (10) Monarchie hat einen Monarchen als \_\_\_\_\_\_\_\_\_\_ (11), aber die Macht liegt beim \_\_\_\_\_\_\_\_\_\_ (12). Wichtig für Stabilität sind \_\_\_\_\_\_\_\_\_\_ (13) und Gewaltenteilung.}
\end{exercice}

\begin{corrige}[Exercice : Reformulation (Lückentext)]
(1) \all{Systeme} (l. 1) — (2) \all{Demokratie} (l. 2/4) — (3) \all{Volk} (l. 3) — (4) \all{Volk} (l. 5) — (5) \all{Bundeskanzler} (l. 6) — (6) \all{Grundgesetz} (l. 7) — (7) \all{Diktatur} (l. 9) — (8) \all{Wahlen} (l. 11) — (9) \all{Opposition} (l. 11) — (10) \all{parlamentarische} (l. 13) — (11) \all{Staatsoberhaupt} (l. 14) — (12) \all{Parlament} (l. 16) — (13) \all{Verfassung} (l. 18).
\end{corrige}

\begin{aretenir}[Vocabulaire clé du résumé (familles de mots)]
\begin{center}
\begin{tabularx}{\linewidth}{|X|X|X|}
\hline
\rowcolor{popblueL} \textbf{Nom (article + pluriel)} & \textbf{Verbe associé} & \textbf{Adjectif / Participe} \\ \hline
\all{das System, die Systeme} & \all{systematisieren} & \all{systematisch} \\ \hline
\all{die Demokratie, -n} & \all{demokratisieren} & \all{demokratisch} \\ \hline
\all{die Diktatur, -en} & \all{diktieren} & \all{diktatorisch} \\ \hline
\all{die Wahl, -en} & \all{wählen (wählte, hat gewählt)} & \all{gewählt} \\ \hline
\all{das Grundgesetz, die Grundgesetze} & — & \all{grundgesetzlich} \\ \hline
\all{die Gewaltenteilung, -en} & \all{die Gewalten teilen} & \all{geteilt} \\ \hline
\all{die Monarchie, -n} & \all{regieren} & \all{monarchisch} \\ \hline
\all{die Opposition, -en} & \all{opponieren} & \all{oppositionell} \\ \hline
\all{die Verfassung, -en} & \all{verfassen} & \all{verfassungsmäßig} \\ \hline
\all{der Präsident, -en / die Präsidentin, -nen} & \all{präsidieren} & \all{präsidial} \\ \hline
\all{die Partei, -en} & \all{sich einer Partei anschließen} & \all{parteilich / parteilos} \\ \hline
\end{tabularx}
\end{center}
\end{aretenir}

\begin{attention}[Faux amis et pièges lexicaux du thème politique]
\begin{itemize}
\item \all{die Gewalt} signifie « la violence » \textbf{et} « le pouvoir » (dans \all{die Staatsgewalt}, \all{die Gewaltenteilung}) : ne traduisez pas toujours par « violence » !
\item \all{die Wahl} = « l'élection » ; « le choix » se dit aussi \all{die Wahl} (\all{die Wahl haben} = avoir le choix), mais « voter » = \all{wählen} ou \all{wählen gehen}.
\item \all{die Partei} = « le parti » (politique) ; « la partie » (d'un jeu) = \all{die Partie}.
\item \all{herrschen} = « régner, dominer » (sans complément d'objet direct : \all{über ein Land herrschen}) ; ne pas confondre avec \all{regieren} (« gouverner », avec accusatif possible).
\item Majuscules obligatoires : \all{das Volk}, \all{die Opposition}, \all{die Medien} (pluriel de \all{das Medium}).
\end{itemize}
\end{attention}

\courssub{Question d'ouverture : expression personnelle guidée}

Cette question prépare l'expression orale et écrite (dernière partie de l'épreuve). Vous devez donner votre avis \textbf{en allemand}, en justifiant avec \textbf{\cle{weil} (subordonnée, verbe conjugué à la fin) ou \cle{denn} (coordination, verbe en 2e position)}.

\begin{methode}[Argumenter son choix : \textit{weil} vs \textit{denn}]
\begin{enumerate}
\item Choisissez votre système : \all{die Demokratie}, \all{die parlamentarische Monarchie} (évitez de défendre la dictature !).
\item Construisez deux arguments (liberté, stabilité, participation, contrôle du pouvoir).
\item \textbf{Structure \textit{weil} (subordonnée) :} \all{Ich finde die Demokratie besser, weil das Volk die Regierung wählen kann.} — le verbe conjugué \all{kann} va à la \textbf{fin}.
\item \textbf{Structure \textit{denn} (coordination) :} \all{Ich finde die Demokratie besser, denn das Volk kann die Regierung wählen.} — le verbe conjugué \all{kann} reste en \textbf{2e position} après \all{denn}.
\item \textbf{Conseil :} alternez \textit{weil} et \textit{denn} pour montrer au correcteur la maîtrise des deux constructions.
\end{enumerate}
\end{methode}

\begin{exemplebox}[Modèles de réponse à la question d'ouverture]
\textbf{Frage :} \all{Welches System findest du besser? Warum?}

\textbf{Antwort A (mit \textit{weil}) :} \all{Ich finde die parlamentarische Demokratie am besten, weil die Bürger ihre Vertreter frei wählen können und die Grundrechte durch das Grundgesetz geschützt sind. Außerdem gibt es eine Gewaltenteilung, die Machtmissbrauch verhindert.}

\textbf{Antwort B (mit \textit{denn}) :} \all{Die Demokratie ist besser als die Diktatur, denn in einer Diktatur gibt es keine Meinungsfreiheit, und die Opposition wird verfolgt. In einer Demokratie kann man seine Meinung sagen, ohne Angst zu haben.}
\end{exemplebox}

\begin{exercice}[À votre tour ! Écrivez 2 à 3 phrases en allemand.]
\all{Welches politische System (Demokratie, Diktatur, parlamentarische Monarchie) findest du am besten? Begründe deine Meinung mit zwei Argumenten.} Utilisez \textbf{une fois \textit{weil} et une fois \textit{denn}}.
\end{exercice}

\begin{corrige}[Exercice : Expression personnelle (proposition de corrigé)]
\all{Ich bevorzuge die Demokratie, weil sie die Menschenrechte und die Pressefreiheit garantiert. Sie ist mir am liebsten, denn nur hier kann die Opposition legal arbeiten und die Regierung kontrollieren.}
\end{corrige}

\courssub{Stratégie de repérage : tableau « Frage / Antwortstrategie »}

Pour automatiser le réflexe « mot-clé $\rightarrow$ zone du texte », voici une carte mentale opératoire. Apprenez-la par cœur pour gagner du temps le jour de l'examen.

\begin{popfigure}
\begin{tikzpicture}[
  node distance=8mm and 12mm,
  main/.style={draw=popblue, thick, rounded corners, fill=popblueL, text width=3.2cm, align=center, font=\small},
  sub/.style={draw=poporange, thick, rounded corners, fill=poporangeL, text width=3.5cm, align=center, font=\small\itshape},
  strat/.style={draw=popgreen, thick, rounded corners, fill=popgreenL, text width=4cm, align=center, font=\small},
  arrow/.style={-Stealth, thick, popdark}
]
% Central node
\node[main, align=center] (center){\textbf{FRAGE-WÖRTER}\\(Mots-clés de la question)};

% Branches
\node[sub, above left=of center, align=center] (wer){\textbf{Wer?} (Qui ?)\\Personne, acteur};
\node[strat, left=of wer] (werS) {Chercher : le sujet de la phrase, noms propres, pronoms. Ex. : \all{Der Bundestag}, \all{Das Volk}};

\node[sub, left=of center, align=center] (was){\textbf{Was?} (Quoi ?)\\Action, objet, concept};
\node[strat, left=of was] (wasS) {Chercher : verbe + accusatif, noms abstraits (\all{Grundrechte}, \all{Macht})};

\node[sub, below left=of center, align=center] (warum){\textbf{Warum? / Weshalb?}\\Cause, raison};
\node[strat, left=of warum] (warumS) {Chercher : \all{weil}, \all{denn}, \all{damit} ; structure cause $\rightarrow$ effet};

\node[sub, above right=of center, align=center] (wie){\textbf{Wie? / Welche Rolle?}\\Manière, fonction};
\node[strat, right=of wie] (wieS) {Chercher : adverbes, verbes modaux, fonctions (\all{repräsentativ}, \all{regieren})};

\node[sub, right=of center, align=center] (wo){\textbf{Wo? / In welchem Land?}\\Lieu, exemple};
\node[strat, right=of wo] (woS) {Chercher : noms de pays. Ex. : \all{Deutschland}, \all{Großbritannien}};

\node[sub, below right=of center, align=center] (wann){\textbf{Wann? / Wie oft?}\\Temps, fréquence};
\node[strat, right=of wann] (wannS) {Chercher : expressions de temps. Ex. : \all{alle vier Jahre}, \all{meist}};

% Arrows
\draw[arrow] (center) -- (wer); \draw[arrow] (wer) -- (werS);
\draw[arrow] (center) -- (was); \draw[arrow] (was) -- (wasS);
\draw[arrow] (center) -- (warum); \draw[arrow] (warum) -- (warumS);
\draw[arrow] (center) -- (wie); \draw[arrow] (wie) -- (wieS);
\draw[arrow] (center) -- (wo); \draw[arrow] (wo) -- (woS);
\draw[arrow] (center) -- (wann); \draw[arrow] (wann) -- (wannS);
\end{tikzpicture}
\legende{Stratégie de repérage : associer le mot-clé de la question (Fragewort) à la zone syntaxique du texte où se trouve la réponse (Antwortstrategie).}
\end{popfigure}

\begin{savaistu}[Le saviez-vous ? — Le \textit{Grundgesetz} allemand]
Le \all{Grundgesetz} (Loi fondamentale) de 1949 n'était au départ qu'une constitution \textbf{provisoire} pour l'Allemagne de l'Ouest — d'où le nom \emph{Grundgesetz} (« loi de base ») et non \emph{Verfassung} (« constitution »). Il devait durer jusqu'à la réunification : son article 146 prévoyait qu'il perdrait sa validité le jour où le peuple allemand se donnerait une constitution par une décision libre. Lors de la réunification de 1990, on a choisi de \textbf{conserver le Grundgesetz} (avec l'adhésion des Länder de l'Est via l'article 23) plutôt que d'en rédiger un nouveau. C'est pourquoi l'Allemagne est encore régie aujourd'hui par un texte qui porte le nom de « loi fondamentale » — un cas unique au monde !
\end{savaistu}

\coursec{Lexique thématique : die politischen Systeme}

Après avoir compris le texte \all{„Politische Systeme in der Welt“} dans son ensemble, nous allons maintenant nous arrêter sur ses mots-clés. Un texte sur les systèmes politiques ne peut être ni compris finement ni résumé correctement sans maîtriser le vocabulaire spécialisé : les noms des régimes, des institutions, des acteurs politiques et des élections. Dans cette section, nous allons observer ces mots dans des phrases authentiques, les classer par champs lexicaux, apprendre systématiquement leur \cle{article} et leur \cle{pluriel}, explorer leurs \cle{familles de mots} et mémoriser les \cle{expressions figées} qui reviennent dans tous les sujets d'examen.

\begin{definition}[Qu'est-ce qu'un champ lexical ?]
Un \cle{champ lexical} est l'ensemble des mots qui se rapportent à une même notion ou à un même domaine de la réalité. Ici, le domaine est celui de la \all{Politik} (la politique). Apprendre le vocabulaire par champs — les systèmes, les institutions, les élections — est bien plus efficace qu'une liste désordonnée : les mots se renforcent mutuellement en mémoire.
\end{definition}

\courssub{Champ 1 : les systèmes politiques}

Observons d'abord ces phrases, tirées du texte support et de la presse germanophone :

\all{Deutschland ist eine Demokratie und eine Bundesrepublik.} « L'Allemagne est une démocratie et une république fédérale. »

\all{In einer Diktatur hat nur eine Person oder eine kleine Gruppe die Macht.} « Dans une dictature, une seule personne ou un petit groupe détient le pouvoir. »

\all{Großbritannien ist eine Monarchie, aber der König regiert nicht wirklich.} « La Grande-Bretagne est une monarchie, mais le roi ne gouverne pas réellement. »

On remarque immédiatement que tous ces noms de systèmes sont \textbf{féminins} (\all{die}) sauf \all{der Staat} et \all{das System}. C'est une aide précieuse : \all{-ie}, \all{-ur} et \all{-tät} sont des terminaisons presque toujours féminines en allemand.

\begin{definition}[Les systèmes politiques — vocabulaire de base]
\begin{center}
\begin{tabularx}{\linewidth}{|l|l|l|X|}
\hline
\rowcolor{popblueL} \textbf{Mot} & \textbf{Article} & \textbf{Pluriel} & \textbf{Traduction} \\
\hline
Demokratie & die & die Demokratien & la démocratie \\
\hline
Diktatur & die & die Diktaturen & la dictature \\
\hline
Monarchie & die & die Monarchien & la monarchie \\
\hline
Republik & die & die Republiken & la république \\
\hline
Staat & der & die Staaten & l'État \\
\hline
System & das & die Systeme & le système \\
\hline
\end{tabularx}
\end{center}
\end{definition}

\begin{exemplebox}[Les systèmes en contexte]
\all{Die Schweiz ist eine Republik mit einer langen demokratischen Tradition.} « La Suisse est une république avec une longue tradition démocratique. »

\all{Der Staat schützt seine Bürger.} « L'État protège ses citoyens. »

\all{Jedes politische System hat Vorteile und Nachteile.} « Chaque système politique a des avantages et des inconvénients. »
\end{exemplebox}

\begin{attention}[Faux amis et pièges de genre]
\begin{itemize}
\item \all{der Staat} est \textbf{masculin} alors que « l'État » prend souvent un article qui fait hésiter les élèves ; dites toujours \all{der Staat}, jamais \all{das Staat}.
\item \all{die Republik} se dit avec \all{die} : on dit \all{die Bundesrepublik Deutschland} (la République fédérale d'Allemagne).
\item Attention au pluriel de \all{der Staat} : \all{die Staaten} (les États), avec un -\textbf{en} qui rappelle la déclinaison : \all{die Mitgliedstaaten der EU} « les États membres de l'UE ».
\end{itemize}
\end{attention}

\courssub{Champ 2 : les institutions et les personnes}

Qui gouverne ? Où se prennent les décisions ? Voici les acteurs et les lieux du pouvoir.

\begin{definition}[Institutions et personnes — vocabulaire]
\begin{center}
\begin{tabularx}{\linewidth}{|l|l|l|X|}
\hline
\rowcolor{popblueL} \textbf{Mot} & \textbf{Article} & \textbf{Pluriel} & \textbf{Traduction} \\
\hline
Präsident & der & die Präsidenten & le président \\
\hline
Bundeskanzler & der & die Bundeskanzler & le chancelier fédéral \\
\hline
Regierung & die & die Regierungen & le gouvernement \\
\hline
Parlament & das & die Parlamente & le parlement \\
\hline
Bundestag & der & --- (nom propre) & le Bundestag (parlement allemand) \\
\hline
Partei & die & die Parteien & le parti (politique) \\
\hline
Bürger / Bürgerin & der / die & die Bürger / die Bürgerinnen & le citoyen / la citoyenne \\
\hline
Politiker & der & die Politiker & l'homme politique, le politicien \\
\hline
\end{tabularx}
\end{center}
\end{definition}

\begin{propriete}[Le masculin et le féminin des personnes : le suffixe -in]
Pour les noms de personnes, l'allemand forme le féminin en ajoutant \cle{-in} au masculin, souvent avec tréma (Umlaut) au pluriel : \all{der Bürger} → \all{die Bürgerin} (pluriel \all{die Bürgerinnen}) ; \all{der Politiker} → \all{die Politikerin} ; \all{der Bundeskanzler} → \all{die Bundeskanzlerin}. C'est plus systématique qu'en français, où « le président / la présidente » varie selon les usages.
\end{propriete}

\begin{exemplebox}[Les institutions en contexte]
\all{Das Parlament diskutiert über ein neues Gesetz.} « Le parlement débat d'une nouvelle loi. »

\all{Die Regierung hat gestern eine wichtige Entscheidung getroffen.} « Le gouvernement a pris hier une décision importante. »

\all{Die Bürger haben das Recht zu wählen.} « Les citoyens ont le droit de voter. »

\all{Viele junge Politiker engagieren sich für die Umwelt.} « Beaucoup de jeunes hommes politiques s'engagent pour l'environnement. »
\end{exemplebox}

\begin{savaistu}[Qui fait quoi en Allemagne ?]
Le Parlement allemand s'appelle \all{der Bundestag} ; il siège à Berlin, dans le \all{Reichstagsgebäude} (le palais du Reichstag), surmonté d'une coupole de verre ouverte aux visiteurs — symbole de la transparence démocratique. La chancelière ou le chancelier (\all{der Bundeskanzler} / \all{die Bundeskanzlerin}) est le \textbf{chef du gouvernement} : c'est lui ou elle qui dirige réellement la politique du pays. \all{Der Bundespräsident} (le président fédéral), lui, a un rôle surtout \textbf{représentatif} : il représente l'État, signe les lois et prononce des discours, mais ne gouverne pas au quotidien. C'est une différence importante avec la France, où le président de la République exerce un pouvoir politique fort. L'Autriche fonctionne de manière comparable ; la Suisse, elle, est dirigée par un conseil de sept membres, le \all{Bundesrat}.
\end{savaistu}

\courssub{Champ 3 : les élections et la vie politique}

C'est le champ le plus riche et le plus fréquent dans les sujets du baccalauréat. Observons :

\all{Das Volk wählt den Präsidenten.} « Le peuple élit le président. »

\all{Die Partei hat die Wahl gewonnen.} « Le parti a gagné les élections. »

\all{Die Verfassung garantiert die Rechte der Bürger.} « La constitution garantit les droits des citoyens. »

\begin{definition}[Élections et vie politique — vocabulaire]
\begin{center}
\begin{tabularx}{\linewidth}{|l|l|l|X|}
\hline
\rowcolor{popblueL} \textbf{Mot} & \textbf{Article} & \textbf{Pluriel} & \textbf{Traduction} \\
\hline
Wahl & die & die Wahlen & l'élection, le choix \\
\hline
Wähler & der & die Wähler & l'électeur \\
\hline
Stimme & die & die Stimmen & la voix, le vote (suffrage) \\
\hline
Opposition & die & die Oppositionen & l'opposition \\
\hline
Macht & die & (die Mächte) & le pouvoir \\
\hline
Verfassung & die & die Verfassungen & la constitution \\
\hline
Gesetz & das & die Gesetze & la loi \\
\hline
Recht & das & die Rechte & le droit \\
\hline
Freiheit & die & die Freiheiten & la liberté \\
\hline
\end{tabularx}
\end{center}
Verbes associés : \all{wählen} (wählte, hat gewählt) « élire, choisir, voter » ; \all{abstimmen} (stimmte ab, hat abgestimmt) « voter, mettre aux voix » (verbe à particule séparable : \all{Der Bundestag stimmt morgen ab.}).
\end{definition}

\begin{attention}[der Wähler n'est pas die Wahl !]
\all{Der Wähler} = l'\textbf{électeur} (la personne qui vote) ; \all{die Wahl} = l'\textbf{élection} (l'événement). Ne confondez pas : \all{Die Wähler gehen zur Wahl.} « Les électeurs vont aux élections. » Autre piège : le verbe \all{wählen} a deux sens. Dans un contexte politique, il signifie « élire, choisir » ; au téléphone, il signifie « composer un numéro » : \all{Wählen Sie die 110!} « Composez le 110 ! » C'est le \textbf{contexte} qui décide du sens.
\end{attention}

\begin{attention}[das Recht ≠ das Gesetz]
Ces deux mots sont souvent confondus par les francophones car le français « droit » et « loi » se recoupent partiellement. \all{Das Gesetz} (pluriel \all{die Gesetze}) désigne une \textbf{loi concrète} votée par le parlement : \all{Das Parlament hat ein neues Gesetz beschlossen.} « Le parlement a adopté une nouvelle loi. » \all{Das Recht} (pluriel \all{die Rechte}) désigne le \textbf{droit} au sens de droit individuel ou de système juridique : \all{die Menschenrechte} « les droits de l'homme », \all{das Wahlrecht} « le droit de vote ». Retenez la formule : le parlement vote des \all{Gesetze}, la constitution garantit des \all{Rechte}.
\end{attention}

\courssub{Les familles de mots : un mot en cache trois}

L'allemand construit ses mots de façon très régulière. Connaître un verbe, c'est souvent connaître toute sa famille. Voici les quatre familles essentielles de notre thème :

\begin{center}
\begin{tabularx}{\linewidth}{|X|X|X|}
\hline
\rowcolor{popblueL} \textbf{Verbe / adjectif} & \textbf{Noms dérivés} & \textbf{Traductions} \\
\hline
wählen (élire) & die Wahl, der Wähler, die Wiederwahl & l'élection, l'électeur, la réélection \\
\hline
herrschen (régner, dominer) & der Herrscher, die Herrschaft & le souverain/dirigeant, la domination \\
\hline
regieren (gouverner) & die Regierung, der Regierungschef & le gouvernement, le chef du gouvernement \\
\hline
frei (libre) & die Freiheit, freiheitlich & la liberté, libéral (qui garantit les libertés) \\
\hline
\end{tabularx}
\end{center}

\begin{exemplebox}[Les familles en action]
\all{Nach vier Jahren will der Präsident die Wiederwahl.} « Après quatre ans, le président veut la réélection. »

\all{Der Herrscher übt seine Herrschaft über das Land aus.} « Le souverain exerce sa domination sur le pays. »

\all{Der Regierungschef leitet die Regierung.} « Le chef du gouvernement dirige le gouvernement. »

\all{Die Pressefreiheit ist ein Teil der freiheitlichen Demokratie.} « La liberté de la presse fait partie de la démocratie libérale. »
\end{exemplebox}

\begin{popfigure}
\begin{center}
\begin{tikzpicture}[mindmap, grow cyclic,
  every node/.style={concept, align=center, font=\small},
  root concept/.append style={concept color=popblue, font=\bfseries},
  level 1/.append style={level distance=3.4cm, sibling angle=72},
  level 1 concept/.append style={concept color=poporangeL}]
\node[root concept, align=center]{die Wahl\\ (l'élection)}
  child { node[align=center]{wählen\\ (élire, voter)} }
  child { node[align=center]{der Wähler\\ (l'électeur)} }
  child { node[align=center]{die Stimme\\ (le vote, la voix)} }
  child { node[align=center]{abstimmen\\ (voter, mettre aux voix)} }
  child { node[align=center]{eine Wahl gewinnen / verlieren\\ (gagner / perdre une élection)} };
\end{tikzpicture}
\end{center}
\legende{Carte mentale lexicale : la famille de \all{die Wahl} et ses expressions.}
\end{popfigure}

\courssub{Les expressions utiles : parler comme un journaliste}

Pour résumer un texte ou rédiger un paragraphe sur la politique, il faut connaître les \cle{collocations}, c'est-à-dire les combinaisons figées verbe + nom. Les voici :

\begin{center}
\begin{tabularx}{\linewidth}{|X|X|}
\hline
\rowcolor{popblueL} \textbf{Deutsch} & \textbf{Français} \\
\hline
eine Wahl gewinnen & gagner une élection \\
\hline
eine Wahl verlieren & perdre une élection \\
\hline
an die Macht kommen & arriver au pouvoir \\
\hline
die Macht ausüben & exercer le pouvoir \\
\hline
für ein Gesetz stimmen & voter pour une loi \\
\hline
gegen ein Gesetz stimmen & voter contre une loi \\
\hline
\end{tabularx}
\end{center}

\begin{exemplebox}[Les expressions en contexte]
\all{1933 ist Hitler an die Macht gekommen.} « En 1933, Hitler est arrivé au pouvoir. »

\all{In einer Demokratie übt das Volk die Macht aus.} « Dans une démocratie, c'est le peuple qui exerce le pouvoir. »

\all{Die Opposition hat gegen das Gesetz gestimmt.} « L'opposition a voté contre la loi. »

\all{Wer die Wahl verliert, muss das Ergebnis akzeptieren.} « Celui qui perd l'élection doit accepter le résultat. »
\end{exemplebox}

\begin{attention}[La préposition fait tout]
On dit \all{an die Macht kommen} avec l'\textbf{accusatif} (mouvement : on « arrive » au pouvoir) et \all{für/gegen ein Gesetz stimmen} avec \all{für} ou \all{gegen} + accusatif. Erreur fréquente des francophones : traduire « voter pour » par \all{stimmen für} sans article ou avec un \textbf{datif}. Phrase modèle à mémoriser : \all{Die Abgeordneten stimmen für das Gesetz.} « Les députés votent pour la loi. »
\end{attention}

\courssub{Texte d'entraînement : les élections dans un pays imaginaire}

Lisons maintenant un court texte qui réemploie tout le vocabulaire de cette section — exactement le type de document que vous rencontrerez à l'examen.

\begin{textebox}[Wahlen in Musterland]
Musterland ist eine Republik mit etwa zwanzig Millionen Bürgern. Alle vier Jahre finden Wahlen statt: Das Volk wählt das Parlament, und das Parlament wählt dann den Präsidenten. Die Verfassung des Landes garantiert die Rechte der Bürger, zum Beispiel die Meinungsfreiheit und die Pressefreiheit.

In diesem Jahr gibt es drei große Parteien. Die liberale Partei will die Wirtschaft stärken, die soziale Partei verspricht bessere Schulen und Krankenhäuser, und die grüne Partei kämpft für die Umwelt. Viele Wähler sind noch unsicher: „Ich weiß nicht, wen ich wählen soll“, sagt eine junge Bürgerin aus der Hauptstadt. „Alle Politiker versprechen viel vor der Wahl.“

Am Wahltag gehen Millionen Menschen in die Wahllokale. Jeder Wähler hat eine Stimme. Am Abend ist das Ergebnis klar: Die soziale Partei hat die Wahl gewonnen, die liberale Partei hat verloren und wird jetzt die Opposition bilden. Der neue Präsident erklärt: „Ich danke allen Wählern. Meine Regierung wird die Macht für alle Bürger ausüben, auch für die, die mich nicht gewählt haben.“

Das Parlament wird in den nächsten Wochen über neue Gesetze abstimmen – über Bildung, Gesundheit und Arbeit. Die Demokratie in Musterland funktioniert, sagen die Experten, aber die Bürger müssen wachsam bleiben: Freiheit und Rechte sind nie selbstverständlich.
\end{textebox}

\textbf{Glossaire.} \all{finden ... statt} : avoir lieu (verbe à particule séparable) ; \all{die Meinungsfreiheit} : la liberté d'opinion ; \all{die Pressefreiheit} : la liberté de la presse ; \all{die Wirtschaft} : l'économie ; \all{stärken} : renforcer ; \all{versprechen} (verspricht, versprach, hat versprochen) : promettre ; \all{kämpfen für} : lutter pour ; \all{unsicher} : indécis ; \all{das Wahllokal (-e)} : le bureau de vote ; \all{das Ergebnis (-se)} : le résultat ; \all{bilden} : former, constituer ; \all{wachsam} : vigilant ; \all{selbstverständlich} : évident, qui va de soi.

\textbf{Questions de compréhension.}
\begin{enumerate}
\item \all{Wer wählt den Präsidenten in Musterland?} (Qui élit le président à Musterland ?)
\item \all{Welche Rechte garantiert die Verfassung?} (Quels droits la constitution garantit-elle ?)
\item \all{Welche Partei hat die Wahl gewonnen, und was wird die liberale Partei jetzt?} (Quel parti a gagné l'élection, et que devient le parti libéral ?)
\item \all{Warum müssen die Bürger wachsam bleiben?} (Pourquoi les citoyens doivent-ils rester vigilants ?)
\end{enumerate}

\begin{corrige}[Questions de compréhension]
\begin{enumerate}
\item \all{Das Parlament wählt den Präsidenten.} (C'est le parlement qui élit le président — et non le peuple directement : c'est une république parlementaire.)
\item \all{Die Verfassung garantiert die Rechte der Bürger, zum Beispiel die Meinungsfreiheit und die Pressefreiheit.} (Elle garantit les droits des citoyens, par exemple la liberté d'opinion et la liberté de la presse.)
\item \all{Die soziale Partei hat die Wahl gewonnen. Die liberale Partei wird jetzt die Opposition bilden.} (Le parti social a gagné ; le parti libéral va former l'opposition.)
\item \all{Sie müssen wachsam bleiben, weil Freiheit und Rechte nie selbstverständlich sind.} (Parce que la liberté et les droits ne vont jamais de soi.)
\end{enumerate}
\end{corrige}

\courssub{Entraînement guidé}

\begin{exoresolu}[Complétez avec le bon mot (article et forme)]
Énoncé. Complétez les phrases avec le mot qui convient, tiré du vocabulaire de la leçon : \all{Wahl, Wähler, Gesetz, Recht, Regierung, Opposition, Macht}.
\begin{enumerate}
\item \all{Das Parlament stimmt morgen über ein neues \dots\ ab.}
\item \all{Jeder \dots\ hat nur eine Stimme.}
\item \all{Die \dots\ hat die Wahl verloren und kritisiert jetzt die Regierung.}
\item \all{Die Meinungsfreiheit ist ein wichtiges \dots\ der Bürger.}
\item \all{Nach der \dots\ bildet der Sieger eine neue \dots.}
\item \all{Der Diktator übt die \dots\ allein aus.}
\end{enumerate}
\tcblower
Solution détaillée.
\begin{enumerate}
\item \all{Gesetz} — le parlement vote des \emph{lois} (\all{über ein Gesetz abstimmen}) ; \all{Recht} serait un faux sens ici.
\item \all{Wähler} — c'est l'\emph{électeur} (la personne) qui possède une voix ; \all{die Wahl} désignerait l'événement.
\item \all{Opposition} — le parti qui perd et critique le gouvernement joue le rôle d'opposition.
\item \all{Recht} — la liberté d'opinion est un \emph{droit} individuel garanti par la constitution, pas une loi votée.
\item \all{Wahl} ; \all{Regierung} — après l'\emph{élection}, le vainqueur forme un nouveau \emph{gouvernement}.
\item \all{Macht} — expression figée \all{die Macht ausüben} « exercer le pouvoir ».
\end{enumerate}
\end{exoresolu}

\begin{exercice}[À vous de jouer]
\begin{enumerate}
\item Traduisez en allemand : a) « Le peuple élit le président tous les cinq ans. » b) « Le parti a perdu les élections. » c) « Les citoyens votent pour une nouvelle loi. » d) « La constitution garantit la liberté de la presse. »
\item Donnez l'article et le pluriel de : \all{Partei, Staat, Verfassung, Gesetz, Bürgerin, System}.
\item Formez le nom dérivé : \all{wählen} (l'événement), \all{regieren} (l'institution), \all{herrschen} (la personne), \all{frei} (le concept).
\item Rédigez trois phrases sur les élections au Cameroun en employant au moins trois expressions de la leçon (\all{eine Wahl gewinnen}, \all{an die Macht kommen}, \all{für/gegen ein Gesetz stimmen}, \all{die Macht ausüben}).
\end{enumerate}
\end{exercice}

\begin{aretenir}[L'essentiel du lexique politique]
\begin{itemize}
\item Les noms de systèmes en \all{-ie}, \all{-ur} sont féminins : \all{die Demokratie, die Diktatur, die Monarchie, die Republik} ; mais \all{der Staat} et \all{das System}.
\item \all{der Wähler} = la personne qui vote ; \all{die Wahl} = l'événement électoral ; \all{die Stimme} = le suffrage (et la voix).
\item \all{das Gesetz} (die Gesetze) = la loi votée ; \all{das Recht} (die Rechte) = le droit garanti.
\item Expressions à connaître par cœur : \all{eine Wahl gewinnen/verlieren}, \all{an die Macht kommen}, \all{die Macht ausüben}, \all{für/gegen ein Gesetz stimmen}.
\item Familles de mots : \all{wählen → die Wahl, der Wähler, die Wiederwahl} ; \all{regieren → die Regierung} ; \all{herrschen → der Herrscher, die Herrschaft} ; \all{frei → die Freiheit}.
\end{itemize}
\end{aretenir}

Ce vocabulaire est désormais votre boîte à outils. Dans la section suivante, nous verrons comment ces mots se combinent dans le \all{Nominalstil}, ce style nominal si caractéristique des textes politiques et journalistiques allemands.

\coursec{Grammaire 1 : le style nominal (der Nominalstil)}

Dans la section précédente, nous avons découvert le vocabulaire des systèmes politiques : \all{die Demokratie}, \all{die Wahl}, \all{die Partei}, \all{die Verfassung}. Vous avez peut-être remarqué que beaucoup de ces mots sont des \cle{noms} qui désignent pourtant des \emph{actions} : \all{die Wahl} (le fait de choisir, d'élire), \all{die Regierung} (le fait de gouverner). Ce n'est pas un hasard : la langue allemande écrite aime exprimer les actions par des noms plutôt que par des verbes. C'est ce qu'on appelle le \cle{style nominal} (\all{der Nominalstil}), une caractéristique essentielle des textes de presse et des sujets du Baccalauréat.

\courssub{Observation : deux façons de dire la même chose}

Reprenons une expression tirée de notre texte support « Politische Systeme in der Welt » :

\begin{center}
\all{die Wahl der Regierung durch das Volk}\\
« l'élection du gouvernement par le peuple »
\end{center}

Comparons-la avec une phrase verbale complète :

\begin{center}
\all{Das Volk wählt die Regierung.}\\
« Le peuple élit le gouvernement. »
\end{center}

\begin{exemplebox}[Observation guidée]
Regardez bien ce qui a changé entre les deux formulations :
\begin{itemize}
\item le \textbf{verbe} \all{wählt} est devenu un \textbf{nom} : \all{die Wahl} ;
\item le \textbf{sujet} \all{das Volk} est devenu un complément introduit par \all{durch} (+ datif) : \all{durch das Volk} ;
\item le \textbf{complément d'objet} \all{die Regierung} est devenu un \textbf{génitif} : \all{der Regierung} ;
\item le tout ne forme plus une phrase, mais un \textbf{groupe nominal} qui peut s'insérer dans une phrase plus grande : \all{Nach der Wahl der Regierung durch das Volk beginnt die Arbeit.} « Après l'élection du gouvernement par le peuple, le travail commence. »
\end{itemize}
\end{exemplebox}

\begin{popfigure}
\begin{center}
\begin{tikzpicture}[node distance=1.2cm]
\node[draw=popblue, fill=popblueL, rounded corners, align=center, text width=5.6cm, inner sep=6pt] (verbal) {\textbf{Verbalsatz}\\ \all{Das Volk wählt die Regierung.}\\ \footnotesize sujet + verbe + COD};
\node[draw=poporange, fill=poporangeL, rounded corners, align=center, text width=5.6cm, inner sep=6pt, right=2.2cm of verbal] (nominal) {\textbf{Nominalstil}\\ \all{die Wahl der Regierung durch das Volk}\\ \footnotesize nom d'action + génitif + \all{durch}};
\draw[-{Stealth[length=3mm]}, very thick, popgreen] (verbal.north) to[bend left=25] node[above, align=center, font=\footnotesize]{verbe $\rightarrow$ nom} (nominal.north);
\draw[-{Stealth[length=3mm]}, very thick, poppurple] (nominal.south) to[bend left=25] node[below, align=center, font=\footnotesize]{nom $\rightarrow$ verbe\\ (reformulation)} (verbal.south);
\end{tikzpicture}
\legende{La double transformation : phrase verbale $\rightleftarrows$ style nominal}
\end{center}
\end{popfigure}

\courssub{Définition et emplois}

\begin{definition}[Der Nominalstil]
Le \cle{style nominal} (\all{der Nominalstil}) est un style où l'\textbf{action est exprimée par un nom} (appelé \emph{nom d'action}) au lieu d'un verbe : \all{die Wahl} au lieu de \all{wählen}, \all{die Diskussion} au lieu de \all{diskutieren}, \all{die Entscheidung} au lieu de \all{entscheiden}. Ce style est très fréquent dans la langue \textbf{écrite} : presse, textes administratifs, lois, discours politiques — et donc dans les textes proposés au Baccalauréat.
\end{definition}

Pourquoi l'allemand aime-t-il tant ce style ? Parce qu'un nom peut être entouré d'articles, d'adjectifs, de prépositions et de génitifs : il permet de condenser beaucoup d'informations en peu de mots. Comparez :

\begin{center}
\begin{tabularx}{\linewidth}{|X|X|}
\hline
\rowcolor{popblueL} \textbf{Style verbal (oral, simple)} & \textbf{Style nominal (écrit, condensé)} \\
\hline
\all{Die Regierung wurde gewählt.} & \all{nach der Wahl der Regierung} \\
\hline
\all{Die Parteien diskutieren das Gesetz.} & \all{während der Diskussion des Gesetzes} \\
\hline
\all{Die Bürger nehmen an den Wahlen teil.} & \all{dank der Teilnahme der Bürger an den Wahlen} \\
\hline
\all{Man hat die Verfassung geändert.} & \all{trotz der Änderung der Verfassung} \\
\hline
\end{tabularx}
\end{center}

\begin{attention}[Comparaison avec le français]
En français, on préfère souvent le \textbf{verbe} : « parce que les citoyens participent aux élections… ». En allemand écrit, le style nominal est \textbf{valorisé} et considéré comme plus élégant : \all{wegen der Teilnahme der Bürger an den Wahlen…}. Au Bac, vous devez donc :
\begin{enumerate}
\item savoir \textbf{reconnaître} un nom d'action pour comprendre le texte ;
\item savoir le \textbf{traduire} souvent par un verbe français ;
\item savoir le \textbf{transformer} en phrase verbale pour reformuler (résumé, questions de compréhension).
\end{enumerate}
Traduire mot à mot « la participation des citoyens aux élections » est correct, mais « parce que les citoyens participent aux élections » est souvent plus naturel en français.
\end{attention}

\courssub{Formation : du verbe au nom d'action}

La plupart des noms d'action sont \textbf{féminins} (\all{die}) et se forment à partir du radical du verbe. Voici les grandes familles :

\begin{center}
\begin{tabularx}{\linewidth}{|l|l|X|}
\hline
\rowcolor{popblueL} \textbf{Verbe} & \textbf{Nom d'action} & \textbf{Formation / Remarque} \\
\hline
\all{wählen} (élire) & \all{die Wahl, -en} & radical sans suffixe \\
\hline
\all{regieren} (gouverner) & \all{die Regierung, -en} & suffixe \all{-ung} \\
\hline
\all{entscheiden} (décider) & \all{die Entscheidung, -en} & suffixe \all{-ung} \\
\hline
\all{diskutieren} (discuter) & \all{die Diskussion, -en} & suffixe \all{-ion} (verbes en \all{-ieren}) \\
\hline
\all{teilnehmen} (participer) & \all{die Teilnahme, -n} & suffixe \all{-nahme} \\
\hline
\all{sich beteiligen} (participer) & \all{die Beteiligung, -en} & suffixe \all{-ung} \\
\hline
\all{gründen} (fonder) & \all{die Gründung, -en} & suffixe \all{-ung} \\
\hline
\all{kämpfen} (lutter) & \all{der Kampf, die Kämpfe} & radical, masculin (exception) \\
\hline
\all{ändern} (modifier) & \all{die Änderung, -en} & suffixe \all{-ung} \\
\hline
\all{meinen} (penser) & \all{die Meinung, -en} & suffixe \all{-ung} \\
\hline
\end{tabularx}
\end{center}

\begin{propriete}[Règles de formation]
\begin{enumerate}
\item Le suffixe le plus productif est \all{-ung} : il donne toujours un nom \textbf{féminin} au pluriel en \all{-en} : \all{regieren} $\rightarrow$ \all{die Regierung}, \all{die Regierungen}.
\item Les verbes en \all{-ieren} (souvent d'origine française ou latine) donnent des noms en \all{-ion} : \all{diskutieren} $\rightarrow$ \all{die Diskussion}, \all{reagieren} $\rightarrow$ \all{die Reaktion}, \all{informieren} $\rightarrow$ \all{die Information}.
\item Certains noms reprennent le radical nu du verbe, parfois avec modification de la voyelle : \all{wählen} $\rightarrow$ \all{die Wahl}, \all{kämpfen} $\rightarrow$ \all{der Kampf}, \all{beginnen} $\rightarrow$ \all{der Beginn}.
\item Les verbes à particule séparable gardent leur particule : \all{teilnehmen} $\rightarrow$ \all{die Teil}\textbf{nahme}, \all{zunehmen} $\rightarrow$ \all{die Zu}\textbf{nahme} (l'augmentation), \all{abnehmen} $\rightarrow$ \all{die Ab}\textbf{nahme} (la diminution).
\end{enumerate}
\end{propriete}

\courssub{La transformation : phrase verbale $\rightarrow$ expression nominale}

\begin{propriete}[Règle de transformation en trois étapes]
Pour transformer une phrase verbale en style nominal :
\begin{enumerate}
\item Le \textbf{verbe} devient un \textbf{nom d'action} avec article : \all{wählen} $\rightarrow$ \all{die Wahl}.
\item Le \textbf{sujet} devient un \textbf{génitif} ou une préposition \all{durch} / \all{von} + datif : \all{das Volk} $\rightarrow$ \all{des Volkes} ou \all{durch das Volk}.
\item Le \textbf{complément} devient un \textbf{génitif} ou la préposition associée au nom : \all{die Regierung} $\rightarrow$ \all{der Regierung}.
\end{enumerate}
\end{propriete}

\begin{exemplebox}[Transformations commentées]
\textbf{a)} \all{Das Parlament diskutiert das Gesetz.} « Le parlement discute la loi. »\\
$\rightarrow$ \all{die Diskussion des Gesetzes im Parlament} « la discussion de la loi au parlement »\\
(verbe $\rightarrow$ \all{die Diskussion} ; COD $\rightarrow$ génitif \all{des Gesetzes} ; sujet $\rightarrow$ complément de lieu \all{im Parlament}.)

\textbf{b)} \all{Die Bürger beteiligen sich an der Politik.} « Les citoyens participent à la politique. »\\
$\rightarrow$ \all{die Beteiligung der Bürger an der Politik} « la participation des citoyens à la politique »\\
(sujet $\rightarrow$ génitif \all{der Bürger} ; la préposition \all{an} + datif est conservée.)

\textbf{c)} \all{Das Volk wählt den Präsidenten.} « Le peuple élit le président. »\\
$\rightarrow$ \all{die Wahl des Präsidenten durch das Volk} « l'élection du président par le peuple »\\
(sujet $\rightarrow$ \all{durch} + datif ; COD $\rightarrow$ génitif.)
\end{exemplebox}

\courssub{Les prépositions associées aux noms d'action}

Chaque nom d'action garde en général la préposition de son verbe d'origine. Il faut les apprendre \textbf{ensemble} :

\begin{center}
\begin{tabularx}{\linewidth}{|l|l|X|}
\hline
\rowcolor{popblueL} \textbf{Verbe + préposition} & \textbf{Nom + préposition} & \textbf{Exemple traduit} \\
\hline
\all{teilnehmen an (+Dat)} & \all{die Teilnahme an (+Dat)} & \all{die Teilnahme an den Wahlen} : la participation aux élections \\
\hline
\all{diskutieren über (+Akk)} & \all{die Diskussion über (+Akk)} & \all{die Diskussion über das Gesetz} : la discussion sur la loi \\
\hline
\all{wählen zu (+Dat)} & \all{die Wahl zu (+Dat)} & \all{die Wahl zum Präsidenten} : l'élection au poste de président \\
\hline
\all{kämpfen gegen/für (+Akk)} & \all{der Kampf gegen/für (+Akk)} & \all{der Kampf für die Demokratie} : le combat pour la démocratie \\
\hline
\all{sich beteiligen an (+Dat)} & \all{die Beteiligung an (+Dat)} & \all{die Beteiligung an der Politik} : la participation à la politique \\
\hline
\all{hoffen auf (+Akk)} & \all{die Hoffnung auf (+Akk)} & \all{die Hoffnung auf Frieden} : l'espoir de paix \\
\hline
\all{denken an (+Akk)} & \all{der Gedanke an (+Akk)} & \all{der Gedanke an die Freiheit} : la pensée de la liberté \\
\hline
\end{tabularx}
\end{center}

\begin{attention}[Pièges fréquents des francophones]
\begin{itemize}
\item \textbf{Ne pas oublier le cas} après la préposition : \all{an den Wahlen} (datif pluriel), jamais \all{an die Wahlen} dans ce sens.
\item \textbf{Le génitif allemand} remplace le « de » français : \all{die Wahl des Präsidenten} = « l'élection \emph{du} président ». Ne dites pas \all{die Wahl von dem Präsidenten} dans un texte soigné (le \all{von} + datif est toléré à l'oral, mais le génitif est la norme écrite).
\item \textbf{Majuscule obligatoire} au nom d'action : \all{die Wahl}, jamais \all{die wahl}.
\item Attention au verbe réfléchi : \all{sich beteiligen} perd son \all{sich} en passant au nom : \all{die Beteiligung} (pas de « die Sichbeteiligung » !).
\end{itemize}
\end{attention}

\courssub{La transformation inverse : du nominal au verbal}

C'est la compétence la plus utile au Bac : face à un groupe nominal dense, le « déplier » en phrase verbale pour vérifier qu'on a compris.

\begin{methode}[Déplier un groupe nominal en 4 étapes]
Prenons : \all{die Diskussion der Abgeordneten über die Verfassung}.
\begin{enumerate}
\item \textbf{Identifier le nom d'action} et retrouver le verbe de base : \all{die Diskussion} $\rightarrow$ \all{diskutieren} (discuter).
\item \textbf{Identifier le sujet} (génitif ou \all{durch}/\all{von} + datif) : \all{der Abgeordneten} = les députés (c'est eux qui discutent).
\item \textbf{Identifier le complément} (préposition) : \all{über die Verfassung} = sur la constitution.
\item \textbf{Reconstruire la phrase verbale} : \all{Die Abgeordneten diskutieren über die Verfassung.} « Les députés discutent de la constitution. »
\end{enumerate}
\end{methode}

\begin{methode}[Deviner le sens d'un nom d'action inconnu]
Face à un nom d'action difficile dans un texte, \textbf{retrouvez le verbe de base} : enlevez le suffixe (\all{-ung}, \all{-ion}, \all{-nahme}) et cherchez le radical.
\begin{itemize}
\item \all{die Entscheidung} $\rightarrow$ \all{entscheiden} = décider $\rightarrow$ « la décision »
\item \all{die Gründung} $\rightarrow$ \all{gründen} = fonder $\rightarrow$ « la fondation »
\item \all{die Änderung} $\rightarrow$ \all{ändern} = modifier $\rightarrow$ « la modification »
\item \all{die Zustimmung} $\rightarrow$ \all{zustimmen} = approuver $\rightarrow$ « l'approbation, le consentement »
\end{itemize}
Cette stratégie vous évite de bloquer sur un mot inconnu pendant l'épreuve de compréhension de l'écrit.
\end{methode}

\courssub{Texte d'application : le style nominal dans un article de presse}

\begin{textebox}[Aus einer Zeitung : „Jugend und Politik“]
\textbf{Die Beteiligung der Jugend an der Politik}

Die Teilnahme junger Menschen an den Wahlen ist in vielen Ländern ein Problem. Nach der Wahl des neuen Parlaments in einem europäischen Land zeigte eine Untersuchung: Nur die Hälfte der Wähler unter 25 Jahren hat an der Abstimmung teilgenommen.

Politiker sehen den Grund in der Enttäuschung vieler Jugendlicher über die Arbeit der Parteien. „Die Diskussion über die Zukunft findet oft ohne die Jugend statt“, sagt eine Expertin. Sie fordert die Gründung von Jugendparlamenten in allen Städten und die Einführung von Politikunterricht in allen Schulen.

Ein Beispiel aus Afrika zeigt eine andere Entwicklung: In einer Stadt in Kamerun führte die Gründung eines Jugendrates zu einer stärkeren Beteiligung der jungen Bürger an den lokalen Entscheidungen. Durch die Zusammenarbeit mit den älteren Ratsmitgliedern lernen die Jugendlichen die Bedeutung der Demokratie kennen.

Die Hoffnung auf Veränderung bleibt: Der Kampf für mehr Demokratie beginnt oft mit der Meinung eines einzelnen Menschen. Die Entscheidung, zur Wahl zu gehen, ist der erste Schritt. Denn die Wahl der Regierung durch das Volk ist das Herz jeder Demokratie — ohne die Teilnahme der Bürger verliert sie ihre Kraft.

\textbf{Glossaire :} die Beteiligung, -en (la participation) — die Untersuchung, -en (l'étude, l'enquête) — die Abstimmung, -en (le scrutin, le vote) — die Enttäuschung, -en (la déception) — fordern (exiger, réclamer) — die Einführung, -en (l'introduction, l'instauration) — der Jugendrat, -räte (le conseil des jeunes) — die Zusammenarbeit (la collaboration) — das Ratsmitglied, -er (le membre du conseil) — die Bedeutung, -en (l'importance, la signification) — die Veränderung, -en (le changement) — der Schritt, -e (le pas).
\end{textebox}

\textbf{Questions de compréhension (oral) :}
\begin{enumerate}
\item \all{Wer hat an der Abstimmung teilgenommen?} (Qui a participé au scrutin ?)
\item \all{Was fordert die Expertin?} (Qu'est-ce que l'experte réclame ?)
\item \all{Wozu führte die Gründung eines Jugendrates in Kamerun?} (À quoi a mené la fondation d'un conseil des jeunes au Cameroun ?)
\item Soulignez dans le texte cinq noms d'action et retrouvez leur verbe de base.
\end{enumerate}

\courssub{Exercice résolu et entraînement}

\begin{exoresolu}[Transformation : verbal $\rightarrow$ nominal]
\textbf{Énoncé.} Transformez les phrases suivantes en style nominal (groupe nominal), comme dans le modèle : \all{Die Bürger nehmen an den Wahlen teil.} $\rightarrow$ \all{die Teilnahme der Bürger an den Wahlen}.
\begin{enumerate}
\item \all{Die Parteien diskutieren über das neue Gesetz.}
\item \all{Das Volk wählt den Präsidenten.}
\item \all{Die Jugendlichen beteiligen sich an der Politik.}
\item \all{Die Opposition kämpft gegen die Diktatur.}
\end{enumerate}
\tcblower
\textbf{Solution détaillée.}
\begin{enumerate}
\item Verbe \all{diskutieren} $\rightarrow$ nom \all{die Diskussion} ; sujet \all{die Parteien} $\rightarrow$ génitif \all{der Parteien} ; préposition \all{über} + accusatif conservée : \textbf{\all{die Diskussion der Parteien über das neue Gesetz}} (la discussion des partis sur la nouvelle loi).
\item Verbe \all{wählen} $\rightarrow$ \all{die Wahl} ; COD \all{den Präsidenten} $\rightarrow$ génitif \all{des Präsidenten} ; sujet \all{das Volk} $\rightarrow$ \all{durch} + datif : \textbf{\all{die Wahl des Präsidenten durch das Volk}} (l'élection du président par le peuple).
\item Verbe réfléchi \all{sich beteiligen} $\rightarrow$ \all{die Beteiligung} (le \all{sich} disparaît) ; sujet $\rightarrow$ génitif ; préposition \all{an} + datif conservée : \textbf{\all{die Beteiligung der Jugendlichen an der Politik}} (la participation des jeunes à la politique).
\item Verbe \all{kämpfen} $\rightarrow$ \all{der Kampf} (masculin, exception !) ; préposition \all{gegen} + accusatif conservée ; sujet $\rightarrow$ génitif : \textbf{\all{der Kampf der Opposition gegen die Diktatur}} (le combat de l'opposition contre la dictature).
\end{enumerate}
\end{exoresolu}

\begin{exercice}[À vous : nominal $\rightarrow$ verbal]
Transformez les groupes nominaux suivants (tirés du texte « Jugend und Politik ») en phrases verbales complètes, puis traduisez-les en français :
\begin{enumerate}
\item \all{die Gründung von Jugendparlamenten in allen Städten}
\item \all{die Einführung von Politikunterricht in allen Schulen}
\item \all{die Zusammenarbeit der Jugendlichen mit den älteren Ratsmitgliedern}
\item \all{die Entscheidung der Bürger, zur Wahl zu gehen}
\end{enumerate}
\end{exercice}

\begin{corrige}[Exercice : nominal $\rightarrow$ verbal]
\begin{enumerate}
\item \all{Man gründet Jugendparlamente in allen Städten.} « On fonde des parlements de jeunes dans toutes les villes. » (verbe : \all{gründen})
\item \all{Man führt Politikunterricht in allen Schulen ein.} « On introduit l'éducation politique dans toutes les écoles. » (verbe à particule : \all{einführen} — attention, la particule \all{ein-} retourne à la fin !)
\item \all{Die Jugendlichen arbeiten mit den älteren Ratsmitgliedern zusammen.} « Les jeunes collaborent avec les membres plus âgés du conseil. » (verbe : \all{zusammenarbeiten})
\item \all{Die Bürger entscheiden sich, zur Wahl zu gehen.} « Les citoyens décident d'aller voter. » (verbe réfléchi : \all{sich entscheiden})
\end{enumerate}
Remarque : en passant du nom au verbe, on retrouve la particule séparable (\all{Einführung} $\rightarrow$ \all{führt … ein}) et le pronom réfléchi (\all{Entscheidung} $\rightarrow$ \all{entscheiden sich}).
\end{corrige}

\begin{aretenir}[L'essentiel du style nominal]
\begin{itemize}
\item Le style nominal exprime l'action par un \textbf{nom} : \all{die Wahl}, \all{die Diskussion}, \all{die Entscheidung}.
\item Formation : verbe + suffixe (\all{-ung}, \all{-ion}, \all{-nahme}) ou radical nu ; ces noms sont presque toujours \textbf{féminins} (exception : \all{der Kampf}, \all{der Beginn}).
\item Transformation : verbe $\rightarrow$ nom ; sujet $\rightarrow$ génitif ou \all{durch}/\all{von} + datif ; complément $\rightarrow$ génitif ou préposition associée.
\item Les prépositions s'apprennent avec le nom : \all{die Teilnahme an}, \all{die Diskussion über}, \all{die Wahl zu}, \all{der Kampf für/gegen}.
\item Au Bac : reconnaître le nom d'action, retrouver son verbe de base, déplier le groupe nominal en phrase verbale pour comprendre et reformuler.
\end{itemize}
\end{aretenir}

\begin{savaistu}[Pourquoi l'allemand aime-t-il tant les noms ?]
Le style nominal est si typique de l'allemand administratif qu'il existe un mot pour le critiquer : \all{die Behördensprache} (la langue des administrations), parfois moquée comme \all{„Amtsdeutsch“}. Les journalistes allemands eux-mêmes conseillent aux élèves de « réveiller les verbes » pour écrire de façon vivante. Mais pour les examens, ce style reste la norme des textes proposés : le maîtriser, c'est tenir la clé de la compréhension de l'écrit. Au Cameroun, où l'allemand est langue étrangère, c'est souvent ce style qui déroute les candidats — plus maintenant !
\end{savaistu}

\coursec{Grammaire 2 : connecteurs logiques, passif et futur}

Dans la section précédente, nous avons vu comment le \cle{style nominal} condense l'information en transformant des verbes en noms (\all{die Wahl} au lieu de \all{wählen}). Mais un texte politique ne se résume pas à des noms : pour construire un raisonnement, exprimer une cause, une opposition ou une conséquence, il faut relier les phrases entre elles. C'est le rôle des \cle{connecteurs logiques}. Nous verrons ensuite deux formes verbales essentielles du discours politique : le \cle{passif} (\all{Der Bundeskanzler wird gewählt.}) et le \cle{futur} (\all{Die Bürger werden wählen.}).

\courssub{Les connecteurs logiques : trois groupes, trois places du verbe}

\textbf{Observation.} Lisons d'abord trois phrases tirées de notre thème :

\all{Die Diktatur unterdrückt die Opposition, denn sie fürchtet Kritik.} « La dictature réprime l'opposition, car elle craint la critique. »

\all{Es gibt keine freien Wahlen, deshalb protestieren die Bürger.} « Il n'y a pas d'élections libres, c'est pourquoi les citoyens manifestent. »

\all{Obwohl das Volk wählt, hat eine Partei die ganze Macht.} « Bien que le peuple vote, un seul parti détient tout le pouvoir. »

Observez la place du verbe conjugué (soulignons-le mentalement) : après \all{denn}, le verbe reste à sa place habituelle ; après \all{deshalb}, le verbe passe \emph{avant} le sujet ; après \all{obwohl}, le verbe est renvoyé \emph{à la fin}. Toute la difficulté — et toute la logique — des connecteurs allemands tient à cette question : \textbf{où va le verbe ?}

\begin{propriete}[Les trois groupes de connecteurs]
En allemand, les connecteurs se classent en trois groupes selon la place qu'ils imposent au verbe conjugué :
\begin{enumerate}
\item \textbf{Coordination sans inversion (position 0)} : \all{und} (et), \all{aber} (mais), \all{denn} (car), \all{oder} (ou), \all{sondern} (mais, après une négation). Le connecteur occupe la « position 0 » : il ne compte pas dans la structure de la phrase. Le verbe reste en \textbf{position 2}, comme si de rien n'était.
\item \textbf{Connecteurs avec inversion (position 1)} : \all{deshalb} (c'est pourquoi), \all{trotzdem} (pourtant, quand même), \all{dann} (ensuite), \all{zuerst} (d'abord). Le connecteur occupe la \textbf{position 1} ; le verbe suit immédiatement en position 2, donc le sujet passe après le verbe : c'est l'\textbf{inversion}.
\item \textbf{Subordination (verbe à la fin)} : \all{weil} (parce que), \all{obwohl} (bien que), \all{dass} (que), \all{wenn} (quand, si), \all{während} (pendant que, tandis que). Le verbe conjugué est renvoyé \textbf{en dernière position} de la subordonnée.
\end{enumerate}
\end{propriete}

\begin{exemplebox}[Exemples traduits — un par groupe]
\all{Die Demokratie schützt die Freiheit, aber sie ist nicht perfekt.} « La démocratie protège la liberté, mais elle n'est pas parfaite. »

\all{Die Wahlen waren nicht frei, deshalb gab es Proteste.} « Les élections n'étaient pas libres, c'est pourquoi il y a eu des manifestations. »

\all{Weil das Volk den Bundestag wählt, spricht man von einer Demokratie.} « Parce que le peuple élit le Bundestag, on parle de démocratie. »
\end{exemplebox}

\textbf{Remarques importantes.}
\begin{itemize}
\item \all{sondern} ne s'emploie qu'après une \emph{négation} pour introduire une correction : \all{Deutschland ist keine Monarchie, sondern eine Republik.} « L'Allemagne n'est pas une monarchie, mais une république. » Sinon, on utilise \all{aber}.
\item Quand la subordonnée (\all{weil}, \all{obwohl}...) est placée \emph{en tête}, la principale suit avec inversion : \all{Obwohl das Volk wählt, \textbf{hat} eine Partei die ganze Macht.} La virgule joue le rôle de position 1.
\item \all{während} peut signifier « pendant que » (temps) ou « tandis que » (opposition) : \all{Während die einen demonstrieren, schweigen die anderen.} « Tandis que les uns manifestent, les autres se taisent. »
\end{itemize}

\begin{center}
\begin{tabularx}{\linewidth}{|l|X|X|X|}
\hline
\rowcolor{popblueL} Groupe & Connecteurs & Place du verbe & Exemple \\
\hline
Coordination (position 0) & und, aber, denn, oder, sondern & position 2, ordre normal & \all{..., denn sie fürchtet Kritik.} \\
\hline
Inversion (position 1) & deshalb, trotzdem, dann, zuerst & verbe puis sujet & \all{..., deshalb protestieren die Bürger.} \\
\hline
Subordination & weil, obwohl, dass, wenn, während & verbe conjugué à la fin & \all{Obwohl das Volk wählt, ...} \\
\hline
\end{tabularx}
\end{center}

\begin{popfigure}
\begin{tikzpicture}[node distance=0.4cm]
\node[draw=popblue, fill=popblueL, rounded corners, align=center, text width=4.6cm, minimum height=3.4cm] (a) {\textbf{Position 0}\\[2pt] \all{und, aber, denn,}\\ \all{oder, sondern}\\[4pt] \footnotesize\all{Sie fürchtet Kritik,}\\ \footnotesize\all{\textbf{denn} sie hat Angst.}\\ \footnotesize verbe en position 2};
\node[draw=poporange, fill=poporangeL, rounded corners, align=center, text width=4.6cm, minimum height=3.4cm, right=of a] (b) {\textbf{Position 1 + inversion}\\[2pt] \all{deshalb, trotzdem,}\\ \all{dann, zuerst}\\[4pt] \footnotesize\all{..., \textbf{deshalb} \emph{protestieren}}\\ \footnotesize\all{\emph{die Bürger}.}\\ \footnotesize verbe avant le sujet};
\node[draw=poppurple, fill=poppurpleL, rounded corners, align=center, text width=4.6cm, minimum height=3.4cm, right=of b] (c) {\textbf{Subordination}\\[2pt] \all{weil, obwohl, dass,}\\ \all{wenn, während}\\[4pt] \footnotesize\all{\textbf{Weil} das Volk den}\\ \footnotesize\all{Bundestag \emph{wählt}, ...}\\ \footnotesize verbe à la fin};
\end{tikzpicture}
\legende{Les trois groupes de connecteurs : la place du verbe change selon le connecteur.}
\end{popfigure}

\begin{attention}[Le piège n° 1 des francophones : deshalb]
En français, « c'est pourquoi les citoyens manifestent » garde l'ordre sujet-verbe. En allemand, après \all{deshalb} (et \all{trotzdem}, \all{dann}...), l'inversion est \textbf{obligatoire} :

\all{deshalb die Bürger protestieren} ✗

\all{deshalb \textbf{protestieren} die Bürger} ✓

Astuce : le verbe conjugué est \emph{toujours} en position 2 dans une phrase principale allemande. Si un connecteur occupe la position 1, le verbe doit le suivre immédiatement.
\end{attention}

\begin{exoresolu}[Reliez les phrases avec le connecteur indiqué]
Énoncé. Reliez : 1) \all{Das Volk demonstriert. Die Regierung reagiert nicht.} (\all{aber}) — 2) \all{Die Wahl war manipuliert. Die Opposition klagt.} (\all{deshalb}) — 3) \all{Die Verfassung garantiert die Freiheit. Die Bürger können frei wählen.} (\all{weil}, subordonnée en tête).
\tcblower
Solution détaillée.
\begin{enumerate}
\item \all{aber} = position 0, ordre normal : \all{Das Volk demonstriert, \textbf{aber} die Regierung \emph{reagiert} nicht.} « Le peuple manifeste, mais le gouvernement ne réagit pas. »
\item \all{deshalb} = position 1, inversion : \all{Die Wahl war manipuliert, \textbf{deshalb} \emph{klagt} die Opposition.} « L'élection était truquée, c'est pourquoi l'opposition porte plainte. »
\item \all{weil} = subordination, verbe à la fin ; subordonnée en tête, donc inversion dans la principale : \all{\textbf{Weil} die Verfassung die Freiheit \emph{garantiert}, \emph{können} die Bürger frei wählen.} « Parce que la constitution garantit la liberté, les citoyens peuvent voter librement. »
\end{enumerate}
\end{exoresolu}

\courssub{Le passif : werden + Partizip II}

\textbf{Observation.} Dans un texte politique, on s'intéresse souvent plus à l'\emph{action} qu'à celui qui la fait. Comparez :

\all{Das Volk wählt den Bundespräsidenten.} (actif) « Le peuple élit le président fédéral. »

\all{Der Bundespräsident wird gewählt.} (passif) « Le président fédéral est élu. »

Au passif, le complément d'objet direct de l'actif (\all{den Bundespräsidenten}) devient le \textbf{sujet} ; l'ancien sujet disparaît ou devient le \textbf{complément d'agent}.

\begin{propriete}[Formation du passif (Vorgangspassiv)]
Le passif allemand se forme avec l'auxiliaire \all{werden} conjugué + le \textbf{Partizip II} du verbe principal :
\begin{itemize}
\item \textbf{Präsens} : \all{Der Bundeskanzler \textbf{wird} gewählt.} « Le chancelier est élu. »
\item \textbf{Präteritum} : \all{Der Bundeskanzler \textbf{wurde} gewählt.} « Le chancelier fut / a été élu. »
\end{itemize}
Conjugaison de \all{werden} au Präsens : \all{ich werde, du wirst, er/sie/es wird, wir werden, ihr werdet, sie/Sie werden}. Au Präteritum : \all{ich wurde, du wurdest, er wurde, wir wurden, ihr wurdet, sie wurden}.

Le complément d'agent est introduit par \all{von} (+ datif) pour une \textbf{personne} ou une institution, et par \all{durch} (+ accusatif) pour un \textbf{moyen}, une cause :
\all{Die Verfassung wird \textbf{vom} Parlament geändert.} « La constitution est modifiée par le parlement. »
\all{Die Nachricht wurde \textbf{durch} das Radio verbreitet.} « La nouvelle a été diffusée par la radio. »
\end{propriete}

\begin{center}
\begin{tabular}{|l|l|l|}
\hline
\rowcolor{popblueL} Temps & Actif & Passif \\
\hline
Präsens & \all{Das Volk wählt.} & \all{Der Präsident wird gewählt.} \\
\hline
Präteritum & \all{Das Volk wählte.} & \all{Der Präsident wurde gewählt.} \\
\hline
\end{tabular}
\end{center}

\begin{exemplebox}[Exemples traduits au passif]
\all{Die Verfassung wird vom Parlament geändert.} « La constitution est modifiée par le parlement. »

\all{Die Gesetze werden vom Bundestag beschlossen.} « Les lois sont votées par le Bundestag. »

\all{1972 wurden viele junge Kameruner zum Deutschlernen eingeladen.} « En 1972, de nombreux jeunes Camerounais furent invités à apprendre l'allemand. »
\end{exemplebox}

\begin{attention}[L'agent : von ou durch, jamais mit]
Les francophones traduisent souvent « par » mécaniquement. Retenez :
\begin{itemize}
\item \textbf{von} + datif = l'agent (personne, institution) : \all{Der Präsident wird \textbf{vom} Volk gewählt.} « Le président est élu par le peuple. »
\item \textbf{durch} + accusatif = le moyen, l'intermédiaire : \all{Die Stadt wurde \textbf{durch} ein Erdbeben zerstört.} « La ville fut détruite par un tremblement de terre. »
\item \all{mit} exprime l'\emph{instrument} et ne peut jamais introduire l'agent du passif : \all{Der Brief wurde \textbf{mit} einem Stift geschrieben.} « La lettre a été écrite avec un stylo. »
\end{itemize}
\end{attention}

\courssub{Le futur (Futur I) : werden + infinitif}

\textbf{Observation.} \all{Die Bürger werden im September wählen.} « Les citoyens voteront en septembre. » Ici, \all{werden} n'est plus un auxiliaire de passif : il exprime le \textbf{futur}.

\begin{propriete}[Formation et emplois du Futur I]
\textbf{Forme} : \all{werden} conjugué au présent (en position 2) + \textbf{infinitif} du verbe principal renvoyé \textbf{en fin de phrase} :
\all{Die Partei \textbf{wird} die nächste Wahl \textbf{gewinnen}.} « Le parti gagnera les prochaines élections. »

\textbf{Emplois} :
\begin{enumerate}
\item la \textbf{prévision}, l'annonce : \all{Die Wahl wird im Oktober stattfinden.} « L'élection aura lieu en octobre. »
\item la \textbf{promesse}, l'intention : \all{Wir werden die Demokratie verteidigen.} « Nous défendrons la démocratie. »
\item la \textbf{supposition} au présent (souvent avec \all{schon}, \all{wohl}, \all{jetzt}) : \all{Er wird jetzt im Parlament sein.} « Il est probablement au parlement maintenant. »
\end{enumerate}
\textbf{Exception d'usage} : contrairement au français, l'allemand utilise très souvent le \emph{présent} + un complément de temps pour le futur : \all{Morgen \textbf{gehe} ich zur Wahl.} « Demain, j'irai voter. » Le Futur I marque alors plutôt la promesse ou la prédiction solennelle.
\end{propriete}

\begin{attention}[Ne pas confondre passif et futur]
Les deux emploient \all{werden}, mais la fin de phrase trahit tout :
\begin{itemize}
\item \all{werden} + \textbf{Partizip II} = passif : \all{Er wird gewählt.} « Il est élu. »
\item \all{werden} + \textbf{infinitif} = futur : \all{Er wird wählen.} « Il votera. »
\item Et les deux peuvent se combiner (futur du passif) : \all{Er wird gewählt werden.} « Il sera élu. »
\end{itemize}
\end{attention}

\begin{popfigure}
\begin{tikzpicture}
\draw[-{Stealth}, thick, popdark] (0,0) -- (12,0);
\foreach \x/\t/\f in {2/{Präteritum}/{\all{wurde gewählt}}, 6/{Präsens}/{\all{wird gewählt}}, 10/{Futur I}/{\all{wird gewählt werden}}} {
  \fill[poporange] (\x,0) circle (0.09);
  \node[above, align=center, font=\small] at (\x,0.15) {\textbf{\t}\\ \f};
}
\node[below, font=\footnotesize\itshape, popmuted] at (6,-0.25) {Le passif à trois temps : \all{werden} porte la marque du temps, le Partizip II reste en fin de phrase.};
\end{tikzpicture}
\legende{Frise des temps du passif : seul l'auxiliaire \all{werden} change.}
\end{popfigure}

\begin{textebox}[Wahlen in Deutschland — ein kurzer Überblick]
Deutschland ist eine parlamentarische Demokratie, denn das Volk wählt den Bundestag. Alle vier Jahre werden die Abgeordneten vom Volk gewählt, und danach wird der Bundeskanzler vom Bundestag bestimmt. Weil keine Partei meistens allein regieren kann, werden oft Koalitionen gebildet. Die Wahlen sind frei und geheim, deshalb akzeptieren fast alle Bürger das Ergebnis. Trotzdem gibt es Diskussionen: Manche Bürger sind unzufrieden, obwohl sie wählen dürfen. In einer Diktatur ist das anders: Dort wird die Opposition unterdrückt, und die Medien werden vom Staat kontrolliert. Während in einer Demokratie die Verfassung die Grundrechte schützt, werden diese Rechte in einer Diktatur oft missachtet. Wer wird die nächste Wahl gewinnen? Das ist schwer zu sagen, aber eines ist sicher: Die Bürger werden wieder zur Wahl gehen, denn Wählen ist nicht nur ein Recht, sondern auch eine Pflicht. Viele junge Deutsche werden zum ersten Mal wählen, und ihre Stimme wird das Ergebnis beeinflussen. Die Demokratie lebt, weil die Menschen mitmachen.
\end{textebox}

\textbf{Glossaire.} \all{der Abgeordnete, -n} : le député ; \all{bestimmen} : désigner ; \all{die Koalition, -en} : la coalition ; \all{geheim} : secret ; \all{das Ergebnis, -se} : le résultat ; \all{unterdrücken} : réprimer ; \all{missachten} : mépriser, violer ; \all{die Pflicht, -en} : le devoir ; \all{die Stimme, -n} : la voix, le vote ; \all{beeinflussen} : influencer ; \all{mitmachen} : participer.

\textbf{Questions de compréhension.} 1) \all{Wer wählt den Bundestag?} 2) \all{Warum werden oft Koalitionen gebildet?} 3) \all{Was passiert mit der Opposition in einer Diktatur?} 4) Relevez dans le texte deux connecteurs de chaque groupe et deux verbes au passif.

\begin{savaistu}[Le Futur I de supposition — un plus pour le Bac]
Le Futur I sert aussi à exprimer une \textbf{hypothèse sur le présent} : \all{Er wird jetzt im Parlament sein.} « Il est probablement au parlement maintenant. » ; \all{Sie wird wohl krank sein.} « Elle est sans doute malade. » Le français traduit alors par un présent + « probablement ». Ce tour, très fréquent dans les textes de presse, surprend les candidats : ne le traduisez pas par un futur !
\end{savaistu}

\begin{methode}[Du texte à la production : commentaire et expression guidée]
Pour réussir le commentaire et la production écrite sur ce thème, suivez ces étapes :
\begin{enumerate}
\item \textbf{Analyse structurelle} : Repérez les connecteurs (\all{denn, deshalb, weil, obwohl...}) pour identifier la logique du texte (cause, conséquence, concession, opposition). Notez la place du verbe pour valider votre analyse.
\item \textbf{Repérage du passif} : Surlignez les formes \all{werden + Partizip II}. Demandez-vous : quel est le sujet patient ? Y a-t-il un agent (\all{von/durch}) ? Cela révèle qui détient le pouvoir dans le texte.
\item \textbf{Repérage du futur} : Identifiez \all{werden + infinitif}. S'agit-il d'une prévision (\all{stattfinden}), d'une promesse (\all{verteidigen}) ou d'une supposition (\all{sein}) ?
\item \textbf{Réponses aux questions} : Réutilisez le vocabulaire du texte (\all{die Wahl, die Verfassung, die Koalition}) et les structures grammaticales (passif pour l'objectivité, connecteurs pour l'argumentation).
\item \textbf{Production guidée (paragraphe argumentatif)} : 
      \begin{itemize}
      \item Thèse : \all{Meiner Meinung nach ist die Demokratie das beste System.} 
      \item Argument 1 (cause) : \all{Denn sie schützt die Menschenrechte.} 
      \item Argument 2 (concession) : \all{Obwohl sie nicht perfekt ist, ...} 
      \item Conséquence : \all{Deshalb müssen wir wählen gehen.} 
      \item Passif pour la loi : \all{Das Grundgesetz wird vom Parlament geschützt.}
      \item Futur pour l'engagement : \all{Wir werden unsere Stimme abgeben.}
      \end{itemize}
\end{enumerate}
\end{methode}

\begin{exoresolu}[Traduisez en allemand]
Énoncé. 1) « Le président est élu par le peuple. » 2) « Les élections étaient libres, c'est pourquoi il n'y a pas eu de manifestations. » 3) « Les citoyens voteront en septembre, parce que la constitution le prévoit. »
\tcblower
Solution détaillée.
\begin{enumerate}
\item Passif au présent, agent = personne, donc \all{von} : \all{Der Präsident wird \textbf{vom} Volk gewählt.}
\item \all{deshalb} en position 1, inversion ; passif au Präteritum pour « il y a eu » : \all{Die Wahlen waren frei, \textbf{deshalb gab es} keine Proteste.}
\item Futur I (infinitif en fin) + subordonnée avec \all{weil} (verbe à la fin) : \all{Die Bürger \textbf{werden} im September \textbf{wählen}, weil die Verfassung das \emph{vorsieht}.}
\end{enumerate}
\end{exoresolu}

\begin{exercice}[À vous]
1) Complétez avec le bon connecteur (\all{denn}, \all{deshalb}, \all{obwohl}) : \all{Das Volk durfte nicht wählen, ... gab es Proteste.} / \all{... die Wahl frei war, blieben viele Bürger zu Hause.} / \all{Die Partei fürchtet die Presse, ... sie kontrolliert die Medien.}
2) Mettez au passif (Präsens puis Präteritum) : \all{Das Parlament ändert die Verfassung.}
3) Traduisez : « Bien que la démocratie ne soit pas parfaite, elle protège la liberté. »
\end{exercice}

\begin{corrige}[Exercice « À vous »]
1) \all{..., \textbf{deshalb} gab es Proteste.} (inversion) ; \all{\textbf{Obwohl} die Wahl frei war, ...} (verbe \all{war} en fin) ; \all{..., \textbf{denn} sie kontrolliert die Medien.} (ordre normal).
2) Präsens : \all{Die Verfassung wird vom Parlament geändert.} Präteritum : \all{Die Verfassung wurde vom Parlament geändert.}
3) \all{Obwohl die Demokratie nicht perfekt ist, schützt sie die Freiheit.} (subordonnée en tête : verbe \all{ist} en fin, puis inversion dans la principale : \all{schützt sie}.)
\end{corrige}

\coursec{Méthode du commentaire et production guidée}

Dans la section précédente, nous avons appris à manier les connecteurs logiques (\all{weil}, \all{deshalb}, \all{zuerst}, \all{danach}...), le passif (\all{die Opposition wird unterdrückt}) et le futur (\all{das Land wird sich entwickeln}). Ces outils ne sont pas de simples exercices de grammaire : ce sont précisément les instruments qui permettent de \cle{commenter} un texte, de le \cle{résumer} et de donner son \cle{avis} de manière structurée. Cette dernière section transforme donc vos acquis grammaticaux en compétence d'expression : le commentaire de texte, exercice central de la classe de Terminale et épreuve décisive au Baccalauréat.

\courssub{La méthode du commentaire de texte en 4 étapes}

Commenter un texte allemand, ce n'est ni le traduire, ni le paraphraser : c'est le présenter, en dégager la structure, analyser sa langue et donner un avis argumenté. Voici la démarche complète, en quatre étapes obligatoires.

\begin{methode}[Commenter un texte : les 4 étapes]
\begin{enumerate}
\item \textbf{Présenter le texte} (\all{der Text}) : indiquer le thème, la source et le type de texte (article de presse, discours, extrait de manuel...). Phrases-gabarits : \all{Der Text handelt von...} / \all{Der Text informiert über...} / \all{Es ist ein Zeitungsartikel über...}.
\item \textbf{Dégager le plan et les idées principales} (\all{der Aufbau}) : repérer les grandes parties et les annoncer avec des connecteurs d'ordre : \all{Zuerst erklärt der Autor..., dann beschreibt er..., danach spricht er über..., schließlich...}.
\item \textbf{Analyser le vocabulaire et les procédés} (\all{die Sprache}) : relever le champ lexical (ici : \all{die Wahl, die Verfassung, die Macht, die Partei}), le style nominal (\all{die Unterdrückung der Opposition} au lieu de \all{die Opposition wird unterdrückt}), les temps et les connecteurs utilisés.
\item \textbf{Donner un avis personnel argumenté} (\all{die eigene Meinung}) : opinion + justification + exemple, avec \all{Meiner Meinung nach...}, \all{Ich bin der Ansicht, dass...}, \all{weil / denn...}, \all{zum Beispiel...}.
\end{enumerate}
\end{methode}

\begin{popfigure}
\begin{tikzpicture}[node distance=0.55cm,
  box/.style={draw=popblue, fill=popblueL, rounded corners, thick, align=center, minimum width=4.2cm, minimum height=1.1cm, text width=4cm},
  arr/.style={-{Stealth[length=3mm]}, thick, poporange}]
\node[box, align=center] (e1){\textbf{1. Présentation}\\ \all{Der Text handelt von...}};
\node[box, below=of e1, align=center] (e2){\textbf{2. Plan / Idées}\\ \all{Zuerst..., dann..., danach...}};
\node[box, below=of e2, align=center] (e3){\textbf{3. Analyse}\\ vocabulaire, style nominal, connecteurs};
\node[box, below=of e3, align=center] (e4){\textbf{4. Avis personnel}\\ \all{Meiner Meinung nach...}};
\draw[arr] (e1) -- (e2);
\draw[arr] (e2) -- (e3);
\draw[arr] (e3) -- (e4);
\end{tikzpicture}
\legende{Organigramme : les 4 étapes enchaînées du commentaire de texte.}
\end{popfigure}

\begin{propriete}[Les verbes de présentation et leurs prépositions]
Chaque verbe de présentation impose sa préposition ; on ne les mélange pas :
\begin{itemize}
\item \all{handeln von + Datif} : \all{Der Text handelt von den politischen Systemen.} — Le texte traite des systèmes politiques.
\item \all{informieren über + Accusatif} : \all{Der Text informiert über die Demokratie.} — Le texte informe sur la démocratie.
\item \all{sprechen über + Accusatif} : \all{Der Autor spricht über die Wahl.} — L'auteur parle de l'élection.
\item \all{vorstellen + Accusatif} (verbe à particule séparable) : \all{Der Text stellt drei Systeme vor.} — Le texte présente trois systèmes.
\end{itemize}
\end{propriete}

\begin{aretenir}[Expressions-gabarits du commentaire]
\begin{itemize}
\item Présentation : \all{Der Text handelt von den politischen Systemen.} — Le texte traite des systèmes politiques.
\item Intention : \all{Der Autor zeigt, dass die Demokratie auf freien Wahlen beruht.} — L'auteur montre que la démocratie repose sur des élections libres.
\item Plan : \all{Zuerst erklärt er die Demokratie, dann die Diktatur.} — D'abord il explique la démocratie, ensuite la dictature.
\item Avis : \all{Meiner Meinung nach ist die Demokratie das beste System.} — À mon avis, la démocratie est le meilleur système.
\item Avis renforcé : \all{Ich bin der Ansicht, dass die Bürger wählen müssen.} — Je suis d'avis que les citoyens doivent voter.
\end{itemize}
\end{aretenir}

\begin{attention}[La place du verbe dans les phrases-gabarits]
Après \all{dass}, le verbe conjugué va à la fin : \all{Ich bin der Ansicht, dass die Bürger ihre Regierung frei wählen \textbf{können}.} Ne dites jamais \all{dass die Bürger können wählen}. De même, après \all{Meiner Meinung nach}, le verbe reste en deuxième position : \all{Meiner Meinung nach \textbf{ist} die Demokratie...} (et non \all{Meiner Meinung nach die Demokratie ist}). Attention aussi à \all{beruhen auf + Datif} : \all{auf freien Wahlen beruht} (et non \all{auf freie Wahlen}).
\end{attention}

\courssub{Commentaire modèle rédigé}

Voici un commentaire complet du texte support \all{Politische Systeme in der Welt}, construit selon les quatre étapes. Observez les connecteurs (ici en gras) qui articulent le propos.

\begin{textebox}[Commentaire modèle : „Politische Systeme in der Welt“]
Der Text „Politische Systeme in der Welt“ informiert über drei Staatsformen: die Demokratie, die Diktatur und die Monarchie. \textbf{Zuerst} erklärt der Autor die Demokratie am Beispiel Deutschlands, \textbf{dann} beschreibt er die Diktatur, wo die Opposition unterdrückt wird. \textbf{Danach} spricht er über die Monarchie. Der Text benutzt viele Wörter aus dem politischen Bereich, zum Beispiel die Wahl, die Verfassung und die Macht. \textbf{Meiner Meinung nach} ist die Teilnahme der Bürger sehr wichtig, \textbf{weil} ein Volk ohne politische Rechte keine Zukunft hat.
\end{textebox}

\textbf{Glossaire} : die Staatsform, -en (forme d'État) ; unterdrücken (réprimer) ; der Bereich, -e (domaine) ; die Teilnahme, -n (participation) ; das Recht, -e (droit).

\textbf{Analyse de la structure du modèle} :
\begin{enumerate}
\item \textbf{Présentation} : première phrase (\all{informiert über drei Staatsformen}).
\item \textbf{Plan} : phrases 2 et 3, avec \all{zuerst}, \all{dann}, \all{danach}.
\item \textbf{Analyse de la langue} : phrase 4 (champ lexical politique).
\item \textbf{Avis personnel} : dernière phrase, avec \all{Meiner Meinung nach} + justification en \all{weil} (verbe \all{hat} en fin de subordonnée).
\end{enumerate}

\courssub{Le résumé (die Zusammenfassung)}

Le résumé obéit à des règles strictes, différentes du commentaire : il ne contient \textbf{aucun avis personnel}.

\begin{propriete}[Les 5 règles d'or du résumé allemand]
\begin{enumerate}
\item Rédiger au \textbf{présent} (\all{das Präsens}), même si le texte est au passé.
\item \textbf{Aucune citation} : on ne recopie jamais de phrases entières du texte.
\item \textbf{Aucun avis personnel} : pas de \all{ich finde}, pas de \all{meiner Meinung nach}.
\item Reformuler \textbf{avec ses propres mots} : synonymes et style nominal (\all{die Wahl der Regierung} pour \all{die Bürger wählen die Regierung}).
\item Longueur : \textbf{3 à 5 phrases}, qui reprennent uniquement les idées principales.
\end{enumerate}
\end{propriete}

\begin{attention}[L'erreur classique du francophone]
Dans un résumé, l'élève recopie souvent des phrases du texte ou ajoute « je pense que... ». Les deux sont sanctionnés. La bonne stratégie : repérer les idées principales, les nominaliser, puis les relier par des connecteurs neutres (\all{zuerst}, \all{dann}, \all{außerdem}, \all{schließlich}). Comparez : \all{Die Bürger wählen ihre Regierung frei} (phrase du texte) devient dans le résumé \all{der Autor erklärt die freie Wahl der Regierung durch die Bürger}. Attention aussi au verbe à particule \all{zusammenfassen} : \all{Der Schüler fasst den Text zusammen.} (particule \all{zusammen} rejetée à la fin).
\end{attention}

\begin{exemplebox}[Résumé modèle du texte support]
\all{Der Text stellt drei politische Systeme vor: die Demokratie, die Diktatur und die Monarchie. Zuerst erklärt er die freie Wahl der Regierung in der Demokratie. Dann beschreibt er die Unterdrückung der Opposition in der Diktatur. Schließlich nennt er die Rolle des Königs in der Monarchie.}

« Le texte présente trois systèmes politiques : la démocratie, la dictature et la monarchie. D'abord il explique l'élection libre du gouvernement en démocratie. Ensuite il décrit la répression de l'opposition dans la dictature. Enfin il mentionne le rôle du roi dans la monarchie. » — Remarquez le style nominal (\all{die freie Wahl}, \all{die Unterdrückung}, \all{die Rolle}), le verbe séparable \all{stellt ... vor} et l'absence totale d'avis personnel.
\end{exemplebox}

\courssub{Donner son avis : la structure en 3 temps}

\begin{methode}[Un avis personnel solide : Meinung → Begründung → Beispiel]
\begin{enumerate}
\item \textbf{Meinung} (opinion) : \all{Meiner Meinung nach ist die Demokratie das beste System.} ou \all{Ich bin der Ansicht, dass...}.
\item \textbf{Begründung} (justification) : \all{weil die Bürger ihre Regierung frei wählen können} (verbe à la fin) ou \all{denn die Bürger können ihre Regierung frei wählen} (verbe en position 2).
\item \textbf{Beispiel} (exemple) : \all{Zum Beispiel wählen die Deutschen alle vier Jahre den Bundestag.}
\end{enumerate}
\end{methode}

\begin{exemplebox}[Avis complet, traduit]
\all{Meiner Meinung nach ist die Demokratie das beste System, weil die Bürger ihre Regierung frei wählen können. Zum Beispiel können die Menschen in Deutschland oder in Kamerun ihre Abgeordneten wählen.}

« À mon avis, la démocratie est le meilleur système parce que les citoyens peuvent élire librement leur gouvernement. Par exemple, les gens en Allemagne ou au Cameroun peuvent élire leurs députés. » — \all{der Abgeordnete, -n} (le député) : adjectif substantivé, ici accusatif pluriel \all{ihre Abgeordneten}.
\end{exemplebox}

\begin{center}
\begin{tabularx}{\linewidth}{|X|X|X|}
\hline
\rowcolor{popblueL} Fonction & Deutsch & Français \\
\hline
Présenter & \all{Der Text handelt von...} & Le texte traite de... \\
\hline
Montrer & \all{Der Autor zeigt, dass...} & L'auteur montre que... \\
\hline
Ordonner & \all{zuerst, dann, danach, schließlich} & d'abord, ensuite, puis, enfin \\
\hline
Opiner & \all{Meiner Meinung nach...} & À mon avis... \\
\hline
Opiner (soutenu) & \all{Ich bin der Ansicht, dass...} & Je suis d'avis que... \\
\hline
Justifier & \all{weil... / denn...} & parce que... / car... \\
\hline
Illustrer & \all{zum Beispiel (z. B.)} & par exemple \\
\hline
Conclure & \all{Zusammenfassend kann man sagen...} & En résumé, on peut dire... \\
\hline
\end{tabularx}
\end{center}

\begin{attention}[Faux amis et pièges de registre]
\begin{itemize}
\item \all{die Meinung} = l'opinion, l'avis ; \all{Meiner Meinung nach} se construit avec le nom \textbf{avant} \all{nach} (postposition).
\item \all{behandeln} signifie « traiter (un sujet) » : \all{Der Text behandelt die Demokratie.} — ne pas traduire par « soigner » ici.
\item Évitez \all{ich denke} dans un devoir écrit : préférez \all{ich bin der Ansicht} ou \all{meiner Meinung nach}, plus soutenus.
\end{itemize}
\end{attention}

\courssub{Exercice résolu : rédiger un court commentaire}

\begin{exoresolu}[Commentaire guidé]
Consigne : \all{Schreiben Sie einen kurzen Kommentar (5-6 Sätze) zum Text „Politische Systeme in der Welt“. Benutzen Sie mindestens zwei Konnektoren und einen Satz im Passiv.} — Rédigez un court commentaire (5-6 phrases) sur le texte support, avec au moins deux connecteurs et une phrase au passif.
\tcblower
\textbf{Solution détaillée.} On suit les 4 étapes. (1) Présentation : \all{Der Text handelt von den politischen Systemen der Welt.} (2) Plan : \all{Zuerst erklärt der Autor die Demokratie, dann die Diktatur und die Monarchie.} (3) Analyse avec passif : \all{In der Diktatur wird die Opposition unterdrückt, und der Text benutzt viele politische Wörter wie die Wahl und die Macht.} (4) Avis : \all{Meiner Meinung nach ist die Demokratie das beste System, weil die Bürger frei wählen können.} Vérification : deux connecteurs (\all{zuerst}, \all{dann}, \all{weil}) et un passif (\all{wird unterdrückt}) sont bien présents ; le verbe après \all{weil} est en fin de phrase.
\end{exoresolu}

\courssub{Production guidée : résumé + paragraphe d'opinion}

\begin{exercice}[Production écrite guidée]
\begin{enumerate}
\item \textbf{Résumé} : \all{Fassen Sie den Text „Politische Systeme in der Welt“ in 3 bis 5 Sätzen zusammen.} Résumez le texte support en 3 à 5 phrases, au présent, sans citation ni avis personnel, en utilisant le style nominal.
\item \textbf{Opinion} : \all{Welches politische System ist das beste für ein Land? Schreiben Sie einen Absatz (5-6 Sätze) mit mindestens zwei Konnektoren und einem Satz im Futur oder im Passiv.} Quel système politique est le meilleur pour un pays ? Rédigez un paragraphe de 5-6 phrases avec au moins deux connecteurs et une phrase au futur ou au passif.
\end{enumerate}
\end{exercice}

\begin{corrige}[Production guidée]
\textbf{1. Résumé possible} : \all{Der Text informiert über die Demokratie, die Diktatur und die Monarchie. Zuerst erklärt er die freie Wahl der Regierung in der Demokratie. Dann beschreibt er die Unterdrückung der Opposition in der Diktatur. Schließlich spricht er über die Rolle des Königs.} (Présent, style nominal, aucun avis, 4 phrases.)

\textbf{2. Paragraphe d'opinion possible} : \all{Meiner Meinung nach ist die Demokratie das beste System für ein Land, weil die Bürger ihre Regierung frei wählen können. Außerdem schützt die Verfassung die Menschenrechte. In einer Diktatur wird die Opposition unterdrückt, und das ist schlecht für die Entwicklung. Ich bin der Ansicht, dass die Jugend mehr lernen wird, wenn sie politisch aktiv ist. Deshalb wird die Demokratie in vielen Ländern Afrikas stärker werden.} — Vérification : connecteurs \all{weil}, \all{außerdem}, \all{deshalb} ; un passif (\all{wird unterdrückt}) et un futur (\all{wird stärker werden}) ; après \all{deshalb}, inversion sujet-verbe (\all{wird die Demokratie}).
\end{corrige}

\courssub{Expression orale guidée : le mini-débat}

\begin{experience}[Mini-débat en binômes : „Demokratie oder Diktatur ?“]
Par groupes de deux, chaque élève défend une position à l'aide des cartes d'arguments ci-dessous. Durée : 2 minutes par élève, puis échange de rôles.
\begin{itemize}
\item \textbf{Carte Freiheit} (liberté) : \all{In der Demokratie sind die Menschen frei.} / \all{die Meinungsfreiheit} (la liberté d'expression).
\item \textbf{Carte Sicherheit} (sécurité) : \all{Eine starke Regierung garantiert die Sicherheit.}
\item \textbf{Carte Entwicklung} (développement) : \all{Die Demokratie hilft der wirtschaftlichen Entwicklung.}
\item \textbf{Carte Menschenrechte} (droits de l'homme) : \all{Die Verfassung schützt die Menschenrechte.}
\end{itemize}
Structure imposée : \all{Meiner Meinung nach... weil... Zum Beispiel...} L'autre élève réplique avec \all{Ich bin anderer Ansicht, denn...} (Je suis d'un autre avis, car...). Remarquez : \all{helfen} se construit avec le \textbf{datif} : \all{Die Demokratie hilft der Entwicklung.}
\end{experience}

\begin{savaistu}[Le débat, une tradition germanophone]
En Allemagne, en Autriche et en Suisse, le débat argumenté fait partie de la culture scolaire : les lycéens y apprennent très tôt à formuler une opinion avec \all{Meiner Meinung nach} et à la justifier. La Suisse pratique même la démocratie directe : les citoyens votent régulièrement par référendum (\all{die Volksabstimmung, -en}) sur des questions précises. Au Cameroun, les clubs de débat et les journaux scolaires des lycées germanophones reprennent cette tradition : c'est un excellent entraînement pour l'épreuve orale du Baccalauréat.
\end{savaistu}

\begin{aretenir}[À retenir pour l'examen]
Un bon commentaire = 4 étapes (présentation, plan, analyse, avis) + connecteurs + verbes bien placés. Un bon résumé = présent, propres mots, style nominal, zéro avis personnel. Un bon avis = \all{Meinung} + \all{Begründung} + \all{Beispiel}. Avec ces trois outils et le vocabulaire politique de la leçon (\all{die Demokratie, die Wahl, die Verfassung, die Macht}), vous êtes prêts pour l'épreuve d'expression écrite et orale de la Terminale A.
\end{aretenir}

\newpage

\coursec{Activité d'intégration}

\coursec{AL24 : Les systèmes politiques}

\courssub{Texte support : „Politische Systeme im Vergleich – Deutschland und Kamerun“}
\begin{textebox}[Texte original -- Niveau B1/B2]
\textbf{Politische Systeme im Vergleich: Deutschland und Kamerun}

\all{Die Bundesrepublik Deutschland ist eine parlamentarische Demokratie. Das Staatsoberhaupt ist der Bundespräsident, er hat aber vor allem repräsentative Aufgaben. Die politische Macht liegt beim Bundeskanzler und der Bundesregierung. Der Bundestag wird alle vier Jahre nach dem Prinzip der personalisierten Verhältniswahl gewählt. Die Wähler haben zwei Stimmen: die Erststimme für einen Kandidaten im Wahlkreis und die Zweitstimme für eine Partei. Da keine Partei meist die absolute Mehrheit erhält, sind Koalitionen notwendig. Die Opposition kontrolliert die Regierung im Parlament.}

\all{Die Republik Kamerun ist eine präsidiale Republik. Der Präsident wird alle sieben Jahre direkt vom Volk gewählt und verfügt über weitreichende Exekutivbefugnisse. Er ernennt den Premierminister und die Minister. Das Parlament besteht aus zwei Kammern: der Nationalversammlung und dem Senat. Die Abgeordneten der Nationalversammlung werden in Wahlkreisen nach einem Mehrheitswahlsystem gewählt. Seit der Einführung des Mehrparteiensystems 1990 gibt es viele Parteien, doch die RDPC (Rassemblement Démocratique du Peuple Camerounais) dominiert die politische Landschaft seit Jahrzehnten. Die Verfassung von 1996, revidiert 2008, garantiert formell die Demokratie und die Menschenrechte.}

\all{Beide Länder kennen Wahlen als Legitimationsgrundlage der Macht. In Deutschland führt das Wahlsystem zu Kompromissen und Koalitionen. In Kamerun sichert das Mehrheitswahlsystem oft die Stabilität der Regierungspartei. Internationale Beobachter begleiten die Wahlen in beiden Staaten, um die Einhaltung der Verfassung und die Transparenz des Wahlprozesses zu sichern.}
\end{textebox}

\courssub{Compréhension du texte}
\begin{exercice}[Compréhension globale]
\begin{enumerate}
    \item Quel est le type de texte ? (article de presse, extrait de manuel, rapport officiel, essai personnel)
    \item Quels sont les deux pays comparés ?
    \item Quel est le thème central du texte ?
\end{enumerate}
\end{exercice}

\begin{exercice}[Compréhension détaillée]
\textbf{Allemagne :} Répondez en allemand par des phrases complètes.
\begin{enumerate}
    \item Welche Art von Demokratie ist die Bundesrepublik Deutschland?
    \item Wer hat die politische Macht in Deutschland?
    \item Wie oft wird der Bundestag gewählt?
    \item Erklären Sie das Wahlsystem (Erststimme / Zweitstimme).
    \item Warum sind Koalitionen in Deutschland notwendig?
\end{enumerate}
\textbf{Cameroun :} Répondez en allemand par des phrases complètes.
\begin{enumerate}
    \setcounter{enumi}{5}
    \item Welche Staatsform hat Kamerun?
    \item Wie oft wird der Präsident gewählt und von wem?
    \item Aus welchen zwei Kammern besteht das Parlament?
    \item Nach welchem Wahlsystem werden die Abgeordneten der Nationalversammlung gewählt?
    \item Welche Partei dominiert die Politik seit 1990?
\end{enumerate}
\textbf{Comparaison :}
\begin{enumerate}
    \setcounter{enumi}{10}
    \item Nennen Sie einen Unterschied und eine Gemeinsamkeit der beiden Wahlsysteme.
    \item Welche Rolle spielen internationale Beobachter?
\end{enumerate}
\end{exercice}

\begin{corrige}[Exercice Compréhension]
\begin{enumerate}
    \item Es ist ein informativer Fachtext / ein Vergleichstext.
    \item Deutschland und Kamerun.
    \item Politische Systeme / Demokratieformen und Wahlsysteme.
    \item Deutschland ist eine parlamentarische Demokratie.
    \item Die politische Macht liegt beim Bundeskanzler und der Bundesregierung.
    \item Der Bundestag wird alle vier Jahre gewählt.
    \item Die Erststimme gilt einem Kandidaten im Wahlkreis, die Zweitstimme einer Partei (Landesliste).
    \item Weil keine Partei meist die absolute Mehrheit erhält (Verhältniswahl).
    \item Kamerun ist eine präsidiale Republik.
    \item Der Präsident wird alle sieben Jahre direkt vom Volk gewählt.
    \item Aus der Nationalversammlung und dem Senat.
    \item Nach einem Mehrheitswahlsystem in Wahlkreisen.
    \item Die RDPC (Rassemblement Démocratique du Peuple Camerounais).
    \item Unterschied: Deutschland hat Verhältniswahl (Koalitionen), Kamerun Mehrheitswahl (Dominanz einer Partei). Gemeinsamkeit: Beide wählen ihr Parlament und ihren Präsidenten (bzw. Kanzler indirekt) durch Wahlen. / Beide haben eine Verfassung.
    \item Sie sichern die Einhaltung der Verfassung und die Transparenz des Wahlprozesses.
\end{enumerate}
\end{corrige}

\courssub{Lexique thématique : Die politischen Systeme}
\begin{propriete}[Vocabulaire de base -- Wortschatz: Politische Systeme (mit Artikel, Plural, Derivaten)]
\begin{center}
\begin{tabularx}{\linewidth}{|X|X|X|}
\hline
\rowcolor{popblueL} \textbf{Deutsch (Artikel, Plural)} & \textbf{Français} & \textbf{Famille de mots / Collocations} \\
\hline
die Demokratie, -n & la démocratie & demokratisch (adj.), demokratisieren (v.), der Demokrat (-en) \\
\hline
die Diktatur, -en & la dictature & diktatorisch (adj.), der Diktator (-en), diktieren (v.) \\
\hline
die Wahl, -en & l'élection, le scrutin & wählen (v., irr.: wählt, wählte, hat gewählt), der Wähler (-), die Wählerin (-nen), die Wahlurne (-n), der Wahlkampf (-e) \\
\hline
die Partei, -en & le parti politique & parteiisch (adj.), der Parteitag (-e), die Parteizentrale (-n), parteiunabhängig \\
\hline
der Präsident, -en & le président & die Präsidentin (-nen), präsidial (adj.), das Präsidium (-ien), präsidieren (v.) \\
\hline
die Verfassung, -en & la constitution & verfassungsgemäß / verfassungswidrig (adj.), verfassungsgebend (adj.) \\
\hline
das Parlament, -e & le parlement & parlamentarisch (adj.), der Parlamentarier (-), die Parlamentswahl (-en) \\
\hline
die Stimme, -n & la voix (vote) & stimmen (v.), stimmberechtigt (adj.), die Stimmabgabe, die Stimmenthaltung \\
\hline
der Wahlkreis, -e & la circonscription électorale & der Wahlkreisbewerber (-), die Wahlkreiseinteilung \\
\hline
die Koalition, -en & la coalition & koalieren (v.), der Koalitionsvertrag (-e), die Koalitionsverhandlungen \\
\hline
die Opposition & l'opposition & oppositionell (adj.), der Oppositionsführer (-), die Oppositionsarbeit \\
\hline
das Grundgesetz, -e & la Loi fondamentale (ALL) & grundgesetzlich (adj.), Artikel des Grundgesetzes \\
\hline
der Bundeskanzler, - & le chancelier fédéral & die Bundeskanzlerin (-nen), das Bundeskanzleramt \\
\hline
der Bundespräsident, -en & le président fédéral (ALL) & das Bundespräsidialamt \\
\hline
die Legislative & le pouvoir législatif & legislativ (adj.), der Gesetzgeber (-) \\
\hline
die Exekutive & le pouvoir exécutif & exekutiv (adj.), die Exekutivgewalt \\
\hline
die Judikative & le pouvoir judiciaire & judikativ (adj.), die richterliche Gewalt \\
\hline
\end{tabularx}
\end{center}
\emph{Exercice :} Former 5 noms composés avec \cle{Wahl-}, \cle{Partei-}, \cle{Verfassungs-}, \cle{Präsidenten-}, \cle{Koalitions-} (ex. \all{der Wahlkampf}, \all{die Parteizentrale}).
\end{propriete}

\begin{exercice}[Entraînement lexical]
\begin{enumerate}
    \item Complétez en allemand : \all{In einer \_\_\_\_\_\_\_\_\_\_ wählt das Volk seine Vertreter.} (Demokratie) \\
    \item Transformez en nom (Nominalstil) : \all{die Parteien koalieren} $\rightarrow$ \all{die \_\_\_\_\_\_\_\_\_\_ der Parteien} (Koalition) \\
    \item Donnez le Partizip II : \all{wählen} $\rightarrow$ \_\_\_\_\_\_\_\_\_\_ (gewählt), \all{bestimmen} $\rightarrow$ \_\_\_\_\_\_\_\_\_\_ (bestimmt), \all{stattfinden} $\rightarrow$ \_\_\_\_\_\_\_\_\_\_ (stattgefunden).
    \item Traduisez : « Le président est élu pour sept ans. » $\rightarrow$ \all{Der Präsident wird für sieben Jahre \_\_\_\_\_\_\_\_\_\_.} (gewählt)
\end{enumerate}
\end{exercice}

\courssub{Grammaire 1 : Le style nominal (Der Nominalstil)}
\begin{propriete}[Définition et formation]
Le \cle{Nominalstil} (style nominal) transforme des verbes, adjectifs ou phrases entières en noms (souvent des noms déverbaux en \cle{-ung}, \cle{-ung}, \cle{-heit}, \cle{-keit}, \cle{-nis}, \cle{-tion}, ou simples infinitifs substantivés). Il confère au texte une densité informationnelle, de l'objectivité et de la formalité (presse, administration, sciences).
\begin{center}
\begin{tabularx}{\linewidth}{|X|X|X|}
\hline
\rowcolor{popblueL} \textbf{Verbalstil (oral, narratif)} & \textbf{Transformation} & \textbf{Nominalstil (écrit, formel)} \\
\hline
\all{Die Parteien wählen die Kandidaten.} & Verbe $\rightarrow$ Nom déverbal + \cle{Genitiv} & \all{Die Wahl der Kandidaten durch die Parteien} \\
\hline
\all{Weil das Volk abstimmt, ...} & Conjonction \cle{weil} $\rightarrow$ Préposition \cle{wegen} + \cle{Genitiv} & \all{Wegen der Abstimmung des Volkes ...} \\
\hline
\all{Obwohl es regnet, ...} & Conjonction \cle{obwohl} $\rightarrow$ Préposition \cle{trotz} + \cle{Genitiv} & \all{Trotz des Regens ...} \\
\hline
\all{Der Präsident wird gewählt.} & Passif $\rightarrow$ Nom déverbal + \cle{Genitiv} (sujet logique) & \all{Die Wahl des Präsidenten} \\
\hline
\all{Die Regierung regiert das Land.} & Sujet agent $\rightarrow$ \cle{durch} + Akk / \cle{von} + Dat (rare) & \all{Die Regierung des Landes durch die Regierung} (maladroit) $\rightarrow$ mieux : \all{Die Regierungsführung} \\
\hline
\end{tabularx}
\end{center}
\end{propriete}

\begin{propriete}[Règles clés pour le Bac]
\begin{itemize}
    \item \textbf{Majuscule obligatoire :} Tout nom (substantivé) prend une majuscule : \all{das Wählen} $\rightarrow$ \all{das Wählen} (infinitif substantivé), \all{die Wahl}.
    \item \textbf{Genitif du sujet logique :} \all{Die Wahl \textbf{des Präsidenten}} (le président est celui qu'on élit). \all{Die Ernennung \textbf{des Ministers}}.
    \item \textbf{Prépositions au lieu de conjonctions :}
    \begin{itemize}
        \item \cle{weil} (cause) $\rightarrow$ \cle{wegen} + Gen / \cle{aufgrund} + Gen / \cle{infolge} + Gen.
        \item \cle{obwohl} (concession) $\rightarrow$ \cle{trotz} + Gen.
        \item \cle{damit} (but) $\rightarrow$ \cle{zum} + Infinitif substantivé (\all{zur Wahl}) / \cle{zu} + Infinitif.
        \item \cle{daß} (contenu) $\rightarrow$ \cle{das} + Nom (\all{die Tatsache, dass...} $\rightarrow$ \all{die Tatsache der Wahl}).
    \end{itemize}
    \item \textbf{Verbes supports :} Pour réintroduire un verbe : \cle{stattfinden} (avoir lieu), \cle{durchführen} (effectuer), \cle{erfolgen} (se produire), \cle{gelten} (valoir). Ex: \all{Die Wahl findet statt.} / \all{Die Wahl wird durchgeführt.}
\end{itemize}
\end{propriete}

\begin{exemplebox}[Exemples d'application pour l'article]
\begin{itemize}
    \item \all{Weil die Verfassung die Wahl regelt.} $\rightarrow$ \all{Aufgrund der verfassungsmäßigen Regelung der Wahl.}
    \item \all{Obwohl die Opposition schwach ist, gibt es Wahlen.} $\rightarrow$ $\rightarrow$ \all{Trotz der Schwäche der Opposition finden Wahlen statt.}
    \item \all{Die Wähler geben ihre Stimme ab.} $\rightarrow$ \all{Die Stimmabgabe der Wähler.}
    \item \all{Der Präsident übt die Macht aus.} $\rightarrow$ \all{Die Machtausübung des Präsidenten.}
\end{itemize}
\end{exemplebox}

\courssub{Grammaire 2 : Connecteurs logiques (Konnektoren) \& Structure de phrase}
\begin{propriete}[Tableau récapitulatif -- Place du verbe]
\begin{center}
\begin{tabularx}{\linewidth}{|l|X|X|X|}
\hline
\rowcolor{popblueL} \textbf{Type} & \textbf{Connecteur} & \textbf{Sens / Fonction} & \textbf{Structure (Place du verbe)} \\
\hline
\textbf{Koordination} (0 / Position 2) & \cle{aber} & opposition / mais & Satz 1, \all{aber} Satz 2 (V2). \\
& \cle{dennoch} / \cle{trotzdem} & concession / pourtant & Satz 1. \all{Dennoch} Satz 2 (V2). \\
& \cle{deshalb} / \cle{deswegen} & conséquence / donc & Satz 1. \cle{Deshalb} Satz 2 (V2). \\
& \cle{zudem} / \cle{außerdem} & addition / de plus & Satz 1. \cle{Zudem} Satz 2 (V2). \\
& \cle{hingegen} / \cle{im Gegensatz dazu} & contraste / alors que & Satz 1. \cle{Hingegen} Satz 2 (V2). \\
\hline
\textbf{Subordination} (Verbe à la fin) & \cle{weil} & cause / parce que & \cle{Weil} ... \textbf{verbe\_fin}. \\
& \cle{obwohl} & concession / bien que & \cle{Obwohl} ... \textbf{verbe\_fin}. \\
& \cle{dass} & que (complément) & \cle{Dass} ... \textbf{verbe\_fin}. \\
& \cle{damit} & but / afin que & \cle{Damit} ... \textbf{verbe\_fin}. \\
& \cle{falls} / \cle{wenn} & condition / si & \cle{Falls} ... \textbf{verbe\_fin}. \\
\hline
\textbf{Prépositions} (Cas : Génitif) & \cle{wegen} + Gen & à cause de & \all{Wegen des Regens} (Groupe nominal, pas de verbe). \\
& \cle{trotz} + Gen & malgré & \all{Trotz der Diktatur} ... \\
& \cle{aufgrund} + Gen & en raison de & \all{Aufgrund der Verfassung} ... \\
& \cle{während} + Gen & pendant / alors que & \all{Während der Sitzung} ... \\
\hline
\end{tabularx}
\end{center}
\emph{Stratégie pour la comparaison :} Utilisez \cle{Obwohl} (concession: \all{Obwohl Kamerun eine Demokratie ist...}) + \cle{deshalb} (conséquence: \all{...deshalb fehlt Alternanz.}) ou \cle{hingegen} (contraste: \all{In Deutschland hingegen...}) + \cle{weil} (cause: \all{...weil das System Kompromisse erfordert.}).
\end{propriete}

\courssub{Conjugaison : Passif (Vorgangspassiv) \& Futur I}
\begin{propriete}[Passif Présent (Präsens Passiv) -- Processus, action subie]
\cle{werden} (conjugué, position 2) + \cle{Partizip II} (à la fin de la phrase / proposition).
\begin{center}
\begin{tabular}{|l|l|l|l|}
\hline
\rowcolor{popblueL} \textbf{Personne} & \textbf{werden} & \textbf{Exemple : wählen $\rightarrow$ gewählt} & \textbf{Traduction} \\
\hline
ich & werde & \all{ich werde gewählt} & je suis élu \\
du & wirst & \all{du wirst gewählt} & tu es élu \\
\cle{er/sie/es} & \cle{wird} & \cle{er wird gewählt} & \textbf{il est élu} \\
wir & werden & \all{wir werden gewählt} & nous sommes élus \\
ihr & werdet & \all{ihr werdet gewählt} & vous êtes élus \\
\cle{sie/Sie} & \cle{werden} & \cle{sie werden gewählt} & \textbf{ils sont élus} \\
\hline
\end{tabular}
\end{center}
\textbf{Agent (optionnel) :} \cle{von} + Datif (personne) / \cle{durch} + Accusatif (instrument/moyen/institution). Ex: \all{Der Präsident wird \textbf{vom Volk} gewählt.} / \all{Das Gesetz wird \textbf{durch das Parlament} beschlossen.}
\end{propriete}

\begin{propriete}[Futur I -- Prédiction, annonce officielle, promesse]
\cle{werden} (conjugué, position 2) + \cle{Infinitif} (à la fin).
\begin{center}
\begin{tabular}{|l|l|l|l|}
\hline
\rowcolor{popblueL} \textbf{Personne} & \textbf{werden} & \textbf{Exemple : stattfinden / regieren} & \textbf{Traduction} \\
\hline
\cle{es} & \cle{wird} & \cle{Die Wahl wird 2025 stattfinden.} & \textbf{L'élection aura lieu en 2025.} \\
\cle{wir} & \cle{werden} & \cle{Wir werden wählen.} & Nous voterons. \\
\cle{sie} & \cle{werden} & \cle{Sie werden koalieren.} & Ils formeront une coalition. \\
\hline
\end{tabular}
\end{center}
\end{propriete}

\begin{attention}[Différence Présent / Futur I]
En allemand, le \textbf{Présent + indicateur de temps} (\all{morgen, 2025, nächstes Jahr}) exprime un \textbf{fait planifié, certain} : \all{Die Wahl findet 2025 statt.}
Le \textbf{Futur I} exprime une \textbf{prédiction, une hypothèse forte, une annonce solennelle} : \all{Die Wahl wird 2025 stattfinden.} (Communiqué officiel).
\emph{Consigne Bac :} Si le sujet exige le Futur I (\emph{„eine Vorhersage machen“, „offiziell ankündigen“}), l'utiliser. Sinon, le Présent est plus naturel.


\end{attention}

\begin{attention}[Pièges fréquents francophones -- Check-list de révision]
\begin{itemize}
    \item \textbf{Passif vs Actif / Devenir :} \all{Er wird Präsident} (Actif : Il \emph{devient} président) $\neq$ \all{Er wird gewählt} (Passif : Il \emph{est élu}). \cle{werden} auxiliaire passif + Partizip II ; \cle{werden} verbe plein + Nom/Adjectif (sans Partizip II).
    \item \textbf{Partizip II à la fin absolue :} \all{Der Präsident wird vom Volk \textbf{gewählt}.} (Jamais \all{*wird gewählt der Präsident}).
    \item \textbf{Accord du Partizip II :} \textbf{Pas d'accord} en allemand (invariable). \all{Die Abgeordneten werden gewählt.} (Pas *gewählte).
    \item \textbf{Connecteurs \& Ordre :} \cle{Deshalb} (V2) : \all{Deshalb \textbf{wählen} wir.} | \cle{Weil} (Verbe fin) : \all{Weil wir \textbf{wählen}...}
    \item \textbf{Majuscules systématiques :} \all{die wahl} $\rightarrow$ \textbf{\all{die Wahl}} ; \all{der präsident} $\rightarrow$ \textbf{\all{der Präsident}} ; \all{verfassungsgemäß} $\rightarrow$ \textbf{\all{verfassungsgemäß}} (adjectif, minuscule) MAIS \all{die Verfassungsmäßigkeit} (nom, majuscule).
    \item \textbf{Umlauts \& ß :} \all{Präsident} (pas d'Umlaut), \all{Wähler} (ä), \all{Stimme} (pas), \all{Verfassung} (pas), \all{größer} (ß après voyelle longue/diphtongue), \all{dass} (ss après voyelle brève). \cle{Präsidentschaftswahl} (ä).
\end{itemize}

\end{attention}

\courssub{Méthode : Commentaire structuré \& Production guidée (Activité d'intégration)}

\courssub{Objectifs de l'activité}
\noindent
À l'issue de cette activité d'intégration, l'élève sera capable de :
\begin{itemize}
    \item \textbf{Mobiliser} le lexique politique de la leçon (\cle{die Demokratie}, \cle{die Wahl}, \cle{die Partei}, \cle{der Präsident}, \cle{die Verfassung}, \cle{das Parlament}, \cle{der Wähler}, \cle{die Stimme}, \cle{der Wahlkreis}, \cle{die Koalition}, \cle{die Opposition}) dans une production écrite cohérente.
    \item \textbf{Appliquer} le \cle{style nominal} (Nominalstil) pour conférer un ton journalistique et objectif à son article (titres, compléments prépositionnels).
    \item \textbf{Structurer} son propos à l'aide de \cle{connecteurs logiques} (\all{aber}, \all{weil}, \all{deshalb}, \all{obwohl}, \all{zudem}, \all{dennoch}, \all{hingegen}) pour articuler arguments, concessions et conséquences.
    \item \textbf{Conjuguer} correctement le \cle{passif} (Vorgangspassiv : \all{werden} + Partizip II) pour décrire le processus électoral et le \cle{futur I} (\all{werden} + Infinitif) pour annoncer les échéances à venir.
    \item \textbf{Produire} un texte argumentatif court (10--12 lignes) respectant les conventions du genre journalistique (titre, chapeau, paragraphes) et le présenter à l'oral.
\end{itemize}

\courssub{Situation}
\noindent
Le lycée bilingue de \textbf{Yaoundé} (ou de \textbf{Douala}) lance son premier journal germanophone en ligne : \emph{„Der Schulbote“}. À l'approche des élections présidentielles et législatives au Cameroun, la rédaction cherche de jeunes correspondants pour couvrir l'actualité politique sous un angle comparatif. Vous êtes \textbf{Junior-Reporter} pour la rubrique \emph{„Politik \& Gesellschaft“}.

\begin{textebox}[Briefing de la rédaction -- Auftrag der Redaktion]
\textbf{Thema:} „Wahlen in unserem Land – Ein Vergleich mit Deutschland“

\all{Liebe Junior-Reporter,}

\all{für die nächste Ausgabe unseres Schuljournals „Der Schulbote“ brauchen wir einen Artikel über das politische System Kameruns und die anstehenden Wahlen. Eure Leser sind deutschsprachige Schüler in Kamerun, in Deutschland, in Österreich und in der Schweiz. Sie kennen das deutsche System gut, aber wissen wenig über die kamerunische Realität.}

\all{\textbf{Eure Mission:} Schreibt einen kurzen Presseartikel (10--12 Zeilen), der informiert, vergleicht und einen Ausblick gibt.}

\all{\textbf{Redaktionsschluss:} Nächste Woche Freitag.}
\all{\textbf{Format:} Titel, Lead-Absatz (Wer? Was? Wann? Wo?), Hauptteil, Ausblick.}
\end{textebox}

\noindent
Pour vous aider à saisir le ton et la structure attendus, voici un \textbf{exemple de lead} (chapô) d'un article du \emph{Tagesschau} (simplifié pour le niveau B1/B2) :

\begin{textebox}[Modèle de lead journalistique -- Beispiel Lead]
\textbf{Bundestagswahl 2025: Parteien starten in den Wahlkampf}

\all{In Deutschland wird der Bundestag alle vier Jahre gewählt. Die Parteien nominieren ihre Kandidaten und präsentieren ihre Programme. Millionen Wähler entscheiden über die politische Richtung des Landes. Die Wahl findet nach den Prinzipien der personalisierten Verhältniswahl statt.}
\end{textebox}

\courssub{Consignes}
\noindent
Travaillez \textbf{par groupes de 3 élèves} (Rédacteur en chef / Rédacteur / Correcteur). Suivez les étapes ci-dessous. Chaque étape fait l'objet d'une vérification par l'enseignant avant de passer à la suivante.

\begin{enumerate}
    \item \textbf{Analyse du briefing \& Planification (10 min)} \\
    Lisez le briefing (\emph{Auftrag der Redaktion}). Identifiez les \textbf{contraintes obligatoires} (liste de contrôle) :
    \begin{itemize}
        \item [$\square$] Présenter le système politique camerounais (République, présidentialisme, rôle du Parlement).
        \item [$\square$] Expliquer le mode de scrutin du Président / des députés (\textbf{au moins 2 phrases au passif} : \all{Der Präsident wird gewählt...}, \all{Die Abgeordneten werden bestimmt...}).
        \item [$\square$] Annoncer la date des prochaines élections (\textbf{1 phrase au futur I} : \all{Die nächsten Wahlen werden 2025 stattfinden.}).
        \item [$\square$] Comparer avec l'Allemagne (\textbf{au moins 2 connecteurs logiques} parmi : \all{aber}, \all{weil}, \all{deshalb}, \all{obwohl}, \all{zudem}, \all{dennoch}, \all{hingegen}).
        \item [$\square$] Intégrer \textbf{8 mots minimum} du lexique de la leçon (voir \emph{Ressources}).
        \item [$\square$] Utiliser le \textbf{style nominal} au moins 2 fois (titre, début de phrase, complément prépositionnel).
    \end{itemize}
    Répartissez les rôles : qui écrit le lead ? qui gère la comparaison ? qui vérifie la grammaire (passif/futur/nominalstil) ?

    \item \textbf{Rédaction collaborative de l'article (25 min)} \\
    Rédigez l'article \textbf{directement en allemand} sur une feuille A4 (ou document partagé). Respectez la structure :
    \begin{enumerate}
        \item \textbf{Titre} (court, informatif, style nominal : \all{Bundespräsidentenwahl in Kamerun 2025} ou \all{Präsidentschaftswahlen 2025: Kamerun im Vergleich}).
        \item \textbf{Lead} (2--3 phrases : situation actuelle, enjeu).
        \item \textbf{Corps du texte} (2 paragraphes) :
            \begin{itemize}
                \item Paragraphe A : Système camerounais + Passif (élection Président/Parlement).
                \item Paragraphe B : Comparaison Allemagne/Cameroun + Connecteurs.
            \end{itemize}
        \item \textbf{Conclusion / Ausblick} (1 phrase au Futur I + ouverture).
    \end{enumerate}
    \textbf{Astuce :} Utilisez le \emph{Nominalstil} pour les titres et les débuts de phrase (\all{Die Wahl des Präsidenten erfolgt...} au lieu de \all{Man wählt den Präsidenten...} ; \all{Aufgrund der Verfassung...} au lieu de \all{Weil die Verfassung das regelt...}).

    \item \textbf{Révision linguistique croisée (10 min)} \\
    Échangez votre brouillon avec un autre groupe. Vérifiez \textbf{uniquement} les points de grammaire ciblés :
    \begin{itemize}
        \item Passif : \all{werden} + Partizip II (accord ? non, invariable. Place du Partizip II à la fin ?).
        \item Futur I : \all{werden} + Infinitif (à la fin).
        \item Connecteurs : place du verbe (conjonction subordonnante = verbe à la fin ; conjonction de coordination/adverbe = verbe en 2e position).
        \item Majuscules sur \textbf{tous les noms} (\cle{die Wahl}, \cle{der Wähler}, \cle{die Stimme}, \cle{das Parlament}).
        \item Orthographe des Umlauts (\all{Ä}, \all{Ö}, \all{Ü}) et \all{ß} (\cle{Straße}, \cle{größer}, \cle{weiß} vs \cle{dass}, \cle{muss}).
    \end{itemize}
    Annotez les corrections au crayon vert. Rendez le brouillon.

    \item \textbf{Finalisation \& Mise en page (10 min)} \\
    Intégrez les corrections. Tapez la version finale (police Times New Roman, 12 pt, interligne 1.5). Ajoutez une \emph{Infobox} latérale avec 3 mots-clés définis en allemand simple (ex. \all{der Wahlkreis -- das Gebiet, in dem gewählt wird}).

    \item \textbf{Présentation orale \& Écoute active (15 min)} \\
    Chaque groupe désigne un \textbf{lecteur}. Celui-ci lit l'article devant la classe (débit journalistique, articulation, intonation descendante). Les autres élèves ont une \textbf{fiche d'écoute} : ils cochent les 8 mots du lexique qu'ils entendent et notent les 2 connecteurs logiques utilisés pour la comparaison. À la fin, débriefing collectif : quel article était le plus objectif ? Le plus fluide ? Le plus riche lexicalement ?
\end{enumerate}

\noindent
\textbf{Variante orale (si temps restant ou évaluation orale) :} \emph{„Ein Journalist interviewt einen Politiker“}. En binôme (A = Journaliste, B = Porte-parole du parti au pouvoir / de l'opposition). Préparer 5 questions/réponses intégrant : passif (\all{Wann wird der Präsident gewählt?}), futur (\all{Was wird Ihre Partei tun?}), connecteurs (\all{Obwohl wir die Mehrheit haben...}), vocabulaire (\all{Verfassung, Koalition, Opposition}). Jeu de rôle de 2 min devant la classe.

\courssub{Ressources à mobiliser (Rappel synthétique)}
\noindent
Les tableaux de vocabulaire, grammaire (Nominalstil, Konnektoren) et conjugaison (Passif, Futur I) ci-dessus (sections Lexique, Grammaire 1, Grammaire 2, Conjugaison) constituent la boîte à outils obligatoire pour réussir l'activité.

\courssub{Proposition de production et corrigé commenté}
\begin{exoresolu}[Article modèle : „Präsidentschaftswahlen in Kamerun 2025 – Ein Vergleich“]
\textbf{Énoncé :} Rédiger l'article final attendu (10--12 lignes) respectant toutes les contraintes.

\tcblower

\textbf{Production modèle (Version élève « Excellence »)} \\

\textbf{Präsidentschaftswahlen in Kamerun 2025: Zwischen Stabilität und Reformbedarf}

\all{In Kamerun finden im Jahr 2025 die nächsten Präsidentschaftswahlen statt. Das Land ist eine präsidiale Republik, in der der Präsident alle sieben Jahre direkt vom Volk gewählt wird. Die Abgeordneten des Parlaments werden durch ein Mehrheitswahlsystem in den Wahlkreisen bestimmt.}

\all{Obwohl Kamerun eine formale Demokratie mit Mehrparteiensystem ist, dominiert eine Partei die politische Landschaft seit Jahrzehnten. Deshalb fehlt oft eine echte Alternanz an der Macht. In Deutschland hingegen regiert meist eine Koalition, weil das Verhältniswahlsystem Kompromisse erfordert. Zudem kontrolliert der Bundestag die Regierung stärker als das kamerunische Parlament den Präsidenten.}

\all{Die kommende Wahl wird ein wichtiger Test für die Verfassungstreue sein. Internationale Beobachter werden den Prozess genau beobachten.}

\vspace{0.5cm}
\textbf{Commentaire détaillé des choix de langue (Corrigé commenté)} \\

\begin{enumerate}
    \item \textbf{Respect des contraintes quantitatives :} 11 lignes (format A4 standard), 4 paragraphes distincts (Titre, Lead, Corps A, Corps B, Ausblick).
    \item \textbf{Lexique (10 mots intégrés -- $\ge$ 8 exigés) :}
    \begin{itemize}
        \item \cle{Präsidentschaftswahlen} (composé nominal, style nominal)
        \item \cle{Republik} (système)
        \item \cle{Präsident} / \cle{gewählt} (passif)
        \item \cle{Abgeordneten} / \cle{Parlament} (institutions)
        \item \cle{Mehrheitswahlsystem} / \cle{Wahlkreise} (vocabulaire précis scrutin)
        \item \cle{Demokratie} / \cle{Mehrparteiensystem} (concepts)
        \item \cle{Partei} / \cle{Koalition} (acteurs)
        \item \cle{Bundestag} / \cle{Regierung} (comparaison ALL)
        \item \cle{Verfassungstreue} (style nominal, dérivé de \cle{Verfassung})
        \item \cle{Beobachter} / \cle{beobachten} (acteurs élection)
    \end{itemize}
    \item \textbf{Passif (2 occurrences + 1 Partizip II adjectival) :}
    \begin{itemize}
        \item \all{der Präsident ... \textbf{wird ... gewählt}} (Präsens Passiv, 3e pers. sg. -- contrainte respectée).
        \item \cle{Die Abgeordneten ... \textbf{werden ... bestimmt}} (Präsens Passiv, 3e pers. pl. -- contrainte respectée).
        \item \all{ein ... \textbf{bestimmt}es System} (Partizip II comme adjectif, style nominal implicite) -- bonus.
    \end{itemize}
    \item \textbf{Futur I (2 occurrences) :}
    \begin{itemize}
        \item \cle{Die kommende Wahl \textbf{wird ... sein}} (Futur I, prédiction formelle -- contrainte respectée).
        \item \cle{Internationale Beobachter \textbf{werden ... beobachten}} (Futur I, 3e pl. -- annonce).
    \end{itemize}
    \item \textbf{Connecteurs logiques (5 occurrences -- $\ge$ 2 exigés) :}
    \begin{itemize}
        \item \cle{Obwohl} (Subordination, concession, verbe à la fin : \all{ist}) -- Paragraphe B, phrase 1.
        \item \cle{Deshalb} (Coordination, conséquence, V2 : \all{fehlt}) -- Paragraphe B, phrase 2.
        \item \cle{hingegen} (Adverbe de liaison, contraste, V2 : \all{regiert}) -- Paragraphe B, phrase 3.
        \item \cle{weil} (Subordination, cause, verbe à la fin : \all{erfordert}) -- Paragraphe B, phrase 3.
        \item \cle{Zudem} (Coordination, addition, V2 : \all{kontrolliert}) -- Paragraphe B, phrase 4.
    \end{itemize}
    \item \textbf{Style nominal (Nominalstil) :}
    \begin{itemize}
        \item Titre : \all{Präsidentschaftswahlen ...} (Nominalisation de \all{Präsident wählen}).
        \item Lead : \all{die nächsten Präsidentschaftswahlen stattfinden} (verbal) mais structure nominale \all{In Kamerun finden ... statt}.
        \item Corps A : \all{durch ein Mehrheitswahlsystem} (Préposition + Nom au lieu de \all{weil man ein Mehrheitswahlsystem hat}).
        \item Corps B : \all{die politische Landschaft} (métaphore nominale).
        \item Conclusion : \cle{Verfassungstreue} (Nom abstrait en \cle{-e} / \cle{-keit} au lieu de \all{treu zur Verfassung sein}).
    \end{itemize}
    \item \textbf{Orthographe / Ponctuation :} Majuscules systématiques sur tous les noms (\all{Land, Jahr, Wahlen, Republik, Präsident, Volk, Abgeordnete, Parlament, Wahlkreise, Demokratie, Mehrparteiensystem, Partei, Jahrzehnte, Alternanz, Macht, Deutschland, Koalition, Verhältniswahlsystem, Kompromisse, Bundestag, Regierung, Parlament, Präsident, Wahl, Test, Verfassungstreue, Prozess, Beobachter}). Guillemets allemands non nécessaires ici (pas de citation directe). Umlauts corrects (\all{Ä} dans \all{Präsidentschaftswahlen}, \all{Ä} dans \all{Mehrheitswahlsystem}, \all{Ä} dans \all{Verhältniswahlsystem}, \all{Ö} dans \all{Koalition} -- non, \all{Koalition} sans Umlaut, \all{Ü} dans \all{Überwachung} absent ici. Mots présents : \all{Jahre, wählen, Abgeordneten, Mehrheitswahlsystem, Wahlkreise, Jahrzehnte, Verhältniswahlsystem, Beobachter} -- Umlauts sur \all{Ä} uniquement. Correct).
\end{enumerate}
\end{exoresolu}

\courssub{Bilan de l'activité}
\noindent
Cette activité de \emph{Junior-Reporter} a permis de \textbf{réinvestir de manière contextualisée} l'ensemble des savoirs de la leçon AL24 :

\begin{itemize}
    \item \textbf{Compétence lexicale :} Le passage de la liste de vocabulaire à l'emploi en contexte (collocations : \all{Präsidentschaftswahlen stattfinden}, \all{Mehrheitswahlsystem anwenden}, \cle{Verfassungstreue prüfen}) ancre la mémorisation à long terme.
    \item \textbf{Compétence grammaticale :} La contrainte explicite (\emph{« 2 phrases au passif, 1 au futur, 2 connecteurs, 2 nominalstil »}) force l'élève à \textbf{choisir consciemment} la structure syntaxique adaptée au sens (processus $\rightarrow$ passif ; prédiction officielle $\rightarrow$ futur ; argumentation $\rightarrow$ connecteurs ; densité $\rightarrow$ nominalstil). La révision par les pairs (étape 3) développe l'auto-correction et le regard métalinguistique.
    \item \textbf{Compétence scripturale :} L'écriture journalistique (style nominal, structure en pyramide inversée : Lead $\rightarrow$ Détails $\rightarrow$ Contexte $\rightarrow$ Ausblick) est une compétence transférable (Bac, vie professionnelle, études supérieures).
    \item \textbf{Compétence interculturelle :} La comparaison structurée \textbf{Cameroun -- Allemagne} (\cle{obwohl} / \cle{weil} / \cle{deshalb} / \cle{zudem} / \cle{hingegen}) oblige à dépasser le cliché pour analyser des faits institutionnels (scrutin uninominal majoritaire vs proportionnel personnalisé, présidentialisme fort vs parlementarisme chancelier, parti dominant vs culture de coalition).
    \item \textbf{Compétence orale :} La lecture à voix haute (étape 5) travaille la prosodie de l'allemand (accent tonique, intonation descendante sur les phrases déclaratives, liaison \all{r}-vocalique) et l'écoute sélective (repérage de mots-clés).
\end{itemize}

\begin{aretenir}[Ce qu'il faut retenir de l'AL24 -- Intégration]
\begin{itemize}
    \item Le \cle{Nominalstil} est la marque de l'allemand écrit formel (presse, admin, sciences) : \textbf{Verbe $\rightarrow$ Nom} (\all{wählen $\rightarrow$ die Wahl}), \textbf{Conjonction $\rightarrow$ Préposition} (\all{weil $\rightarrow$ wegen} + Génitif, \all{obwohl $\rightarrow$ trotz} + Génitif).
    \item Les \cle{Konnektoren} ne sont pas interchangeables : ils dictent la \textbf{place du verbe} (V2 ou Verbe fin). Maîtriser \cle{aber/dennoch/deshalb/zudem/hingegen} (V2) vs \cle{weil/obwohl/dass/damit} (Verbe fin) est vital pour le Bac (expression écrite / traduction thème/version).
    \item Le \cle{Passiv} (\all{werden} + Partizip II) déplace le focus sur l'\textbf{action/objet} (le Président, la loi) et efface l'agent (souvent inconnu ou inutile : \all{von wem?}). Agent : \cle{von} + Dat (personne) / \cle{durch} + Akk (institution/moyen).
    \item Le \cle{Futur I} (\all{werden} + Infinitif) exprime la \textbf{certitude objective, l'annonce officielle} ou la \textbf{prédiction}, contrairement au présent + adverbe de temps (simple fait planifié).
    \item \textbf{Contexte Cameroun :} Utiliser des réalités locales (Yaoundé, Douala, ELECAM, scrutin uninominal majoritaire à un tour, Sénat, Assemblée Nationale, bilingualisme officiel, RDPC) rend la production authentique et motivante.
\end{itemize}
\end{aretenir}

\begin{savaistu}[Le saviez-vous ? -- Cameroun \& Allemagne : Parallèles électoraux]
\begin{itemize}
    \item \textbf{Âge du vote :} 20 ans au Cameroun (depuis 2012, avant 21 ans) vs 18 ans en Allemagne (16 ans pour élections européennes/communales dans certains Länder comme Brandenburg, Bremen, Hambourg, Schleswig-Holstein).
    \item \textbf{Mandat présidentiel :} Cameroun : 7 ans, \textbf{sans limitation du nombre de mandats} (révision constitutionnelle de 2008 ayant supprimé la limite de deux mandats). Allemagne : \cle{Bundespräsident} (chef d'État protocolaire) 5 ans, renouvelable une fois ; le pouvoir exécutif réel est le \cle{Bundeskanzler} (élu par le Bundestag pour 4 ans, pas de limite de mandats).
    \item \textbf{Scrutin :} Cameroun = Uninominal majoritaire à un tour (présidentielle) + Mixte (législatives : majorité + proportionnel). Allemagne = \cle{Personalisierte Verhältniswahl} (double vote : 1. \cle{Erststimme} candidat circonscription, 2. \cle{Zweitstimme} liste régionale parti). C'est ce système qui \textbf{force les coalitions} (\cle{Koalitionen}) en Allemagne (souvent CDU/CSU-SPD, ou Ampel-Koalition SPD-Grünen-FDP, ou Kenya-Koalition CDU/CSU-SPD-Grünen).
    \item \textbf{Observateurs :} Au Cameroun, la mission d'observation de l'Union Africaine (UA) et de la Francophonie (OIF) est centrale. En Allemagne, l'OSCE (Organisation pour la sécurité et la coopération en Europe) envoie des observateurs pour les fédérales (Bundestagswahl).
    \item \textbf{Femmes en politique :} Allemagne : Quotas volontaires dans certains partis (ex. \cle{Frauenquote} 50\% chez les Verts/SPD), \cle{Bundestag} $\approx$ 35\% femmes. Cameroun : Loi sur la parité (2012) peu appliquée, Assemblée Nationale $\approx$ 30\% femmes (progression récente).
\end{itemize}
\end{savaistu}

\newpage

\coursec{Méthodes et exercices résolus}

\courssub{Fiche méthode 1 : Comprendre un texte sur les systèmes politiques}

\begin{methode}[Lire et exploiter un texte politique]
\begin{enumerate}
\item \cle{Lire d'abord le titre et la source} : ils annoncent le thème (ici : \all{Politische Systeme in der Welt}). Souligner mentalement les mots-clés attendus : \all{die Demokratie}, \all{die Diktatur}, \all{die Wahl}, \all{die Verfassung}.
\item \cle{Première lecture globale} : repérer la structure du texte (définition, exemples, comparaison, conclusion). Ne pas s'arrêter sur les mots inconnus.
\item \cle{Deuxième lecture détaillée} : souligner les connecteurs logiques (\all{denn}, \all{weil}, \all{obwohl}, \all{deshalb}, \all{jedoch}) : ils portent l'argumentation et aident à répondre aux questions « Warum...? ».
\item \cle{Exploiter le contexte} pour deviner les mots inconnus : \all{die Staatsform} = \all{Staat} (État) + \all{Form} (forme) → « forme d'État ». Les mots composés allemands se devinent presque toujours par décomposition.
\item \cle{Répondre aux questions} en reprenant la formulation de la question, avec une phrase complète, sans copier intégralement le texte : reformuler avec ses propres mots.
\item \cle{Vérifier} : place du verbe en 2\textsuperscript{e} position, majuscules aux noms, accord des articles.
\end{enumerate}
\end{methode}

\begin{exoresolu}[Compréhension de l'écrit — type Bac]
Texte : \all{In einer Demokratie wählen die Bürger ihre Regierung in freien Wahlen. Die Verfassung garantiert die Grundrechte, zum Beispiel die Meinungsfreiheit. In einer Diktatur dagegen hat eine Person oder eine Partei die ganze Macht. Oppositionelle werden oft verfolgt, und die Presse wird kontrolliert.}

\textbf{Fragen :} 1. Wer wählt in einer Demokratie die Regierung? 2. Was garantiert die Verfassung? 3. Was passiert mit Oppositionellen in einer Diktatur?
\tcblower
\textbf{Analyse :} La question 1 porte sur le sujet de la première phrase (\all{die Bürger}). La question 2 cherche le complément d'objet direct de \all{garantieren}. La question 3 exige de repérer la phrase au passif (\all{werden verfolgt}).

\textbf{Réponses rédigées :}

1. \all{In einer Demokratie wählen die Bürger ihre Regierung in freien Wahlen.} (« En démocratie, ce sont les citoyens qui élisent leur gouvernement lors d'élections libres. ») — On reprend le sujet exact du texte.

2. \all{Die Verfassung garantiert die Grundrechte, zum Beispiel die Meinungsfreiheit.} (« La constitution garantit les droits fondamentaux, par exemple la liberté d'opinion. »)

3. \all{In einer Diktatur werden Oppositionelle oft verfolgt.} (« En dictature, les opposants sont souvent persécutés. ») — Phrase au passif : \all{werden} + participe II, le verbe conjugué reste en 2\textsuperscript{e} position.

\textbf{Conclusion :} chaque réponse est une phrase complète, avec verbe en 2\textsuperscript{e} position et majuscules aux noms ; on reformule sans recopier tout le paragraphe.
\end{exoresolu}

\courssub{Fiche méthode 2 : Employer le vocabulaire des systèmes politiques}

\begin{methode}[Construire et utiliser le lexique thématique]
\begin{enumerate}
\item \cle{Apprendre chaque nom avec son article et son pluriel} : \all{die Demokratie, -n} ; \all{die Diktatur, -en} ; \all{die Wahl, -en} ; \all{die Partei, -en} ; \all{der Präsident, -en} ; \all{die Verfassung, -en}.
\item \cle{Mémoriser les couples verbe + nom} (collocations) : \all{eine Wahl gewinnen/verlieren} (gagner/perdre une élection), \all{zur Wahl gehen} (aller voter), \all{die Macht übernehmen} (prendre le pouvoir), \all{eine Regierung bilden} (former un gouvernement).
\item \cle{Exploiter les familles de mots} : \all{wählen} (voter, élire) → \all{die Wahl} (l'élection) → \all{der Wähler} (l'électeur) → \all{wählbar} (éligible) ; \all{herrschen} (régner, dominer) → \all{der Diktator} → \all{die Diktatur}.
\item \cle{Associer chaque mot à un contexte} : au Cameroun comme en Allemagne, \all{die Bürger wählen den Präsidenten}.
\item \cle{Réutiliser le lexique dans ses propres phrases} : un mot n'est acquis que s'il est employé dans une production personnelle.
\end{enumerate}
\end{methode}

\begin{exoresolu}[Vocabulaire — phrases à compléter]
Complétez avec : \all{Verfassung}, \all{Wahl}, \all{Partei}, \all{Diktatur}, \all{Bürger}.

1. \all{Die ... garantiert die Rechte der Bürger.} 2. \all{Bei der nächsten ... will ich zum ersten Mal wählen.} 3. \all{In einer ... gibt es keine freie Presse.} 4. \all{Die größte ... hat die meisten Stimmen bekommen.} 5. \all{Die ... gehen alle vier Jahre zur Wahl.}
\tcblower
\textbf{Raisonnement :} on identifie d'abord le sens attendu, puis on vérifie l'article et le cas.

1. \all{Die \textbf{Verfassung} garantiert die Rechte der Bürger.} — C'est la constitution qui garantit les droits ; nom féminin, sujet au nominatif.

2. \all{Bei der nächsten \textbf{Wahl} will ich zum ersten Mal wählen.} — « Lors des prochaines élections » ; après la préposition \all{bei} → datif : \all{bei der Wahl}.

3. \all{In einer \textbf{Diktatur} gibt es keine freie Presse.} — Absence de presse libre = dictature ; \all{in} + datif (localisation, réponse à « où ? ») : \all{in einer Diktatur}.

4. \all{Die größte \textbf{Partei} hat die meisten Stimmen bekommen.} — Le parti qui obtient le plus de voix.

5. \all{Die \textbf{Bürger} gehen alle vier Jahre zur Wahl.} — Les citoyens votent ; pluriel de \all{der Bürger} (identique au singulier).

\textbf{Conclusion :} toujours vérifier le cas exigé par la préposition ou la fonction du nom avant d'écrire l'article.
\end{exoresolu}

\courssub{Fiche méthode 3 : Traduire (thème) avec le style nominal}

\begin{propriete}[Le Nominalstil : principe]
L'allemand écrit (presse, textes politiques et administratifs) préfère souvent un \cle{nom d'action} à un verbe, là où le français utilise une proposition verbale. On obtient le schéma : \cle{préposition + article + nom d'action (+ complément du nom au génitif)}.

\begin{center}
\begin{tabularx}{\linewidth}{|X|X|X|}\hline
\rowcolor{popblueL} Français (verbe) & Deutsch (Nominalstil) & Remarque \\ \hline
lutter contre la dictature & der Kampf gegen die Diktatur & \all{gegen} + accusatif \\ \hline
après les élections & nach den Wahlen & \all{nach} + datif pluriel \\ \hline
la formation du gouvernement & die Bildung der Regierung & complément au génitif \\ \hline
la protection des droits & der Schutz der Rechte & \all{der Schutz} (nom en \all{-ung} ? non : nom masculin court) \\ \hline
le développement du pays & die Entwicklung des Landes & génitif : \all{des Landes} \\ \hline
\end{tabularx}
\end{center}
\end{propriete}

\begin{methode}[Passer du français au Nominalstil allemand]
\begin{enumerate}
\item \cle{Repérer dans la phrase française les verbes d'action} que l'allemand préfère exprimer par un nom : « lutter » → \all{der Kampf} ; « décider » → \all{die Entscheidung} ; « développer » → \all{die Entwicklung}.
\item \cle{Construire le couple préposition + nom} : « après les élections » → \all{nach den Wahlen} ; « lors de la formation du gouvernement » → \all{bei der Bildung der Regierung}.
\item \cle{Choisir le bon cas} : le nom suit le cas exigé par la préposition (\all{nach} + datif, \all{trotz} + génitif, \all{für} + accusatif).
\item \cle{Former le complément du nom au génitif} : \all{die Bildung der Regierung} (la formation du gouvernement), \all{der Schutz der Grundrechte} (la protection des droits fondamentaux).
\item \cle{Vérifier la place du verbe} : même avec un long groupe nominal en tête, le verbe conjugué reste en 2\textsuperscript{e} position.
\end{enumerate}
\end{methode}

\begin{exoresolu}[Thème — traduction avec Nominalstil]
Traduisez en allemand : « Après les élections, la formation du gouvernement a commencé. La protection des droits des citoyens est importante pour la démocratie. »
\tcblower
\textbf{Analyse :}
\begin{itemize}
\item « Après les élections » → \all{nach den Wahlen} (préposition \all{nach} + datif pluriel : \all{den Wahlen}).
\item « la formation du gouvernement » → style nominal : \all{die Bildung der Regierung} (complément du nom au génitif : \all{der Regierung}).
\item « a commencé » → passé composé : \all{hat begonnen} (verbe fort \all{beginnen}, participe II \all{begonnen}, auxiliaire \all{haben}).
\item « la protection des droits des citoyens » → \all{der Schutz der Rechte der Bürger} (double génitif).
\end{itemize}

\textbf{Traduction :}

\all{Nach den Wahlen hat die Bildung der Regierung begonnen. Der Schutz der Rechte der Bürger ist für die Demokratie wichtig.}

« Après les élections, la formation du gouvernement a commencé. La protection des droits des citoyens est importante pour la démocratie. »

\textbf{Justifications :} \all{Nach den Wahlen} occupe la 1\textsuperscript{re} position, donc inversion : le verbe \all{hat} vient en 2\textsuperscript{e} position, avant le sujet \all{die Bildung der Regierung}. \all{für} + accusatif : \all{für die Demokratie}. Majuscules à tous les noms : \all{Wahlen}, \all{Bildung}, \all{Regierung}, \all{Schutz}, \all{Rechte}, \all{Bürger}, \all{Demokratie}.
\end{exoresolu}

\courssub{Fiche méthode 4 : Utiliser les connecteurs logiques}

\begin{propriete}[Connecteurs : trois comportements à ne pas confondre]
\begin{itemize}
\item \cle{Conjonctions de coordination} (\all{und}, \all{aber}, \all{denn}, \all{oder}, \all{sondern}) : elles n'occupent pas de position dans la phrase ; le verbe reste en 2\textsuperscript{e} position. \all{Er geht zur Wahl, denn er ist Bürger.}
\item \cle{Conjonctions de subordination} (\all{weil}, \all{obwohl}, \all{wenn}, \all{dass}, \all{als}) : le verbe conjugué passe à la fin de la subordonnée. \all{Er geht zur Wahl, weil er Bürger ist.}
\item \cle{Adverbes-connecteurs} (\all{deshalb}, \all{darum}, \all{jedoch}, \all{außerdem}, \all{trotzdem}, \all{dann}) : placés en 1\textsuperscript{re} position, ils provoquent l'inversion du sujet. \all{Deshalb gehen viele Bürger zur Wahl.}
\end{itemize}
\end{propriete}

\begin{methode}[Choisir et placer le bon connecteur]
\begin{enumerate}
\item \cle{Identifier la relation logique} : cause (\all{denn}, \all{weil}), conséquence (\all{deshalb}, \all{darum}), opposition (\all{aber}, \all{jedoch}, \all{obwohl}), addition (\all{und}, \all{außerdem}), temps (\all{dann}, \all{danach}).
\item \cle{Distinguer coordination et subordination} : \all{denn} et \all{weil} traduisent tous les deux « parce que », mais \all{denn} laisse le verbe en 2\textsuperscript{e} position, \all{weil} l'envoie à la fin.
\item \cle{Placer les adverbes-connecteurs} (\all{deshalb}, \all{jedoch}, \all{außerdem}) en 1\textsuperscript{re} position → inversion du sujet : \all{Deshalb gehen viele Bürger zur Wahl.}
\item \cle{Varier les connecteurs} dans une production : ne pas répéter \all{und} ; alterner \all{deshalb}, \all{außerdem}, \all{jedoch}.
\end{enumerate}
\end{methode}

\begin{exoresolu}[Connecteurs — transformation de phrases]
Reliez les phrases avec le connecteur indiqué :

1. \all{Die Bürger gehen zur Wahl. Sie wollen die Regierung bestimmen.} (\all{weil})

2. \all{In der Diktatur gibt es keine freien Wahlen. Die Opposition wird verfolgt.} (\all{außerdem})

3. \all{Die Partei hat viele Stimmen bekommen. Sie hat die Wahl nicht gewonnen.} (\all{obwohl})
\tcblower
\textbf{Raisonnement :} pour chaque phrase, on détermine si le connecteur est une conjonction de subordination (verbe à la fin) ou un adverbe (inversion).

1. \all{Die Bürger gehen zur Wahl, weil sie die Regierung bestimmen wollen.} — \all{weil} = subordonnée : le verbe conjugué \all{wollen} passe à la fin, après l'infinitif \all{bestimmen}. (« Les citoyens vont voter parce qu'ils veulent désigner le gouvernement. »)

2. \all{In der Diktatur gibt es keine freien Wahlen. Außerdem wird die Opposition verfolgt.} — \all{außerdem} = adverbe en 1\textsuperscript{re} position : inversion, \all{wird} avant le sujet \all{die Opposition}. (« En dictature, il n'y a pas d'élections libres. De plus, l'opposition est persécutée. »)

3. \all{Obwohl die Partei viele Stimmen bekommen hat, hat sie die Wahl nicht gewonnen.} — \all{obwohl} en tête : verbe conjugué \all{hat} à la fin de la subordonnée ; dans la principale qui suit, le verbe \all{hat} occupe la 1\textsuperscript{re} position (la subordonnée entière compte comme position 1). (« Bien que le parti ait obtenu beaucoup de voix, il n'a pas gagné l'élection. »)

\textbf{Conclusion :} la règle d'or : subordination = verbe à la fin ; adverbe-connecteur = inversion du sujet.
\end{exoresolu}

\courssub{Fiche méthode 5 : Passif et futur dans un texte politique}

\begin{propriete}[Conjugaison de \all{werden} (auxiliaire du passif et du futur)]
\begin{center}
\begin{tabular}{|l|l|l|}\hline
\rowcolor{popblueL} Personne & Präsens & Präteritum \\ \hline
ich & werde & wurde \\ \hline
du & wirst & wurdest \\ \hline
er/sie/es & wird & wurde \\ \hline
wir & werden & wurden \\ \hline
ihr & werdet & wurdet \\ \hline
sie/Sie & werden & wurden \\ \hline
\end{tabular}
\end{center}

Participe II : \all{geworden} — mais au \cle{Perfekt du passif}, on utilise la forme courte \all{worden} : \all{Der Präsident ist gewählt worden.} (Le président a été élu.)
\end{propriete}

\begin{methode}[Former et employer le passif et le futur]
\begin{enumerate}
\item \cle{Passif (Vorgangspassiv)} : \all{werden} + participe II. Le sujet subit l'action : \all{Der Präsident wird gewählt.} (Le président est élu.)
\item \cle{Conjuguer} \all{werden} au temps voulu : Präsens \all{wird gewählt} ; Präteritum \all{wurde gewählt} ; Perfekt \all{ist gewählt worden} (attention : \all{worden}, pas \all{geworden}).
\item \cle{Introduire l'agent} avec \all{von} + datif : \all{Der Präsident wird vom Volk gewählt.} (\all{von dem} = \all{vom}).
\item \cle{Futur I} : \all{werden} + infinitif à la fin : \all{Die Bürger werden einen neuen Präsidenten wählen.} (Les citoyens éliront un nouveau président.)
\item \cle{Choisir le passif} quand l'agent est inconnu ou sans importance : \all{Die Verfassung wurde verabschiedet.} (La constitution a été adoptée.)
\end{enumerate}
\end{methode}

\begin{exoresolu}[Passif et futur — transformations]
1. Mettez au passif (Präsens) : \all{Das Volk wählt den Präsidenten.}

2. Mettez au passif (Präteritum) : \all{Das Parlament verabschiedete das Gesetz.}

3. Mettez au futur : \all{Die Partei bildet eine neue Regierung.}
\tcblower
\textbf{Raisonnement :} au passif, le complément d'objet direct (\all{den Präsidenten}) devient le sujet au nominatif ; l'ancien sujet devient complément d'agent avec \all{von} ou disparaît.

1. \all{Der Präsident wird vom Volk gewählt.} — \all{den Präsidenten} (accusatif) → \all{der Präsident} (nominatif) ; \all{wählen} → participe II \all{gewählt} ; \all{von dem Volk} = \all{vom Volk}. (« Le président est élu par le peuple. »)

2. \all{Das Gesetz wurde vom Parlament verabschiedet.} — Präteritum de \all{werden} : \all{wurde} ; \all{verabschieden} est un verbe faible à préfixe inséparable \all{ver-} : le participe II ne prend donc pas de \all{ge-} → \all{verabschiedet}. (« La loi fut adoptée par le parlement. »)

3. \all{Die Partei wird eine neue Regierung bilden.} — Futur I : \all{wird} en 2\textsuperscript{e} position + infinitif \all{bilden} à la fin ; \all{eine neue Regierung} reste à l'accusatif. (« Le parti formera un nouveau gouvernement. »)

\textbf{Conclusion :} ne jamais confondre \all{werden} auxiliaire du passif (« être ») et \all{werden} auxiliaire du futur : le passif exige un participe II, le futur un infinitif.
\end{exoresolu}

\begin{savaistu}[Et en Allemagne ?]
En République fédérale d'Allemagne, le chef de l'État est le \all{Bundespräsident}, élu pour cinq ans non pas directement par le peuple, mais par une assemblée spéciale, la \all{Bundesversammlung}. Le vrai chef du gouvernement est le ou la \all{Bundeskanzler(in)}, élu(e) par le \all{Bundestag}, le parlement. La constitution allemande de 1949 s'appelle le \all{Grundgesetz} (« Loi fondamentale ») : son premier article affirme que la dignité humaine est inviolable. Comparer les systèmes politiques est un sujet classique du commentaire au Bac.
\end{savaistu}

\courssub{Fiche méthode 6 : Produire un court paragraphe argumenté}

\begin{methode}[Rédiger un paragraphe sur les systèmes politiques]
\begin{enumerate}
\item \cle{Structurer en 4 à 6 phrases} : phrase d'introduction (le thème), 2-3 arguments avec connecteurs, exemple (Cameroun, Allemagne), phrase de conclusion.
\item \cle{Plan type} : \all{Zuerst} (d'abord) → définition ; \all{Dann/Außerdem} → argument ; \all{Aber/Jedoch} → nuance ; \all{Deshalb} → conclusion.
\item \cle{Employer le vocabulaire du chapitre} : au moins 4 mots du lexique (\all{die Wahl}, \all{die Verfassung}, \all{die Partei}, \all{die Bürger}...).
\item \cle{Varier les structures} : une phrase au passif, une subordonnée avec \all{weil} ou \all{obwohl}, une phrase au futur.
\item \cle{Relire} : verbe en 2\textsuperscript{e} position, majuscules, accords, Umlaute.
\end{enumerate}
\end{methode}

\begin{exoresolu}[Production écrite guidée]
Rédigez un court paragraphe (5-6 phrases) : « Warum ist die Demokratie wichtig? »
\tcblower
\textbf{Préparation :} plan = définition → droits des citoyens → rôle des élections → nuance → conclusion. Vocabulaire visé : \all{die Demokratie}, \all{die Bürger}, \all{die Wahl}, \all{die Verfassung}, \all{die Meinungsfreiheit}.

\textbf{Production modèle :}

\all{Die Demokratie ist eine Staatsform, in der die Bürger die Macht haben. Zuerst wählen sie in freien Wahlen ihre Regierung. Außerdem garantiert die Verfassung wichtige Grundrechte, zum Beispiel die Meinungsfreiheit. In einer Diktatur dagegen werden die Bürger oft kontrolliert, weil eine Person allein die Macht hat. Deshalb ist die Demokratie für ein Land wie Kamerun oder Deutschland sehr wichtig: Sie schützt die Freiheit aller Menschen.}

« La démocratie est une forme d'État dans laquelle les citoyens détiennent le pouvoir. D'abord, ils élisent leur gouvernement lors d'élections libres. De plus, la constitution garantit d'importants droits fondamentaux, par exemple la liberté d'opinion. Dans une dictature, en revanche, les citoyens sont souvent contrôlés, parce qu'une seule personne détient le pouvoir. C'est pourquoi la démocratie est très importante pour un pays comme le Cameroun ou l'Allemagne : elle protège la liberté de tous les hommes. »

\textbf{Points de grammaire vérifiés :} relative avec \all{in der} (datif féminin, localisation) ; passif \all{werden kontrolliert} ; subordonnée \all{weil eine Person allein die Macht hat} (verbe à la fin) ; connecteurs variés \all{zuerst}, \all{außerdem}, \all{dagegen}, \all{deshalb} avec inversion correcte ; \all{Kamerun} est neutre : \all{ein Land wie Kamerun}.
\end{exoresolu}

\begin{attention}[Les 5 erreurs qui coûtent des points au Bac]
\begin{enumerate}
\item \cle{Oublier la majuscule aux noms} : écrire \all{die wahl} au lieu de \all{die Wahl}. En allemand, TOUS les noms prennent la majuscule : \all{die Demokratie}, \all{der Präsident}, \all{die Verfassung}.
\item \cle{Mal placer le verbe dans la subordonnée} : \all{... weil er ist Bürger} est FAUX → \all{... weil er Bürger ist.} Après \all{weil}, \all{obwohl}, \all{dass}, \all{wenn}, le verbe conjugué va à la fin.
\item \cle{Confondre passif et futur} : \all{Der Präsident wird wählen} = « le président élira » (futur) ; \all{Der Präsident wird gewählt} = « le président est élu » (passif). Participe II ≠ infinitif.
\item \cle{Calquer le français} : « les citoyens vont aux élections » ne se traduit pas mot à mot ; on dit \all{die Bürger gehen zur Wahl} (\all{zu der} = \all{zur}). De même, « prendre le pouvoir » = \all{die Macht übernehmen}, pas un calque avec \all{nehmen} seul.
\item \cle{Se tromper de cas après les prépositions} : \all{nach die Wahl} est FAUX → \all{nach der Wahl} (datif). Réviser : \all{nach}, \all{von}, \all{bei}, \all{mit} + datif ; \all{für}, \all{gegen}, \all{ohne} + accusatif.
\end{enumerate}
\end{attention}

\newpage

\coursec{Exercices}

\courssub{Je vérifie mes connaissances}

\begin{exercice}[QCM : Systèmes politiques et vocabulaire]
\emph{Choisissez la bonne réponse. Une seule proposition est correcte.}
\begin{enumerate}
    \item Quel terme désigne la « constitution » en allemand ?
    \begin{itemize}
        \item[A)] das Gesetz
        \item[B)] die Verfassung
        \item[C)] der Staat
        \item[D)] die Partei
    \end{itemize}
    \item Dans une démocratie, le pouvoir émane du peuple. Comment dit-on « le peuple » ?
    \begin{itemize}
        \item[A)] der Bürger
        \item[B)] der Wähler
        \item[C)] das Volk
        \item[D)] die Nation
    \end{itemize}
    \item Le verbe \all{wählen} (élire/choisir) est un verbe fort. Quelle est sa forme au \emph{Präteritum} pour \all{er} ?
    \begin{itemize}
        \item[A)] wählte
        \item[B)] wählt
        \item[C)] gewählt
        \item[D)] wählte
    \end{itemize}
    \item Quel connecteur logique exprime une \emph{conséquence} (donc) et place le verbe en deuxième position ?
    \begin{itemize}
        \item[A)] weil
        \item[B)] obwohl
        \item[C)] deshalb
        \item[D)] dass
    \end{itemize}
    \item La phrase \all{Der Präsident wird vom Volk gewählt.} est au :
    \begin{itemize}
        \item[A)] Passif Präteritum
        \item[B)] Passif Präsens
        \item[C)] Actif Präsens
        \item[D)] Futur I
    \end{itemize}
\end{enumerate}
\end{exercice}

\begin{exercice}[Vrai ou Faux ? Justifiez en allemand]
\emph{Lisez les affirmations suivantes. Indiquez si elles sont \textbf{richtig} (vraies) ou \textbf{falsch} (fausses) et justifiez votre réponse par une phrase complète en allemand.}
\begin{enumerate}
    \item In einer Diktatur gibt es freie und geheime Wahlen.
    \item Die Verfassung ist das Grundgesetz eines Staates.
    \item Eine Partei ist eine Gruppe von Menschen, die keine politischen Ziele hat.
    \item Der Nominalstil verwendet viele Verben und wenige Substantive.
    \item „Trotzdem“ leitet einen Nebensatz ein und schiebt das Verb ans Ende.
\end{enumerate}
\end{exercice}

\begin{exercice}[Vocabulaire : Articles, pluriels et familles de mots]
\emph{Complétez le tableau ci-dessous pour les 12 mots du thème. Puis complétez les phrases à trous avec le mot approprié (forme correcte).}

\begin{center}
\begin{tabularx}{\linewidth}{|X|X|X|X|}
\hline
\rowcolor{popblueL} \textbf{Mot (Singulier)} & \textbf{Article} & \textbf{Pluriel} & \textbf{Traduction} \\ \hline
Wahl & & & \\ \hline
Verfassung & & & \\ \hline
Partei & & & \\ \hline
Präsident & & & \\ \hline
Diktatur & & & \\ \hline
Demokratie & & & \\ \hline
Regierung & & & \\ \hline
Opposition & & & \\ \hline
Kanzler & & & \\ \hline
Bürger & & & \\ \hline
Recht & & & \\ \hline
Stimme & & & \\ \hline
\end{tabularx}
\end{center}

\medskip
\emph{Phrases à trous :}
\begin{enumerate}
    \item Die \_\_\_\_\_\_\_\_\_\_\_\_ des Präsidenten findet alle fünf Jahre statt.
    \item Die \_\_\_\_\_\_\_\_\_\_\_\_ garantiert die Grundrechte der Bürger.
    \item Er ist Mitglied einer großen politischen \_\_\_\_\_\_\_\_\_\_\_\_.
    \item In einer \_\_\_\_\_\_\_\_\_\_\_\_ gibt es keine Meinungsfreiheit.
    \item Jeder \_\_\_\_\_\_\_\_\_\_\_\_ hat das Recht zu wählen.
\end{enumerate}
\end{exercice}

\begin{exercice}[Conjugaison : Passif et Futur I]
\emph{Mettez les verbes entre parenthèses à la forme demandée (Passif Präsens, Passif Präteritum ou Futur I). Attention à l'accord du sujet.}
\begin{enumerate}
    \item (Passiv Präsens) Der neue Kanzler \_\_\_\_\_\_\_\_\_\_\_\_ (wählen) vom Parlament.
    \item (Passiv Präteritum) Das Gesetz \_\_\_\_\_\_\_\_\_\_\_\_ (beschließen) gestern vom Bundestag.
    \item (Futur I) Die Wähler \_\_\_\_\_\_\_\_\_\_\_\_ (entscheiden) über die Zukunft des Landes.
    \item (Passiv Präsens) Die Stimmen \_\_\_\_\_\_\_\_\_\_\_\_ (zählen) gerade.
    \item (Passiv Präteritum) Die Regierung \_\_\_\_\_\_\_\_\_\_\_\_ (bilden) nach der Wahl.
    \item (Futur I) Wir \_\_\_\_\_\_\_\_\_\_\_\_ (diskutieren) das Ergebnis morgen in der Schule.
    \item (Passiv Präsens) Die Demokratie \_\_\_\_\_\_\_\_\_\_\_\_ (schützen) durch die Verfassung.
    \item (Futur I) Die Opposition \_\_\_\_\_\_\_\_\_\_\_\_ (kritisieren) den Haushaltsplan.
    \item (Passiv Präteritum) Der Vertrag \_\_\_\_\_\_\_\_\_\_\_\_ (unterschreiben) vom Außenminister.
    \item (Futur I) Ich \_\_\_\_\_\_\_\_\_\_\_\_ (informieren) mich über die Programme der Parteien.
\end{enumerate}
\end{exercice}

\courssub{Je m'entraîne}

\begin{exercice}[Traduction : Thème (Français $\rightarrow$ Allemand)]
\emph{Traduisez les phrases suivantes en allemand. Respectez l'ordre des mots, la casse des noms et la conjugaison (passif, futur, connecteurs).}
\begin{enumerate}
    \item L'Allemagne est une démocratie parlementaire. (Präsens Aktiv)
    \item Le chancelier est élu par le Bundestag. (Passiv Präsens)
    \item Les citoyens votent librement et secrètement. (Präsens Aktiv, adv. de manière)
    \item Bien qu'il y ait des partis nombreux, le gouvernement est stable. (\all{obwohl} + Nebensatz)
    \item Demain, les résultats seront annoncés à 18 heures. (Futur I Passiv : \all{werden bekannt gegeben})
    \item Parce que la constitution protège les droits de l'homme, il n'y a pas de dictature. (\all{weil} + Nebensatz)
    \item Le président a peu de pouvoir politique, mais il représente le pays. (\all{aber} / \all{denn})
    \item Nous allons discuter des élections en classe. (Futur I)
\end{enumerate}
\end{exercice}

\begin{exercice}[Transformation : Style nominal $\leftrightarrow$ Style verbal]
\emph{Transformez les groupes nominaux (Nominalstil) en phrases verbales complètes (Verbalstil) avec un sujet personnel (man, das Volk, der Bundestag, etc.). Conjuguez le verbe au temps indiqué.}
\begin{enumerate}
    \item (Präsens) die Wahl des Präsidenten durch das Volk $\rightarrow$ \_\_\_\_\_\_\_\_\_\_\_\_
    \item (Präteritum) die Bildung der Regierung nach der Wahl $\rightarrow$ \_\_\_\_\_\_\_\_\_\_\_\_
    \item (Perfekt) die Verabschiedung des Gesetzes im Parlament $\rightarrow$ \_\_\_\_\_\_\_\_\_\_\_\_
    \item (Präsens) die Garantie der Menschenrechte in der Verfassung $\rightarrow$ \_\_\_\_\_\_\_\_\_\_\_\_
    \item (Futur I) die Kritik der Opposition an der Politik $\rightarrow$ \_\_\_\_\_\_\_\_\_\_\_\_
\end{enumerate}
\medskip
\emph{Transformez maintenant les phrases verbales en groupes nominaux (style nominal / Nominalstil).}
\begin{enumerate}
    \item Das Volk wählt den Präsidenten. $\rightarrow$ \_\_\_\_\_\_\_\_\_\_\_\_
    \item Der Bundestag beschließt das Gesetz. $\rightarrow$ \_\_\_\_\_\_\_\_\_\_\_\_
    \item Die Verfassung schützt die Freiheit. $\rightarrow$ \_\_\_\_\_\_\_\_\_\_\_\_
    \item Die Partei nominiert den Kandidaten. $\rightarrow$ \_\_\_\_\_\_\_\_\_\_\_\_
    \item Die Bürger demonstrieren für ihre Rechte. $\rightarrow$ \_\_\_\_\_\_\_\_\_\_\_\_
\end{enumerate}
\end{exercice}

\begin{exercice}[Connecteurs logiques : Reliez les phrases]
\emph{Reliez les deux phrases en une seule phrase en utilisant le connecteur imposé entre parenthèses. Attention à la structure de la phrase (ordre des mots, casse).}
\begin{enumerate}
    \item Die Demokratie ist wichtig. Sie schützt die Freiheit. (\all{denn})
    \item Es gibt keine freien Wahlen. Das Land ist eine Diktatur. (\all{deshalb})
    \item Der Bürger hat Rechte. Er hat auch Pflichten. (\all{aber})
    \item Die Wahl war fair. Es gab keine Manipulation. (\all{weil})
    \item Die Opposition kritisiert die Regierung. Sie wird nicht verboten. (\all{obwohl})
    \item Man muss wählen gehen. Man beeinflusst die Politik. (\all{damit} + Konjunktiv II ou infinitif avec \all{um ... zu})
    \item Der Kanzler regiert. Er braucht die Mehrheit im Parlament. (\all{weil})
    \item Die Verfassung gilt für alle. Niemand steht über dem Gesetz. (\all{sodass})
\end{enumerate}
\end{exercice}

\begin{exercice}[Texte à trous : Les élections en Allemagne]
\emph{Complétez le texte suivant avec les mots manquants (articles, prépositions, connecteurs, formes verbales passif/futur, vocabulaire thématique).}

\begin{textebox}[Wahlen in Deutschland]
In Deutschland \_\_\_\_\_\_\_\_\_\_\_\_ (finden / Futur I) die Bundestagswahl alle vier Jahre statt. Das System \_\_\_\_\_\_\_\_\_\_\_\_ (heißen) personalisierte Verhältniswahl. Jeder Wähler \_\_\_\_\_\_\_\_\_\_\_\_ (haben) zwei Stimmen: die Erststimme für einen Direktkandidaten \_\_\_\_\_\_\_\_\_\_\_\_ (in) seinem Wahlkreis und die Zweitstimme für eine Partei.

\_\_\_\_\_\_\_\_\_\_\_\_ (Weil / Denn) die Zweitstimme entscheidend für die Sitzverteilung ist, \_\_\_\_\_\_\_\_\_\_\_\_ (werden / Passiv Präsens) die Sitze im Bundestag proportional verteilt. Eine Partei muss \_\_\_\_\_\_\_\_\_\_\_\_ (mindestens) 5\% der Stimmen erreichen \_\_\_\_\_\_\_\_\_\_\_\_ (oder) drei Direktmandate gewinnen, \_\_\_\_\_\_\_\_\_\_\_\_ (damit) sie in den Bundestag einzieht.

Nach der Wahl \_\_\_\_\_\_\_\_\_\_\_\_ (werden / Futur I Passiv) der Bundeskanzler \_\_\_\_\_\_\_\_\_\_\_\_ (vom Bundestag / wählen). \_\_\_\_\_\_\_\_\_\_\_\_ (Obwohl) der Bundespräsident das Staatsoberhaupt ist, \_\_\_\_\_\_\_\_\_\_\_\_ (haben / er) die politische Macht der Kanzler. Die Verfassung, das Grundgesetz, \_\_\_\_\_\_\_\_\_\_\_\_ (schützen / Passiv Präsens) die Demokratie vor Extremismus.
\end{textebox}
\end{exercice}

\courssub{Je me prépare au Bac}

\begin{exercice}[Compréhension écrite type Bac : « Die Bundestagswahl 2025 »]
\emph{Lisez le texte suivant attentivement, puis répondez aux questions en allemand (1 à 5) et aux items Vrai/Faux (6 à 8). Pour les items Vrai/Faux, citez la phrase du texte qui justifie votre réponse.}

\begin{textebox}[Die Bundestagswahl 2025: Ein Blick auf das politische System]
Am 23. Februar 2025 finden in Deutschland vorgezogene Neuwahlen zum Bundestag statt. Der Bundeskanzler Olaf Scholz hatte die Vertrauensfrage gestellt und verloren, woraufhin der Bundespräsident den Bundestag auflöste. Dies zeigt die Funktionsfähigkeit der parlamentarischen Demokratie: Wenn eine Regierung keine Mehrheit mehr hat, kann das Parlament neu gewählt werden.

Das Wahlsystem bleibt die personalisierte Verhältniswahl. Mit der Erststimme wählen die Bürger einen Direktkandidaten in ihrem Wahlkreis, mit der Zweitstimme eine Landesliste einer Partei. Nur Parteien, die mindestens fünf Prozent der Zweitstimmen erreichen (Fünf-Prozent-Hürde) oder drei Direktmandate gewinnen, ziehen in den Bundestag ein. Diese Hürde soll die Zersplitterung des Parlaments verhindern und stabile Mehrheiten ermöglichen.

Wahlberechtigt sind alle deutschen Staatsbürger ab 18 Jahren, die seit mindestens drei Monaten in Deutschland wohnen. Die Wahl ist allgemein, unmittelbar, frei, gleich und geheim – die fünf Grundsätze des Wahlrechts nach Artikel 38 des Grundgesetzes. Briefwahl ist für alle ohne Angabe von Gründen möglich, was die Wahlbeteiligung oft erhöht.

Nach der Wahl bildet sich der neue Bundestag. Der Bundespräsident schlägt einen Kandidaten für das Kanzleramt vor, der dann vom Bundestag ohne Aussprache geheim gewählt wird. Erreicht kein Kandidat die absolute Mehrheit (Kanzlermehrheit), gibt es weitere Wahlgänge. Gelingt auch dann keine Wahl, kann der Bundespräsident einen Minderheitskanzler ernennen oder den Bundestag auflösen.

Die neue Regierung steht vor großen Herausforderungen: Wirtschaftskrise, Klimaschutz, Migration und die Unterstützung der Ukraine. Die Wähler erwarten Lösungen. Die Demokratie lebt vom Wettbewerb der Ideen und vom Respekt vor dem politischen Gegner – im Gegensatz zur Diktatur, wo Andersdenkende unterdrückt werden.
\end{textebox}

\medskip
\textbf{Fragen de compréhension (répondez en allemand, phrases complètes) :}
\begin{enumerate}
    \item Warum finden im Februar 2025 vorgezogene Neuwahlen statt?
    \item Erklären Sie den Unterschied zwischen Erststimme und Zweitstimme.
    \item Welche Funktion hat die Fünf-Prozent-Hürde?
    \item Nennen Sie die fünf Wahlgrundsätze laut Artikel 38 Grundgesetz.
    \item Was passiert, wenn im Bundestag kein Kandidat die absolute Mehrheit für das Kanzleramt erreicht?
\end{enumerate}

\textbf{Vrai ou Faux ? Justifiez par une citation du texte :}
\begin{enumerate}
    \setcounter{enumi}{5}
    \item Der Bundeskanzler wird direkt vom Volk gewählt.
    \item Die Briefwahl ist nur mit triftigem Grund erlaubt.
    \item In einer Diktatur werden Andersdenkende respektiert.
\end{enumerate}
\end{exercice}

\begin{exercice}[Thème et Version]
\emph{\textbf{Partie A (Thème) :} Traduisez le paragraphe suivant en allemand. Utilisez le style nominal, le passif, le futur et au moins trois connecteurs logiques (z. B. \all{weil}, \all{deshalb}, \all{obwohl}, \all{denn}).}

« La démocratie est un système politique dans lequel le pouvoir vient du peuple. Parce que les citoyens élisent leurs représentants, le parlement contrôle le gouvernement. Bien que les décisions soient souvent lentes, la liberté est protégée par la constitution. C'est pourquoi je préfère la démocratie à la dictature. Demain, nous voterons pour l'avenir de notre pays. »

\medskip
\emph{\textbf{Partie B (Version) :} Traduisez le texte suivant en français. Attention au style nominal (\all{die Wahl des Kanzlers...}) et aux formes passives.}

\begin{textebox}[Version]
Die Wahl des Bundeskanzlers durch den Bundestag ist ein zentraler Akt der parlamentarischen Demokratie. Der Kandidat wird ohne Aussprache geheim gewählt. Erreicht er die notwendige Mehrheit, wird er vom Bundespräsidenten ernannt. Wird keine Mehrheit erzielt, kann der Präsident den Bundestag auflösen. Die Verfassung sichert diesem Verfahren den Vorrang vor instabilen Koalitionen zu.
\end{textebox}
\end{exercice}

\begin{exercice}[Production écrite type Bac : Sujet de rédaction]
\emph{Traitez le sujet suivant en allemand. Longueur attendue : \textbf{8 à 10 lignes} (environ 80--100 mots). Votre production doit comporter :}
\begin{itemize}
    \item \textbf{Au moins 3 connecteurs logiques} différents (ex. \all{weil}, \all{deshalb}, \all{obwohl}, \all{denn}, \all{trotzdem}, \all{sodass}).
    \item \textbf{Au moins 2 formes de passif} (Präsens ou Präteritum).
    \item \textbf{Au moins 1 forme de Futur I}.
    \item \textbf{Du vocabulaire thématique} (Demokratie, Diktatur, Wahl, Verfassung, Partei, Präsident, Recht, Freiheit, Regierung, Opposition).
    \item \textbf{Un avis personnel argumenté} (Ich finde, Meiner Meinung nach, Ich bin der Ansicht, dass...).
\end{itemize}

\medskip
\textbf{Sujet :}
„Vergleichen Sie die Demokratie und die Diktatur. Welches System bevorzugen Sie und warum? Nennen Sie konkrete Unterschiede bei Wahlen, Rechten der Bürger und Machtverteilung.“
\end{exercice}

\newpage

\coursec{Corrigés des exercices}

\coursec{Corrigés des exercices — Leçon AL24 : Les systèmes politiques}

\begin{corrige}[Exercice 1 — QCM : Systèmes politiques et vocabulaire]
\begin{enumerate}
    \item \textbf{B) die Verfassung}. \all{das Gesetz} = la loi ; \all{der Staat} = l'État ; \all{die Partei} = le parti. La constitution = \all{die Verfassung} (en Allemagne : \all{das Grundgesetz}).
    \item \textbf{C) das Volk}. \all{der Bürger} = le citoyen ; \all{der Wähler} = l'électeur ; \all{die Nation} = la nation. « Le pouvoir émane du peuple » = \all{Alle Staatsgewalt geht vom Volke aus.} (article 20 du \all{Grundgesetz}).
    \item \textbf{A) wählte}. \all{wählen} est un verbe faible (régulier) : \all{er wählte} au Präteritum, \all{er hat gewählt} au Perfekt. Attention : l'énoncé le présentait comme fort, c'est un piège — il est faible.
    \item \textbf{C) deshalb}. \all{deshalb} est un adverbe conjonctif de conséquence qui occupe la première position et provoque l'inversion (verbe en 2\textsuperscript{e} position) : \all{Es gibt keine Wahl, deshalb ist es eine Diktatur.} \all{weil}, \all{obwohl}, \all{dass} introduisent des subordonnées (verbe en fin).
    \item \textbf{B) Passif Präsens}. Structure : \all{wird} (auxiliaire \all{werden} au Präsens) + participe II \all{gewählt} + agent introduit par \all{von} (\all{vom Volk}). Le Präteritum serait \all{wurde gewählt}.
\end{enumerate}
\end{corrige}

\begin{corrige}[Exercice 2 — Vrai ou Faux ?]
\begin{enumerate}
    \item \textbf{Falsch.} \all{In einer Diktatur gibt es keine freien und geheimen Wahlen, weil die Opposition unterdrückt wird.} (Dans une dictature, il n'y a pas d'élections libres et secrètes.)
    \item \textbf{Richtig.} \all{Die Verfassung ist das Grundgesetz eines Staates und garantiert die Rechte der Bürger.}
    \item \textbf{Falsch.} \all{Eine Partei ist eine Gruppe von Menschen, die gemeinsame politische Ziele hat.}
    \item \textbf{Falsch.} (L'énoncé affirmait le contraire.) Richtig ist: \all{Der Nominalstil verwendet viele Substantive und wenige Verben.} (Le style nominal utilise beaucoup de noms et peu de verbes ; c'est sa définition.)
    \item \textbf{Falsch.} \all{„Trotzdem“ ist ein Adverb und steht am Anfang des Hauptsatzes; das Verb bleibt an zweiter Stelle.} \all{trotzdem} provoque l'inversion sujet-verbe, il ne renvoie pas le verbe à la fin.
\end{enumerate}
\end{corrige}

\begin{corrige}[Exercice 3 — Vocabulaire : articles, pluriels et familles de mots]
\begin{center}
\begin{tabularx}{\linewidth}{|X|X|X|X|}
\hline
\rowcolor{popblueL} \textbf{Mot} & \textbf{Article} & \textbf{Pluriel} & \textbf{Traduction} \\ \hline
Wahl & die & die Wahlen & l'élection, le choix \\ \hline
Verfassung & die & die Verfassungen & la constitution \\ \hline
Partei & die & die Parteien & le parti \\ \hline
Präsident & der & die Präsidenten & le président \\ \hline
Diktatur & die & die Diktaturen & la dictature \\ \hline
Demokratie & die & die Demokratien & la démocratie \\ \hline
Regierung & die & die Regierungen & le gouvernement \\ \hline
Opposition & die & die Oppositionen & l'opposition \\ \hline
Kanzler & der & die Kanzler & le chancelier \\ \hline
Bürger & der & die Bürger & le citoyen \\ \hline
Recht & das & die Rechte & le droit \\ \hline
Stimme & die & die Stimmen & la voix, la voix électorale \\ \hline
\end{tabularx}
\end{center}

\medskip
\textbf{Phrases à trous :}
\begin{enumerate}
    \item \all{Die \textbf{Wahl} des Präsidenten findet alle fünf Jahre statt.}
    \item \all{Die \textbf{Verfassung} garantiert die Grundrechte der Bürger.}
    \item \all{Er ist Mitglied einer großen politischen \textbf{Partei}.}
    \item \all{In einer \textbf{Diktatur} gibt es keine Meinungsfreiheit.}
    \item \all{Jeder \textbf{Bürger} hat das Recht zu wählen.}
\end{enumerate}
\textbf{Points de vigilance :} \all{der Kanzler} et \all{der Bürger} sont identiques au pluriel (déclinaison faible en \all{-n} seulement aux cas obliques : \all{den Kanzler}). \all{das Recht} fait son pluriel en \all{-e} : \all{die Rechte}.
\end{corrige}

\begin{corrige}[Exercice 4 — Conjugaison : Passif et Futur I]
\begin{enumerate}
    \item \all{Der neue Kanzler \textbf{wird gewählt} vom Parlament.} (Passiv Präsens : \all{wird} + Partizip II)
    \item \all{Das Gesetz \textbf{wurde beschlossen} gestern vom Bundestag.} (Passiv Präteritum : \all{wurde} + Partizip II en \all{-en}, verbe fort/irrégulier \all{beschließen})
    \item \all{Die Wähler \textbf{werden entscheiden} über die Zukunft des Landes.} (Futur I : \all{werden} + infinitif en fin)
    \item \all{Die Stimmen \textbf{werden gerade gezählt}.} (ordre standard : auxiliaire + adverbe + participe II)
    \item \all{Die Regierung \textbf{wurde gebildet} nach der Wahl.} (participe de \all{bilden} : \all{gebildet})
    \item \all{Wir \textbf{werden diskutieren} das Ergebnis morgen in der Schule.}
    \item \all{Die Demokratie \textbf{wird geschützt} durch die Verfassung.} (participe de \all{schützen} : \all{geschützt})
    \item \all{Die Opposition \textbf{wird kritisieren} den Haushaltsplan.}
    \item \all{Der Vertrag \textbf{wurde unterschrieben} vom Außenminister.} (participe II en \all{-en}, verbe fort/irrégulier \all{unterschreiben})
    \item \all{Ich \textbf{werde mich informieren} über die Programme der Parteien.} (verbe réfléchi \all{sich informieren} : le pronom réfléchi \all{mich} suit le verbe conjugué \all{werde}).
\end{enumerate}
\textbf{Points de vigilance :} ne pas confondre l'auxiliaire du passif (\all{werden/wurden} + Partizip II) avec celui du futur (\all{werden} + Infinitiv). Au passif, l'agent est introduit par \all{von} + datif (\all{vom Bundestag}).
\end{corrige}

\begin{corrige}[Exercice 5 — Traduction : Thème]
\begin{enumerate}
    \item \all{Deutschland ist eine parlamentarische Demokratie.}
    \item \all{Der Kanzler wird vom Bundestag gewählt.} (« par le Bundestag » = \all{vom Bundestag}, agent au passif.)
    \item \all{Die Bürger wählen frei und geheim.} (adverbes de manière invariables, placés après le verbe.)
    \item \all{Obwohl es viele Parteien gibt, ist die Regierung stabil.} (subordonnée en tête : verbe conjugué en fin de subordonnée, puis inversion dans la principale.)
    \item \all{Morgen werden die Ergebnisse um 18 Uhr bekannt gegeben werden.} (Futur I passif : \all{werden} + Partizip II \all{bekannt gegeben} + Infinitif \all{werden} — double infinitif en fin de phrase. On utilise souvent le Présens passif avec valeur future : \all{Morgen werden die Ergebnisse um 18 Uhr bekannt gegeben.})
    \item \all{Weil die Verfassung die Menschenrechte schützt, gibt es keine Diktatur.}
    \item \all{Der Präsident hat wenig politische Macht, aber er vertritt das Land.} (\all{aber} : coordination, verbe en 2\textsuperscript{e} position.)
    \item \all{Wir werden in der Klasse über die Wahlen diskutieren.} (\all{über} + accusatif avec \all{diskutieren}.)
\end{enumerate}
\textbf{Barème indicatif :} 2 points par phrase (1 pt vocabulaire/cas, 1 pt morphosyntaxe : ordre des mots, conjugaison).
\end{corrige}

\begin{corrige}[Exercice 6 — Transformation : Nominalstil $\leftrightarrow$ Verbalstil]
\textbf{A. Du style nominal au style verbal :}
\begin{enumerate}
    \item \all{Das Volk wählt den Präsidenten.}
    \item \all{Man bildete die Regierung nach der Wahl.} ou \all{Die Regierung wurde nach der Wahl gebildet.}
    \item \all{Das Parlament hat das Gesetz verabschiedet.}
    \item \all{Die Verfassung garantiert die Menschenrechte.}
    \item \all{Die Opposition wird die Politik kritisieren.}
\end{enumerate}
\textbf{B. Du style verbal au style nominal :}
\begin{enumerate}
    \item \all{die Wahl des Präsidenten durch das Volk}
    \item \all{der Beschluss des Gesetzes durch den Bundestag}
    \item \all{der Schutz der Freiheit durch die Verfassung}
    \item \all{die Nominierung des Kandidaten durch die Partei}
    \item \all{die Demonstration der Bürger für ihre Rechte}
\end{enumerate}
\textbf{Points de vigilance :} au Nominalstil, le complément d'objet devient un génitif (\all{des Präsidenten}, \all{des Gesetzes}) et l'agent est introduit par \all{durch} + accusatif. Les noms d'action prennent une majuscule et un article.
\end{corrige}

\begin{corrige}[Exercice 7 — Connecteurs logiques]
\begin{enumerate}
    \item \all{Die Demokratie ist wichtig, denn sie schützt die Freiheit.} (\all{denn} : coordination, pas d'inversion.)
    \item \all{Es gibt keine freien Wahlen, deshalb ist das Land eine Diktatur.} (\all{deshalb} : inversion sujet-verbe.)
    \item \all{Der Bürger hat Rechte, aber er hat auch Pflichten.}
    \item \all{Die Wahl war fair, weil es keine Manipulation gab.} (\all{weil} : verbe conjugué \all{gab} en fin.)
    \item \all{Obwohl die Opposition die Regierung kritisiert, wird sie nicht verboten.}
    \item \all{Man muss wählen gehen, um die Politik zu beeinflussen.} ou \all{Man muss wählen gehen, damit man die Politik beeinflussen kann.}
    \item \all{Der Kanzler regiert, weil er die Mehrheit im Parlament hat.} (formulation correcte : indicatif présent, \all{weil} + verbe en fin.)
    \item \all{Die Verfassung gilt für alle, sodass niemand über dem Gesetz steht.} (\all{sodass} : conséquence, verbe en fin.)
\end{enumerate}
\textbf{Points de vigilance :} distinguer les coordinateurs (\all{denn}, \all{aber} : verbe en 2\textsuperscript{e} position), les adverbes conjonctifs (\all{deshalb}, \all{trotzdem} : inversion) et les subordonnants (\all{weil}, \all{obwohl}, \all{sodass}, \all{damit} : verbe en fin).
\end{corrige}

\begin{corrige}[Exercice 8 — Texte à trous : Les élections en Allemagne]
\begin{textebox}[Wahlen in Deutschland — Texte complété]
In Deutschland \textbf{findet} die Bundestagswahl alle vier Jahre statt. Das System \textbf{heißt} personalisierte Verhältniswahl. Jeder Wähler \textbf{hat} zwei Stimmen: die Erststimme für einen Direktkandidaten \textbf{in} seinem Wahlkreis und die Zweitstimme für eine Partei.

\textbf{Weil} die Zweitstimme entscheidend für die Sitzverteilung ist, \textbf{werden} die Sitze im Bundestag proportional \textbf{verteilt}. Eine Partei muss \textbf{mindestens} 5\% der Stimmen erreichen \textbf{oder} drei Direktmandate gewinnen, \textbf{damit} sie in den Bundestag einzieht.

Nach der Wahl \textbf{wird} der Bundeskanzler \textbf{vom Bundestag gewählt}. \textbf{Obwohl} der Bundespräsident das Staatsoberhaupt ist, \textbf{hat} die politische Macht der Kanzler. Die Verfassung, das Grundgesetz, \textbf{schützt} die Demokratie vor Extremismus.
\end{textebox}
\textbf{Remarques :} \all{findet ... statt} : verbe à particule séparable, \all{statt} en fin. La subordonnée \all{weil...} en tête impose l'inversion : \all{werden die Sitze ... verteilt}. Après \all{obwohl} en tête, même inversion : \all{hat die politische Macht der Kanzler} (le sujet réel est \all{der Kanzler}).
\end{corrige}

\begin{corrige}[Exercice 9 — Compréhension écrite type Bac]
\textbf{Réponses aux questions (propositions) :}
\begin{enumerate}
    \item \all{Weil Bundeskanzler Olaf Scholz die Vertrauensfrage verloren hat und der Bundespräsident daraufhin den Bundestag aufgelöst hat.}
    \item \all{Mit der Erststimme wählt man einen Direktkandidaten im Wahlkreis, mit der Zweitstimme eine Landesliste einer Partei.}
    \item \all{Die Fünf-Prozent-Hürde soll die Zersplitterung des Parlaments verhindern und stabile Mehrheiten ermöglichen.}
    \item \all{Die Wahl ist allgemein, unmittelbar, frei, gleich und geheim.}
    \item \all{Dann gibt es weitere Wahlgänge; wenn auch dann keine Wahl gelingt, kann der Bundespräsident einen Minderheitskanzler ernennen oder den Bundestag auflösen.}
\end{enumerate}
\textbf{Vrai ou Faux :}
\begin{enumerate}
    \setcounter{enumi}{5}
    \item \textbf{Falsch.} Citation : \all{„Der Bundespräsident schlägt einen Kandidaten für das Kanzleramt vor, der dann vom Bundestag ohne Aussprache geheim gewählt wird.“} Le chancelier est élu par le Bundestag, pas directement par le peuple.
    \item \textbf{Falsch.} Citation : \all{„Briefwahl ist für alle ohne Angabe von Gründen möglich.“}
    \item \textbf{Falsch.} Citation : \all{„...im Gegensatz zur Diktatur, wo Andersdenkende unterdrückt werden.“}
\end{enumerate}
\textbf{Barème indicatif :} questions 1--5 : 2 pts chacune (1 pt contenu, 1 pt langue) ; Vrai/Faux : 2 pts chacun (1 pt réponse, 1 pt citation exacte). \textbf{Points de vigilance :} répondre par des phrases complètes ; recopier la citation sans la déformer ; ne pas traduire, répondre en allemand.
\end{corrige}

\begin{corrige}[Exercice 10 — Thème et Version]
\textbf{Partie A (Thème) — proposition de traduction :}

\all{Die Demokratie ist ein politisches System, in dem die Macht vom Volk kommt. Weil die Bürger ihre Vertreter wählen, kontrolliert das Parlament die Regierung. Obwohl die Entscheidungen oft langsam sind, wird die Freiheit durch die Verfassung geschützt. Deshalb ziehe ich die Demokratie der Diktatur vor. Morgen werden wir für die Zukunft unseres Landes wählen.}

Points de vigilance : \all{vorziehen} + datif (\all{der Diktatur}) ; passif \all{wird geschützt} ; connecteurs \all{weil}, \all{obwohl} (verbe en fin), \all{deshalb} (inversion).

\textbf{Partie B (Version) — proposition de traduction :}

« L'élection du chancelier fédéral par le Bundestag est un acte central de la démocratie parlementaire. Le candidat est élu au scrutin secret, sans débat. S'il obtient la majorité nécessaire, il est nommé par le président fédéral. Si aucune majorité n'est atteinte, le président peut dissoudre le Bundestag. La constitution garantit à cette procédure la primauté sur les coalitions instables. »

Points de vigilance : rendre le Nominalstil (\all{die Wahl des Bundeskanzlers}) par « l'élection du chancelier » ; les phrases sans sujet introduites par le verbe (\all{Erreicht er...}) sont des conditionnelles : « S'il obtient... ».

\textbf{Barème indicatif :} 10 pts par partie ; -0,5 pt par faute de grammaire, -0,25 pt par faute de vocabulaire.
\end{corrige}

\begin{corrige}[Exercice 11 — Production écrite type Bac]
\textbf{Proposition de rédaction modèle :}

\all{Meiner Meinung nach ist die Demokratie das beste politische System, weil die Macht vom Volk ausgeht. In einer Demokratie werden die Abgeordneten in freien und geheimen Wahlen gewählt, und die Rechte der Bürger werden durch die Verfassung geschützt. In einer Diktatur dagegen gibt es keine freie Opposition, deshalb werden Andersdenkende oft unterdrückt. Obwohl Entscheidungen in einer Demokratie manchmal langsam sind, ist die Gewaltenteilung ein großer Vorteil, denn das Parlament kontrolliert die Regierung. Ich bin der Ansicht, dass die Freiheit wichtiger ist als schnelle Entscheidungen. In Zukunft werde ich immer wählen gehen, weil jede Stimme zählt.}

\textbf{Vérification des critères :}
\begin{itemize}
    \item Connecteurs : \all{weil}, \all{deshalb}, \all{obwohl}, \all{denn} (4 connecteurs).
    \item Passif : \all{werden ... gewählt}, \all{werden ... geschützt}, \all{werden ... unterdrückt}.
    \item Futur I : \all{werde ich immer wählen gehen}.
    \item Vocabulaire thématique : \all{Demokratie, Diktatur, Wahlen, Verfassung, Opposition, Regierung, Parlament}.
    \item Avis personnel : \all{Meiner Meinung nach}, \all{Ich bin der Ansicht, dass...}.
\end{itemize}
\textbf{Barème indicatif (sur 20) :} respect des contraintes (connecteurs, passif, futur, vocabulaire) : 8 pts ; correction grammaticale (ordre des mots, cas, conjugaison) : 6 pts ; richesse du lexique et cohérence de l'argumentation : 4 pts ; présentation et longueur : 2 pts.
\textbf{Points de vigilance :} majuscules à tous les noms ; verbe en fin après \all{weil}, \all{obwohl}, \all{dass} ; inversion après \all{deshalb} ; ne pas écrire \all{*die Demokratie ist gewählt} sans auxiliaire complet.
\end{corrige}

\newpage

\coursec{Fiche bilan}

\begin{popfigure}
\begin{tikzpicture}[
  mindmap,
  grow cyclic,
  every node/.style={concept, execute at begin node=\hskip0pt},
  concept color=popblue,
  level 1/.append style={level distance=4.5cm, sibling angle=360/6},
  level 2/.append style={level distance=3cm, sibling angle=360/4},
  text=white,
  font=\small\bfseries
]
\node[concept, scale=1.2, align=center]{AL24\\Systèmes\\politiques}
  child[concept color=poporange] { node[concept, align=center]{Texte support\\„Politische Systeme“} }
  child[concept color=popgreen] { node[concept, align=center]{Lexique\\thématique} }
  child[concept color=poppurple] { node[concept, align=center]{Grammaire 1\\Nominalstil} }
  child[concept color=poppink] { node[concept, align=center]{Grammaire 2\\Konnektoren} }
  child[concept color=popgold] { node[concept, align=center]{Conjugaison\\Passiv / Futur} }
  child[concept color=popblue] { node[concept, align=center]{Expression\\écrite/orale} };
\end{tikzpicture}
\legende{Carte mentale de la leçon AL24 — Systèmes politiques (Terminale A)}
\end{popfigure}

\begin{bilan}

\end{bilan}

\courssub{Lexique clé : die politischen Systeme}

\begin{center}
\begin{tabularx}{\linewidth}{|X|X|l|}
\hline
\rowcolor{popblueL}
\textbf{Deutsch (Artikel, Plural)} & \textbf{Français} & \textbf{Famille / Dérivés} \\
\hline
die Demokratie, -n & la démocratie & demokratisch, der Demokrat, demokratisieren \\
die Diktatur, -en & la dictature & diktatorisch, der Diktator, diktieren \\
die Wahl, -en & l'élection, le vote & wählen, der Wähler, die Wahlurne \\
die Partei, -en & le parti (politique) & parteiisch, der Parteitag, die Parteizentrale \\
der Präsident, -en & le président & das Präsidium, präsidial, präsidieren \\
die Verfassung, -en & la constitution & verfassungsgemäß, der Verfassungsgerichtshof \\
der Bundeskanzler, - & le chancelier fédéral & die Bundeskanzlerin, das Kanzleramt \\
der Bundestag, -e & le Bundestag (parlement) & der Abgeordnete, die Bundestagswahl \\
das Grundgesetz, -e & la Loi fondamentale & grundgesetzlich \\
die Gewaltenteilung & la séparation des pouvoirs & gewaltenteilig, die Legislative, die Exekutive, die Judikative \\
\hline
\end{tabularx}
\end{center}

\vspace{0.5em}
\courssub{Règles de grammaire essentielles}

\begin{center}
\begin{tabularx}{\linewidth}{|X|X|}
\hline
\rowcolor{popblueL}
\textbf{Notion} & \textbf{Règle / Structure / Exemple} \\
\hline
\cle{Nominalstil} (Style nominal) & Transformer verbes/adjectifs en noms (substantivés).\\
& \all{Die Wahl findet statt.} $\rightarrow$ \all{Die Durchführung der Wahl.} \\
& \textbf{Indices :} Artikel (der/die/das), Prépositions (bei, durch, für, ohne, nach, vor, zu, in, an, auf...), Majuscule. \\
& \textit{But :} Langue administrative, scientifique, journalistique (concision, objectivité). \\
\hline
\cle{Konnektoren} (Connecteurs logiques) & \textbf{Coordination :} und, aber, oder, denn, sondern. (Position 0).\\
& \textbf{Subordination :} weil, dass, obwohl, wenn, als, damit, um...zu. (Verbe à la fin).\\
& \textbf{Adverbes / Prépositions :} deshalb, trotzdem, außerdem, daher, infolge, wegen + Gen., trotz + Gen. (Verbe en 2e position).\\
\hline
\cle{Passiv} (Passif) & \textbf{Vorgangspassiv} (action) : werden + Partizip II. \all{Die Wahl wird organisiert.}\\
& \textbf{Zustandspassif} (état) : sein + Partizip II. \all{Die Wahl ist organisiert.}\\
& \textbf{Temps au programme :} Präsens, Präteritum, Perfekt, Futur I, Konjunktiv II.\\
& \textit{Agent :} von + Dat (personne), durch + Akk (moyen/instrument). \\
\hline
\cle{Futur I} & werden + Infinitif. \all{Die Partei wird gewinnen.} (Prédiction, promesse, probabilité).\\
\cle{Futur II} & werden + Partizip II + haben/sein. \all{Die Wahl wird stattgefunden haben.} (Antériorité future, hypothèse sur le passé). \\
\hline
\end{tabularx}
\end{center}

\vspace{0.5em}
\courssub{Méthodes express}

\begin{itemize}
\item \textbf{Compréhension :} 1. Lire le titre + images $\rightarrow$ hypothèses. 2. Lecture globale (sujet, ton, structure). 3. Lecture détaillée (vocabulaire inconnu $\rightarrow$ déduction contextuelle). 4. Répondre aux questions \all{Wo? Wer? Was? Wann? Warum? Wie?} en phrases complètes.
\item \textbf{Résumé :} Identifier les \cle{idées forces} (1 phrase/paragraphe). Relier par connecteurs (\all{zunächst, dann, schließlich, zusammenfassend}). Style nominal privilégié. Longueur : $\approx$ 1/3 du texte.
\item \textbf{Traduction Thème (Fr $\rightarrow$ All) :} Identifier le verbe $\rightarrow$ conjuguer $\rightarrow$ placer (2e pos / fin). Vérifier cas (Nom/Akk/Dat/Gen) et genre/pluriel des noms. Attention aux \cle{faux amis} (\all{aktuell $\neq$ actuel, werden $\neq$ vouloir}).
\item \textbf{Production guidée :} Plan en 3 étapes (These - Argumente - Fazit). Utiliser \cle{chunks} appris : \all{Meiner Meinung nach...}, \all{Ein Argument dafür ist...}, \all{Man muss bedenken, dass...}.
\end{itemize}



\begin{aretenir}[Checklist avant le Bac]
$\square$ Je sais décliner \cle{die Demokratie, die Wahl, der Präsident, die Verfassung, das Grundgesetz} (singulier/pluriel, 4 cas).\\
$\square$ Je reconnais un \cle{Nominalstil} (nom majuscule + article/préposition) et je sais le transformer en \cle{Verbalstil} (phrase complète avec verbe conjugué).\\
$\square$ Je distingue \cle{Vorgangspassiv} (werden + Partizip II = action) et \cle{Zustandspassiv} (sein + Partizip II = état résultant).\\
$\square$ Je conjugue le \cle{Passiv} au Présent, Prétérit, Perfect, Futur I et Konjunktiv II pour les verbes forts/irréguliers (werden, sein, haben, geben, finden, wählen).\\
$\square$ J'utilise correctement les \cle{Konnektoren} : coordination (verbe 2e), subordination (verbe fin), adverbes (verbe 2e), prépositions génitifs (wegen, trotz, während, infolge).\\
$\square$ Je forme le \cle{Futur I} (werden + Inf.) pour la prédiction et le \cle{Futur II} (werden + Part.II + haben/sein) pour l'antériorité/hypothèse.\\
$\square$ Je maîtrise le vocabulaire thématique (Parti, élection, séparation des pouvoirs, chancelier, Loi fondamentale) avec genre et pluriel.\\
$\square$ Je sais structurer un \cle{résumé} (idées forces + connecteurs + style nominal) et une \cle{prise de position} (These - Argumente - Beispiel - Fazit).\\
$\square$ Je repère les \cle{pièges francophones} : ordre des mots (verbe 2e / fin), majuscules des noms, cas après prépositions, faux amis (konsequent, aktuell, eventuell).\\
$\square$ Je suis capable de présenter oralement (2-3 min) un système politique (Allemagne, Cameroun, Suisse) avec vocabulaire précis et connecteurs logiques.
\end{aretenir}

\findecours

\end{document}
